% Chapter VII, first half: §64–71
% The Stability of Couette Flow

\chapter{THE STABILITY OF COUETTE FLOW}
\label{ch:7}

\section{Introduction}\label{sec:7-64}
In the last five chapters we have considered the onset of instability
under a variety of externally impressed conditions; but the origin of the
instability was always the same: a potentially unstable arrangement of
the fluid resulting from a prevailing adverse temperature gradient.  In this
and in the following two chapters we shall consider cases of instabilities
which have a different origin: a potentially unstable arrangement of flow
resulting from a prevailing adverse gradient of angular momentum.  The
simplest example of such instability occurs in \emph{Couette flow}, i.e.\ in the
steady circular flow of a liquid between two rotating coaxial cylinders.

In this chapter we shall consider the simple Couette flow of inviscid
and viscid fluids.  In Chapter~VIII we shall consider some examples of
more general curved flows; and in Chapter~IX we shall consider the
effect of an axial magnetic field on Couette flow.

\section{The physical problem}\label{sec:7-65}
We consider, then, the permissible stationary circular flows of an
incompressible fluid between two rotating coaxial cylinders.  As we shall
presently see, in the absence of viscosity, the class of such permissible
flows is very wide: indeed, if $\Omega$ denotes the angular velocity of rotation
about the axis, then the equations of motion allow $\Omega$ to be an arbitrary
function of the distance $r$ from the axis, provided the velocities in the
radial and the axial directions are zero.  But if viscosity is present, the
class becomes very restricted; in fact, in the absence of any transverse
pressure gradient, the most general form of $\Omega$ which is allowed is
\begin{equation}
\label{eq:7-1}
\Omega(r) = A + B/r^2,
\end{equation}
where $A$ and $B$ are two constants which are related to the angular
velocities $\Omega_1$ and $\Omega_2$ with which the inner and the outer cylinders are
rotated.  Thus, if $R_1$ and $R_2$ ($> R_1$) are the radii of the two cylinders,
then
\begin{equation}
\label{eq:7-2}
\Omega_1 = A + B/R_1^2 \quad \text{and} \quad \Omega_2 = A + B/R_2^2.
\end{equation}
Solving for $A$ and $B$ in terms of $\Omega_1$ and $\Omega_2$, we have
\begin{equation}
\label{eq:7-3}
A = -\Omega_1 \eta^2 \frac{1-\mu/\eta^2}{1-\eta^2} \quad \text{and} \quad
B = \Omega_1 \frac{R_1^2(1-\mu)}{1-\eta^2},
\end{equation}
where
\begin{equation}
\label{eq:7-4}
\mu = \Omega_2/\Omega_1 \quad \text{and} \quad \eta = R_1/R_2.
\end{equation}

The questions to which we seek answers are the following.  In the
absence of viscosity when $\Omega(r)$ can be an arbitrary function of $r$, what is
the necessary and sufficient condition for the stability of the flow?
What does this condition imply for the distribution of $\Omega$ given by~\eqref{eq:7-1}?
And, finally, how is the condition for inviscid fluids, as applied to~\eqref{eq:7-1},
modified when account is taken of viscosity?  These are the questions
to which the present chapter is addressed.

\section{Rayleigh's criterion}\label{sec:7-66}
Rayleigh stated: \emph{in the absence of viscosity, the necessary and sufficient
condition for a distribution of angular velocity $\Omega(r)$ to be stable is}
\begin{equation}
\label{eq:7-5}
\frac{d}{dr}(r^2\Omega)^2 > 0
\end{equation}
\emph{everywhere in the interval; and, further, that the distribution is unstable if
$(r^2\Omega)^2$ should decrease anywhere inside the interval.}

Since $|r^2\Omega|$ is the angular momentum, per unit mass, of a fluid element
about the axis of rotation, an alternative way of stating Rayleigh's
criterion is: \emph{a stratification of angular momentum about an axis is stable if
and only if it increases monotonically outward.}

Rayleigh did not establish his criterion by an analytical discussion of
the relevant perturbation equations.  Instead, he argued as follows.

The hydrodynamical equations governing an incompressible inviscid
fluid in cylindrical polar coordinates $(r,\theta,z)$ are:\footnote{The considerations which follow are equally applicable if an external force derivable
from an axisymmetric potential is present.}
\begin{equation}
\label{eq:7-6}
\frac{\partial u_r}{\partial t} + (\mathbf{u}.\mathrm{grad})\,u_r - \frac{u_\theta^2}{r}
= -\frac{\partial}{\partial r}\!\left(\frac{p}{\rho}\right),
\end{equation}
\begin{equation}
\label{eq:7-7}
\frac{\partial u_\theta}{\partial t} + (\mathbf{u}.\mathrm{grad})\,u_\theta + \frac{u_r u_\theta}{r}
= -\frac{1}{r}\frac{\partial}{\partial \theta}\!\left(\frac{p}{\rho}\right),
\end{equation}
\begin{equation}
\label{eq:7-8}
\frac{\partial u_z}{\partial t} + (\mathbf{u}.\mathrm{grad})\,u_z
= -\frac{\partial}{\partial z}\!\left(\frac{p}{\rho}\right),
\end{equation}
where
\begin{equation}
\label{eq:7-9}
(\mathbf{u}.\mathrm{grad}) = u_r \frac{\partial}{\partial r}
+ \frac{u_\theta}{r}\frac{\partial}{\partial \theta}
+ u_z \frac{\partial}{\partial z}.
\end{equation}
We also have the equation of continuity,
\begin{equation}
\label{eq:7-10}
\frac{\partial u_r}{\partial r}
+ \frac{u_r}{r}
+ \frac{1}{r}\frac{\partial u_\theta}{\partial \theta}
+ \frac{\partial u_z}{\partial z} = 0.
\end{equation}
Those equations clearly allow the stationary solution,
\begin{equation}
\label{eq:7-11}
u_r = u_z = 0 \quad \text{and} \quad u_\theta = V(r),
\end{equation}
where $V(r)$ is an arbitrary function of $r$.  In terms of $V(r)$, the pressure
distribution is determined by
\begin{equation}
\label{eq:7-12}
p = \rho \int \frac{dr}{r}\, V^2(r).
\end{equation}

For axisymmetric motions to which the rest of the discussion in this
section will be restricted, equations~\eqref{eq:7-6}--\eqref{eq:7-8} reduce to:
\begin{equation}
\label{eq:7-13}
\frac{\partial u_r}{\partial t}
+ u_r \frac{\partial u_r}{\partial r}
+ u_z \frac{\partial u_r}{\partial z}
- \frac{u_\theta^2}{r}
= -\frac{\partial}{\partial r}\!\left(\frac{p}{\rho}\right),
\end{equation}
\begin{equation}
\label{eq:7-14}
\frac{\partial u_\theta}{\partial t}
+ u_r \frac{\partial u_\theta}{\partial r}
+ u_z \frac{\partial u_\theta}{\partial z}
+ \frac{u_r u_\theta}{r} = 0,
\end{equation}
\begin{equation}
\label{eq:7-15}
\frac{\partial u_z}{\partial t}
+ u_r \frac{\partial u_z}{\partial r}
+ u_z \frac{\partial u_z}{\partial z}
= -\frac{\partial}{\partial z}\!\left(\frac{p}{\rho}\right).
\end{equation}
Equation~\eqref{eq:7-14} is equivalent to
\begin{equation}
\label{eq:7-16}
\frac{d}{dt}(r u_\theta) = \frac{d}{dt}(r^2\Omega) = 0.
\end{equation}
Therefore, the angular momentum $L$ ($= r^2\Omega$) of a fluid element, per unit
mass, remains constant as we follow it with its motion.  And the motions
in the radial and the axial directions take place as though $u_\theta$ were absent
and, instead, a force $u_\theta^2 r = L^2/r^3$ were acting in the radial direction.  Since
$L$ is a constant of the motion, we can associate with the force $L^2/r^3$ a
potential energy $\frac{1}{2}L^2/r^2$.  This potential energy is the same as the
kinetic energy associated with the velocity $u_\theta$ in the transverse direction.

Suppose now that we interchange the fluid contained in two elementary
rings, of equal heights and masses, at $r = r_1$ and $r = r_2$ ($> r_1$).  If
$dr_1$ and $dr_2$ are the radial extents of the rings, the equality of their
masses requires $2\pi r_1\, dr_1 = 2\pi r_2\, dr_2 = dS$ (say).  In view of the constancy
of $L$ with the motion, the fluid at $r_1$, after the interchange, will have
the same angular momentum (namely, $L_1$) which it had at $r_1$ before the
interchange; similarly, the fluid at $r_2$, after the interchange, will have the
same angular momentum (namely, $L_2$) which it had at $r_2$.  As a result,
the change in the kinetic energy (or, what is the same thing, the change
in the centrifugal potential energy) is proportional to
\begin{equation}
\label{eq:7-17}
\left(\frac{L_1^2}{r_2^2} + \frac{L_2^2}{r_1^2}\right)
- \left(\frac{L_1^2}{r_1^2} + \frac{L_2^2}{r_2^2}\right)
= (L_2^2 - L_1^2)\!\left(\frac{1}{r_1^2} - \frac{1}{r_2^2}\right)\! dS.
\end{equation}
Remembering that $r_2 > r_1$, we observe that this is positive or negative
according as $L_2^2$ is greater than $L_1^2$ or less than $L_1^2$.  Consequently, if $L^2$
is a monotonic increasing function of $r$, no interchange of fluid rings
such as we have imagined can occur without a source of energy; and this
means stability.  On the other hand, if $L^2$ should decrease anywhere,
then an interchange of fluid rings in this region will result in a liberation
of energy; and this means instability.

While the foregoing arguments of Rayleigh make his criterion a very
likely one by drawing attention to its physical origin, one should still
like to establish it directly from the relevant perturbation equations;
this we shall do in \S~67.  Meantime, it is useful to see what Rayleigh's
criterion implies for the particular distribution of $\Omega$ which is permissible
for viscid fluids.  As the \emph{discriminant} for stability, we shall use
\begin{equation}
\label{eq:7-18}
\Phi(r) = \frac{1}{r^3}\frac{d}{dr}(r^2\Omega)^2
= \frac{2\Omega}{r}\frac{d}{dr}(r^2\Omega).
\end{equation}
Rayleigh's criterion requires
\begin{equation}
\label{eq:7-19}
\Phi(r) > 0 \quad \text{for stability.}
\end{equation}
When $\Omega$ has the form given by~\eqref{eq:7-1},
\begin{equation}
\label{eq:7-20}
\Phi = 4A(A + B/r^4).
\end{equation}
It is convenient to measure $r$ in units of the radius $R_2$ of the outer
cylinder.  Then
\begin{equation}
\label{eq:7-21}
\frac{4AB}{R_2^4}\!\left(\frac{1}{r^4}+\frac{A\,R_2^4}{B}\right).
\end{equation}
With the values of $A$ and $B$ given in~\eqref{eq:7-3}, we have
\begin{equation}
\label{eq:7-22}
\Phi = -4\Omega_1^2\,\eta^4\,\frac{(1-\mu)(1-\mu/\eta^2)}{(1-\eta^2)^2}
\!\left(\frac{1}{r^4} - \kappa\right),
\end{equation}
where
\begin{equation}
\label{eq:7-23}
\kappa = \frac{A\, R_2^4}{B} = \frac{1-\mu/\eta^2}{1-\mu}.
\end{equation}
Clearly,
\begin{equation}
\label{eq:7-24}
\Phi > 0 \quad \text{for} \quad \eta \leqslant r \leqslant 1,
\end{equation}
so long as
\begin{equation}
\label{eq:7-25}
\mu = \Omega_2/\Omega_1 > \eta^2.
\end{equation}
Therefore, Rayleigh's criterion applied to the distribution~\eqref{eq:7-1} requires
that \emph{for stability, the outer cylinder must rotate with an angular speed
greater than $\eta^2$ times that of the inner cylinder and in the same sense.}  In
the $(\Omega_1, \Omega_2)$-plane\footnote{Without loss of generality, we may suppose that the sense of rotation of the inner
cylinder is positive.} (see Fig.~62), the regions of stability are delimited
by the positive $\Omega_2$-axis and the Rayleigh line $\Omega_2 = \Omega_1 \eta^2$.  If the effects
of viscosity are ignored, we must have instability to the left of the
Rayleigh line; and it is of interest to determine how Rayleigh's criterion
is violated for flow represented in this part of the plane.  With $\kappa$ given
by~\eqref{eq:7-23}, we find
\begin{equation}
\label{eq:7-26}
\frac{1}{r^4} - \kappa = \begin{cases}
\mu\,\dfrac{1-\eta^2}{\eta^4(1-\mu)} & \text{for } r = 1,\\[6pt]
\dfrac{(1-\eta^2)}{\eta^4(1-\mu)} & \text{for } r = \eta.
\end{cases}
\end{equation}
Accordingly,
\begin{equation}
\label{eq:7-27}
\Phi(r) < 0 \quad \text{for} \quad \eta \leqslant r \leqslant 1,
\end{equation}
so long as
\begin{equation}
\label{eq:7-28}
0 < \mu < \eta^2.
\end{equation}
Therefore, for $0 < \mu < \eta^2$, Rayleigh's criterion is violated everywhere

\figurePlaceholder{62}{The $(\Omega_1,\Omega_2)$-plane: in the limit of zero viscosity,
the stable and the unstable regions in this plane are separated
by the Rayleigh line, $\Omega_2 = \Omega_1(R_1/R_2)^2$; $R_1$ and $R_2$ are the
radii and $\Omega_1$ and $\Omega_2$ are the angular velocities of the inner
and the outer cylinders, respectively.}

in the fluid; and we may say that instability occurs throughout the
fluid.  On the other hand, if the cylinders are rotating in opposite
directions and $\mu < 0$, $\Phi$ is positive in a part of the fluid adjoining the
outer cylinder; and it is negative in the remaining part of the fluid
adjoining the inner cylinder.  The two parts are separated by the \emph{nodal
surface} on which $\Omega = 0$.  The radius $r_0$ of this surface in the unit $R_2$ is
given by
\begin{equation}
\label{eq:7-29}
r_0 = \frac{1}{\sqrt{\kappa}} = \eta\sqrt{\frac{1+|\mu|}{{\eta^2}+|\mu|}}\,;
\end{equation}
and
\begin{equation}
\label{eq:7-30}
\left.\begin{aligned}
\Phi(r) &> 0 \quad \text{for} \quad r_0 < r < 1\\
\Phi(r) &< 0 \quad \text{for} \quad \eta < r < r_0
\end{aligned}\right\}.
\end{equation}
Thus, when $\mu < 0$, Rayleigh's criterion is violated only interior to the
nodal surface.  The instability which is predicted for all negative $\mu$ in
the absence of viscosity, is therefore, to be associated with this instability
of the flow interior to $r_0$.  On this account we may say that for $\mu < 0$,
the instability is only partial.  Also, as $\mu \to -\infty$, $r_0 \to \eta$ and the part
of the fluid which is unstable according to Rayleigh's criterion becomes
vanishingly small.

If we now consider the effect of viscosity on the stability of the flow
prescribed by~\eqref{eq:7-1}, we must on general grounds expect that it will postpone the onset of instability beyond the point predicted by Rayleigh's
criterion; and that for a given $\kappa$ (i.e.\ $\mu$ and $\eta$) the factor,
\begin{equation}
\label{eq:7-31}
4\Omega_1^2\,\eta^4\,\frac{(1-\mu)(1-\mu/\eta^2)}{(1-\eta^2)^2}
\quad (\mu < \eta^2),
\end{equation}
which occurs in Rayleigh's discriminant must exceed a certain critical
value, depending on $\nu$ and $R_2$, before instability can set in.  The non-dimensional number in terms of which the criterion for stability will be
expressed is, then,
\begin{equation}
\label{eq:7-32}
T = \frac{4\Omega_1^2}{\nu^2}\,R_2^4\,\frac{(1-\mu)(1-\mu/\eta^2)}{(1-\eta^2)^2}.
\end{equation}
This is the proper definition of the \emph{Taylor number} for this problem.  For
a given $\eta$ and $\mu$, instability will set in for a certain critical Taylor number
$T_c$; and the central problem of this subject is the determination of $T_c$ as a
function of $\eta$ and $\mu$.

\section{Analytical discussion of the stability of inviscid Couette flow}
\label{sec:7-67}
We now return to an analytical discussion of the stability of the
stationary flow,
\begin{equation}
\label{eq:7-33}
u_r = u_z = 0 \quad \text{and} \quad u_\theta = V(r) = r\Omega(r),
\end{equation}
which the equations of inviscid motion~\eqref{eq:7-6}--\eqref{eq:7-9} allow.
In~\eqref{eq:7-33}, $V(r)$ is an arbitrary function of $r$.

Consider an infinitesimal perturbation of the flow represented by the
solution~\eqref{eq:7-33}.  Let the perturbed state be described by
\begin{equation}
\label{eq:7-34}
u_r, \quad V\!+\!u_\theta, \quad u_z, \quad \text{and} \quad
\varpi\;(= \delta p/\rho).
\end{equation}
The linear equations governing these perturbations are
\begin{equation}
\label{eq:7-35}
\frac{\partial u_r}{\partial t}
+ \frac{V}{r}\frac{\partial u_r}{\partial \theta}
- 2\frac{V}{r}\,u_\theta
= -\frac{\partial \varpi}{\partial r},
\end{equation}
\begin{equation}
\label{eq:7-36}
\frac{\partial u_\theta}{\partial t}
+ \frac{V}{r}\frac{\partial u_\theta}{\partial \theta}
+ \left(\frac{dV}{dr}+\frac{V}{r}\right) u_r
= -\frac{1}{r}\frac{\partial \varpi}{\partial \theta},
\end{equation}
\begin{equation}
\label{eq:7-37}
\frac{\partial u_z}{\partial t}
+ \frac{V}{r}\frac{\partial u_z}{\partial \theta}
= -\frac{\partial \varpi}{\partial z},
\end{equation}
and the equation of continuity,
\begin{equation}
\label{eq:7-38}
\frac{\partial u_r}{\partial r}
+ \frac{u_r}{r}
+ \frac{1}{r}\frac{\partial u_\theta}{\partial \theta}
+ \frac{\partial u_z}{\partial z} = 0.
\end{equation}

In accordance with the general procedure of treating these problems,
we analyse the disturbance into normal modes.  In the present instance,
it is natural to suppose that the various quantities describing the
perturbation have a $(t,\theta,z)$-dependence given by
\begin{equation}
\label{eq:7-39}
e^{i(\sigma t + m\theta + kz)},
\end{equation}
where $p$ is a constant (which can be complex), $m$ is an integer (which
can be positive, zero, or negative), and $k$ is the wave number of the
disturbance in the $z$-direction.

Let $u_r(r)$, $u_\theta(r)$, $u_z(r)$, and $\varpi(r)$ now denote the amplitudes of the
respective perturbations whose $(t,\theta,z)$-dependence is given by~\eqref{eq:7-39}.
Equations~\eqref{eq:7-35}--\eqref{eq:7-38} then give
\begin{equation}
\label{eq:7-40}
i\sigma u_r - 2\Omega u_\theta = -\frac{d\varpi}{dr},
\end{equation}
\begin{equation}
\label{eq:7-41}
i\sigma u_\theta + \left(\Omega + \frac{d}{dr} r\Omega\right) u_r
= -\frac{im}{r}\,\varpi,
\end{equation}
\begin{equation}
\label{eq:7-42}
i\sigma u_z = -ik\varpi,
\end{equation}
and
\begin{equation}
\label{eq:7-43}
\frac{du_r}{dr} + \frac{u_r}{r} + \frac{im}{r}\, u_\theta + ik u_z = 0,
\end{equation}
where
\begin{equation}
\label{eq:7-44}
\sigma = p + m\Omega.
\end{equation}
(Note that through $\Omega$, $\sigma$ is a function of $r$.)

\subsection*{(a) \emph{The equations in terms of the Lagrangian displacement}}
Consider the variables $\xi_r$, $\xi_\theta$, and $\xi_z$ related to $u_r$, $u_\theta$, and $u_z$ by
\begin{equation}
\label{eq:7-45}
u_r = i\sigma \xi_r, \quad
u_\theta = i\sigma \xi_\theta - r\frac{d\Omega}{dr}\,\xi_r, \quad \text{and} \quad
u_z = i\sigma \xi_z.
\end{equation}
Defined in this manner, $\boldsymbol{\xi}$ is the \emph{Lagrangian displacement} appropriate to
the problem on hand.

It can be readily verified that consistent with the meaning of $\boldsymbol{\xi}$ as the
Lagrangian displacement, the solenoidal character of $\mathbf{u}$ implies the
solenoidal character of $\boldsymbol{\xi}$; thus
\begin{equation}
\label{eq:7-46}
\frac{d\xi_r}{dr} + \frac{\xi_r}{r} + \frac{im}{r}\,\xi_\theta + ik\xi_z = 0.
\end{equation}
In terms of the variables $\xi_r$, $\xi_\theta$, and $\xi_z$, equations~\eqref{eq:7-40}--\eqref{eq:7-42} become
\begin{equation}
\label{eq:7-47}
\left(\sigma^2 - 2r\Omega\frac{d\Omega}{dr}\right)\xi_r + 2i\Omega\sigma\xi_\theta
= \frac{d\varpi}{dr},
\end{equation}
\begin{equation}
\label{eq:7-48}
\sigma^2 \xi_\theta - 2i\Omega\sigma\xi_r = \frac{im}{r}\,\varpi,
\end{equation}
and
\begin{equation}
\label{eq:7-49}
\sigma^2 \xi_z = ik\varpi.
\end{equation}
Multiplying equations~\eqref{eq:7-48} and~\eqref{eq:7-49} by $im/r$ and $ik$, respectively,
adding, and making use of equation~\eqref{eq:7-46}, we obtain
\begin{equation}
\label{eq:7-50}
\sigma^2\!\left(\frac{d\xi_r}{dr}+\frac{\xi_r}{r}\right)
+ 2\frac{m\Omega}{r}\sigma\xi_r
= \left(\frac{m^2}{r^2}+k^2\right)\varpi.
\end{equation}
Next eliminating $\xi_\theta$ between equations~\eqref{eq:7-47} and~\eqref{eq:7-48}, we have
\begin{equation}
\label{eq:7-51}
\left(\sigma^2 - 2r\Omega\frac{d\Omega}{dr}\right)\xi_r
+ \frac{2\Omega}{\sigma}\!\left(2i\Omega\sigma\xi_r
+ \frac{im}{r}\,\varpi\right)
= \frac{d\varpi}{dr}.
\end{equation}
Rearranging the terms in this equation, we obtain
\begin{equation}
\label{eq:7-52}
[\sigma^2 - \Phi(r)]\xi_r = \frac{d\varpi}{dr}
+ \frac{2m\Omega}{r\sigma}\,\varpi,
\end{equation}
where
\begin{equation}
\label{eq:7-53}
\Phi(r) = 2r\Omega\frac{d\Omega}{dr} + 4\Omega^2
= \frac{2\Omega}{r}\frac{d}{dr}(r^2\Omega)
\end{equation}
is the Rayleigh discriminant we have already introduced in \S~66 (equation~\eqref{eq:7-18}).

Equation~\eqref{eq:7-52} must be considered in conjunction with equation~\eqref{eq:7-50}, which we shall rewrite in the form
\begin{equation}
\label{eq:7-54}
\frac{1}{r}\frac{d}{dr}(r\xi_r)
- \frac{2m\Omega}{\sigma r}\,\xi_r
= \frac{1}{\sigma^2}\!\left(\frac{m^2}{r^2}+k^2\right)\varpi.
\end{equation}

If the fluid is confined between two coaxial cylinders of radii $R_1$ and
$R_2$, we must require that the radial components of the velocity vanish
for these values of $r$.  Thus, equations~\eqref{eq:7-52} and~\eqref{eq:7-54} must be considered
together with the boundary conditions
\begin{equation}
\label{eq:7-55}
\xi_r = 0 \quad \text{for } r = R_1 \text{ and } R_2.
\end{equation}

We shall presently see that equations~\eqref{eq:7-52} and~\eqref{eq:7-54}, for real $\sigma$, constitute a self-adjoint system with respect to the boundary conditions~\eqref{eq:7-55}; and that in consequence their solution can be reduced to a variational problem; this was, indeed, the object of introducing the Lagrangian
displacement as the variable.

\subsection*{(b) \emph{The case $m = 0$}}
When $m = 0$, $\sigma = p$ and equations~\eqref{eq:7-52} and~\eqref{eq:7-54} become
\begin{equation}
\label{eq:7-56}
[p^2 - \Phi(r)]\xi_r = \frac{d\varpi}{dr},
\end{equation}
and
\begin{equation}
\label{eq:7-57}
\frac{1}{r}\frac{d}{dr}(r\xi_r) = \frac{k^2}{p^2}\,\varpi.
\end{equation}
Eliminating $\varpi$ between these equations, we obtain
\begin{equation}
\label{eq:7-58}
\frac{d}{dr}\!\left\{
\frac{1}{r}\frac{d}{dr}(r\xi_r)\right\}
- k^2\xi_r = -\frac{k^2}{p^2}\,\Phi(r)\xi_r.
\end{equation}
The solution of this equation, together with the boundary conditions~\eqref{eq:7-55},
constitutes a characteristic value problem of the classical Sturm--Liouville type; and by appealing to the standard theorems of the subject,
we can reach the following conclusion.

\emph{The characteristic values of $k^2\!/p^2$ are all positive if $\Phi(r)$ is everywhere
positive; and they are all negative if $\Phi(r)$ is everywhere negative.  If $\Phi(r)$
should change sign anywhere in the interval $(R_1, R_2)$, then there are two
sets of real characteristic values which have the limit points $+\infty$ and $-\infty$.}

Since a negative $p^2$ means instability, it is clear that the foregoing
statements regarding the sign of the characteristic values of $k^2\!/p^2$ are
equivalent to a restatement of Rayleigh's criterion.

An alternative way of establishing the same results is instructive.
Multiply equation~\eqref{eq:7-58} by $r\xi_r$ and integrate over the range of $r$.  The
integral on the left-hand side can be transformed by an integration by
parts.  We obtain
\begin{equation}
\label{eq:7-59}
\int\!\left\{
\frac{1}{r}\!\left(\frac{d}{dr}r\xi_r\right)^{\!2}
+ k^2 r\xi_r^2
\right\} dr
= \frac{k^2}{p^2}\int \Phi(r)\, r\,\xi_r^2\, dr
\end{equation}
or
\begin{equation}
\label{eq:7-60}
\frac{p^2}{k^2} = \frac{\displaystyle\int \Phi(r)\, r\, \xi_r^2\, dr}
{\displaystyle\int\!\left\{
\frac{1}{r}[\!(d/dr)\, r\xi_r\!]^2 + k^2 r\xi_r^2
\right\} dr}.
\end{equation}
From equation~\eqref{eq:7-60} it is immediately apparent that $p^2\!/k^2$ is positive if
$\Phi(r)$ is everywhere positive, and is negative if $\Phi(r)$ is everywhere negative.
The further result that there exist unstable modes, if $\Phi(r)$ is anywhere
negative, can be deduced from the fact that we may regard equation~\eqref{eq:7-58} as the condition that
\begin{equation}
\label{eq:7-61}
I_1 = \int \Phi(r)\, r\, \xi_r^2\, dr
\end{equation}
is a maximum, or a minimum, for a given
\begin{equation}
\label{eq:7-62}
I_2 = \int\!\left\{
\frac{1}{r}\!\left(\frac{d}{dr}r\xi_r\right)^{\!2}
+ k^2 r\xi_r^2
\right\} dr.
\end{equation}
$p^2\!/k^2$ being, then, the value of $I_1/I_2$.  If $\Phi(r)$ should anywhere be negative,
$I_1$ admits a negative value and therefore a negative minimum, so that
one value at least of $p^2$ is negative, and one mode of disturbance is
unstable.

By comparing this alternative method of establishing Rayleigh's
criterion with the corresponding method of establishing the stability of
an incompressible fluid of variable density (Chapter~X, \S\,92), we observe
that $\Phi(r)$ in this problem, and the density gradient in the other problem,
play entirely equivalent roles.

\subsection*{(c) \emph{The case $m \neq 0$}}
The treatment of the general case is not as straightforward as the
case $m = 0$.  It appears that when $m \neq 0$, it is more convenient to consider
the problem of solving equations~\eqref{eq:7-52} and~\eqref{eq:7-54} subject to the boundary
conditions~\eqref{eq:7-55} as a characteristic value problem for $k^2$ for a given $p$,
rather than as a problem for $p$ for a given $k^2$.  This inversion of the usual
procedure is advantageous in so far as the problem can then be formulated
in terms of a variational principle.  Moreover, we shall find that for real
$p$ the problem is Hermitian and the characteristic values of $k^2$ are real.
The question now arises whether the solution of the stability problem
can be inferred from the solution of the characteristic value problem
inverted in the manner we have proposed.  That it should, in principle,
be possible to infer the solution becomes clear when we realize that there
is a unique dispersion relation which gives $p$ as an analytic function of $k^2$
(with, possibly, more than one branch); and that the relationship is independent of how it is determined.  And so long as the relation between
$p$ and $k^2$ is analytic, we should be able to answer the basic questions
concerning stability by determining the relation between real $p$ and $k^2$
(positive or negative) instead of, as is more usual, by determining the
relation between positive $k^2$ and $p$ (real or complex).

Returning to equations~\eqref{eq:7-52}--\eqref{eq:7-55}, we shall first show that for real $p$
the characteristic values of $k^2$ are real.  To show this, we shall multiply
equation~\eqref{eq:7-52} by $r\xi_r^*$ and integrate over the range of $r$.  We obtain
\begin{equation}
\label{eq:7-63}
\int r[\sigma^2-\Phi(r)]|\xi_r|^2\,dr
= \int \xi_r^*\!\left(r\frac{d\varpi}{dr}+\frac{2m\Omega}{\sigma}\,\varpi\right)dr.
\end{equation}
After an integration by parts, the right-hand side becomes
\begin{equation}
\label{eq:7-64}
-\int \varpi\!\left(\frac{d}{dr}(r\xi_r^*) - \frac{2m\Omega}{\sigma}\,\xi_r^*\right) dr;
\end{equation}
and making use of equation~\eqref{eq:7-54}, we can finally write
\begin{equation}
\label{eq:7-65}
\int r[\sigma^2-\Phi(r)]|\xi_r|^2\,dr
= -\int \frac{1}{\sigma^2}
\!\left(\frac{m^2}{r}+k^2 r\right)|\varpi|^2\,dr.
\end{equation}
By replacing every term in this equation by its complex conjugate, we
infer the identity of $k^2$ and $k^{*2}$; $k^2$ is, therefore, real.  When $k^2$ is real, the
proper functions, $\xi_r$ and $\varpi$, belonging to a particular $k^2$ (and a real $p$)
are also real, and we can rewrite equation~\eqref{eq:7-65} in the manner
\begin{equation}
\label{eq:7-66}
k^2 = \frac{\displaystyle\int
r[(\Phi-\sigma^2)\xi_r^2 - (m^2\varpi^2)/(r^2\sigma^2)]\,dr}
{\displaystyle\int r\,\varpi^2/\sigma^2\,dr}
\quad (\text{say}).
\end{equation}

We shall now show that the expression~\eqref{eq:7-66} for $k^2$ provides the basis
for a variational formulation of the problem.

First, it should be remarked that in evaluating $k^2$ in accordance with~\eqref{eq:7-66}, we must suppose that $\xi_r$ is determined in terms of $\varpi$ by the equation
\begin{equation}
\label{eq:7-67}
(\sigma^2-\Phi)\xi_r = \frac{d\varpi}{dr}+\frac{2m\Omega}{r\sigma}\,\varpi;
\end{equation}
and that an allowed $\varpi$ accordingly satisfies the boundary conditions
\begin{equation}
\label{eq:7-68}
\frac{d\varpi}{dr}+\frac{2m\Omega}{\sigma}\,\varpi = 0
\quad \text{for } r = R_1 \text{ and } R_2.
\end{equation}

Consider now the effect on $k^2$ of an arbitrary variation $\delta\varpi$ in $\varpi$ compatible only with the boundary conditions~\eqref{eq:7-68}.  If $\delta\xi_r$ denotes the corresponding variation in $\xi_r$, then, by~\eqref{eq:7-67}
\begin{equation}
\label{eq:7-69}
(\sigma^2-\Phi)\,\delta\xi_r = \frac{d}{dr}\,\delta\varpi
+ \frac{2m\Omega}{\sigma r}\,\delta\varpi;
\end{equation}
and
\begin{equation}
\label{eq:7-70}
\delta\xi_r = 0 \quad \text{for } r = R_1 \text{ and } R_2.
\end{equation}
From equation~\eqref{eq:7-66} it follows that
\begin{equation}
\label{eq:7-71}
\delta k^2 = \frac{1}{I_2}\bigl(\delta I_1 - \frac{I_1}{I_2}\,\delta I_2\bigr)
= \frac{1}{I_2}\bigl(\delta I_1 - k^2 \delta I_2\bigr),
\end{equation}
where
\begin{equation}
\label{eq:7-72}
\delta I_1 = 2\int r\bigl[(\Phi-\sigma^2)\,\xi_r\,\delta\xi_r
- \frac{m^2\varpi}{r^2\sigma^2}\,\delta\varpi\bigr]\,dr
\end{equation}
and
\begin{equation}
\label{eq:7-73}
\delta I_2 = 2\int r\varpi(\delta\varpi/\sigma^2)\,dr
\end{equation}
are the variations in $I_1$ and $I_2$ resulting from the variation in $\varpi$.

The terms in $\delta\xi_r$ in the expression~\eqref{eq:7-72} for $\delta I_1$ are the same as
\begin{equation}
\label{eq:7-74}
-2\int r\xi_r\!\left(\frac{d}{dr}\,\delta\varpi+\frac{2m\Omega}{\sigma r}\,\delta\varpi\right)dr.
\end{equation}
After an integration by parts,~\eqref{eq:7-74} can be rewritten in the form
\begin{equation}
\label{eq:7-75}
2\int \delta\varpi\!\left(\frac{d}{dr}(r\xi_r)
- \frac{2m\Omega}{\sigma}\,\xi_r\right) dr.
\end{equation}
The first-order change in $k^2$ resulting from the variation in $\varpi$ is, therefore,
\begin{equation}
\label{eq:7-76}
\delta k^2 = \frac{2}{I_2}\int \delta\varpi\!\left(
\frac{1}{r}\frac{d}{dr}(r\xi_r)
- \frac{2m\Omega}{\sigma r}\,\xi_r
- \frac{1}{\sigma^2}\!\left(\frac{m^2}{r^2}+k^2\right)\varpi
\right)r\,dr.
\end{equation}
Consequently, $\delta k^2 = 0$ for an arbitrary variation $\delta\varpi$ compatible with
the boundary conditions~\eqref{eq:7-68} if and only if
\begin{equation}
\label{eq:7-77}
\frac{1}{r}\frac{d}{dr}(r\xi_r)
- \frac{2m\Omega}{\sigma r}\,\xi_r
= \frac{1}{\sigma^2}\!\left(\frac{m^2}{r^2}+k^2\right)\varpi.
\end{equation}
Conversely, whenever equation~\eqref{eq:7-77} is satisfied, $\delta k^2 = 0$.  But equation~\eqref{eq:7-77} is the same as equation~\eqref{eq:7-54}.  Therefore, solving the characteristic
value problem presented by equations~\eqref{eq:7-52}, \eqref{eq:7-54}, and~\eqref{eq:7-55} is equivalent
to finding a maximum or a minimum of
\begin{equation}
\label{eq:7-78}
I_1 = -\int \frac{1}{r^2 + \frac{m^2}{r^2}}\!\left[
\!\left(\frac{d\varpi}{dr}+\frac{2m\Omega}{\sigma r}\,\varpi\right)^{\!2}
+ \frac{m^2\varpi^2}{r^2\sigma^2}\right]\,dr,
\end{equation}
for given
\begin{equation}
\label{eq:7-79}
I_2 = \int dr\; r\varpi^2/\sigma^2,
\end{equation}
for arbitrary variations of $\varpi$ subject only to the boundary conditions
\begin{equation}
\label{eq:7-80}
\frac{d\varpi}{dr}+\frac{2m\Omega}{\sigma}\,\varpi = 0
\quad \text{for } r = R_1 \text{ and } R_2;
\end{equation}
and the ratio $I_1/I_2$ at such a maximum or a minimum is a characteristic
value of $k^2$.

From equation~\eqref{eq:7-66} it is at once apparent that if $\Phi(r)$ is everywhere
negative,\footnote{As is the case when the flow is that permissible for viscid fluids and $0<\mu<\eta^2$
(cf.\ equations~\eqref{eq:7-27a} and~\eqref{eq:7-28}).}
$k^2$ cannot admit a positive characteristic value.  For real $p$'s,
the characteristic values of $k$ are necessarily imaginary.  Therefore
for real $k$'s, $p$ must necessarily be complex; and this means instability.

Consider next the case when $\Phi(r)$ changes sign in the interval $(R_1, R_2)$
and it is negative over a part of the interval and positive over others.
Under these circumstances, $I_1$ can assume a negative value for any
arbitrarily assigned real $p$; and for a given $I_2$, $I_1$ can attain a negative
minimum.  Therefore, for any real $p$, there exists at least one negative
characteristic value for $k^2$.  At the same time, since $\Phi(r)$ is positive over
some other parts of the interval $I_1$ can also assume positive values for
suitably chosen real values of $p$.  For a given $I_2$, there exist values of
$p$ for which $I_1$ can attain a positive maximum.  Therefore, for some
real values of $p$, there exist positive characteristic values for $k^2$.  From
our earlier arguments, these same values of $p$ must also allow negative
characteristic values of $k^2$.  Thus, in this case the dispersion relation
has two branches.  One of them is such that for all real $p$'s the characteristic values of $k^2$ are negative.  The inverse of this latter branch must
lead to complex $p$'s for real $k$'s; and this means instability.

Consider finally the case when $\Phi(r)$ is everywhere positive.  In this
case it is convenient to rewrite equation~\eqref{eq:7-66} in the form
\begin{equation}
\label{eq:7-81}
k^2 = \frac{\displaystyle\int
r\{[(d\varpi/dr)+(2m\Omega/\sigma)\varpi]^2(\Phi-\sigma^2)
- (m^2\varpi^2/r^2\sigma^2)\}\,dr}
{\displaystyle\int r\,\varpi^2/\sigma^2\,dr}.
\end{equation}
It is clear that for all real $p$ such that $\sigma^2 = (p+m\Omega)^2 > \Phi(r)$, the characteristic values of $k^2$ are negative.  For positive characteristic values it is
clearly necessary that $\sigma^2$ be less than $\Phi(r)$ at least over a part of the
interval $(R_1, R_2)$.  Indeed, it is clear that by requiring $\sigma^2$ to be less than
$\Phi(r)$ over a part of the interval, we can obtain for $k^2$ arbitrary positive
values; i.e.\ for all real $k$'s there correspond real $p$'s.  The system is
therefore stable.  The necessary and sufficient character of Rayleigh's
criterion for the stability of inviscid rotational flow is therefore established for arbitrary perturbations.

\section{The periods of oscillation of a rotating column of liquid}
\label{sec:7-68}
Rayleigh's criterion, for the stability of a given distribution of angular
velocity in a cylindrical column of liquid rotating about its axis, was
established in \S~67 from a consideration of the characteristic value problem which relates the frequency of oscillation with the wavelength of
the disturbance in the axial direction.  It is of interest to obtain the
explicit form of this relation in some special cases.  As Lord Kelvin has
remarked in his first investigation of the subject, `Crowds of extremely
interesting cases present themselves'.

\subsection*{(a) \emph{The case $\Omega = \text{constant}$}}
When $\Omega = \text{constant}$,
\begin{equation}
\label{eq:7-82}
\Phi(r) = \frac{1}{r^3}\frac{d}{dr}(r^2\Omega)^2 = 4\Omega^2,
\end{equation}
and the relevant equations are (cf.\ equations~\eqref{eq:7-52} and~\eqref{eq:7-54}):
\begin{equation}
\label{eq:7-83}
(\sigma^2 - 4\Omega^2)\xi_r = \frac{d\varpi}{dr}
+ \frac{2m\Omega}{\sigma r}\,\varpi,
\end{equation}
and
\begin{equation}
\label{eq:7-84}
\frac{1}{r}\frac{d}{dr}(r\xi_r)
- \frac{2m\Omega}{\sigma}\,\xi_r
= \frac{1}{\sigma^2}\!\left(\frac{m^2}{r^2}+k^2\right)\varpi,
\end{equation}
where
\begin{equation}
\label{eq:7-85}
\sigma = p + m\Omega
\end{equation}
is now a constant.  Eliminating $\xi_r$ between equations~\eqref{eq:7-83} and~\eqref{eq:7-84}, we
obtain
\begin{equation}
\label{eq:7-86}
\frac{d^2\varpi}{dr^2}+\frac{1}{r}\frac{d\varpi}{dr}
- \frac{m^2}{r^2}\,\varpi
= -k^2\!\left(\frac{4\Omega^2}{\sigma^2}-1\right)\varpi.
\end{equation}
It is convenient to measure $r$ in units of the radius $R_2$ of the outer
cylinder and let
\begin{equation}
\label{eq:7-87}
a = kR_2 \quad \text{and} \quad R_1 = R_2\eta.
\end{equation}
Equation~\eqref{eq:7-86} then becomes
\begin{equation}
\label{eq:7-88}
\frac{d^2\varpi}{dr^2}+\frac{1}{r}\frac{d\varpi}{dr}
+ \left(a^2\!\left(\frac{4\Omega^2}{\sigma^2}-1\right)
- \frac{m^2}{r^2}\right)\varpi = 0;
\end{equation}
and the boundary conditions are (cf.\ equation~\eqref{eq:7-68})
\begin{equation}
\label{eq:7-89}
\frac{d\varpi}{dr}+\frac{2m\Omega}{\sigma}\,\varpi = 0
\quad \text{for } r = 1 \text{ and } \eta.
\end{equation}
The general solution of equation~\eqref{eq:7-88} is
\begin{equation}
\label{eq:7-90}
\varpi = A J_m(\alpha r) + B Y_m(\alpha r),
\end{equation}
where
\begin{equation}
\label{eq:7-91}
\alpha = a(4\Omega^2/\sigma^2-1)^{1/2},
\end{equation}
$A$ and $B$ are constants of integration, and $J_m$ and $Y_m$ are the Bessel functions of the two kinds of order $m$.  The application of the boundary
conditions~\eqref{eq:7-89} to the solution~\eqref{eq:7-90} leads to a transcendental equation
relating $\alpha$ and $\sigma$ ($= p+m\Omega$).  We shall consider two special cases when
the relation is fairly simple.

\subsubsection*{(i) \emph{The case $\eta = 0$}}
When there is no inner cylinder and $\eta = 0$, the freedom from
singularity at $r = 0$ requires that the term in $Y_m$ in the solution~\eqref{eq:7-90} be
absent.  We then have
\begin{equation}
\label{eq:7-92}
\varpi = A J_m(\alpha r);
\end{equation}
and the boundary condition at $r = 1$ gives
\begin{equation}
\label{eq:7-93}
\alpha J_m'(\alpha) + \frac{2m\Omega}{\sigma}\, J_m(\alpha) = 0,
\end{equation}
where a prime denotes differentiation with respect to the argument.
But according to equation~\eqref{eq:7-91},
\begin{equation}
\label{eq:7-94}
\frac{\sigma}{2\Omega} = \frac{\pm 1}{\sqrt{1+\alpha^2/a^2}}.
\end{equation}
Using this relation to eliminate $\Omega/\sigma$ in equation~\eqref{eq:7-93}, we obtain
\begin{equation}
\label{eq:7-95}
\alpha J_m'(\alpha) \pm m(1+\alpha^2/a^2)^{1/2} J_m(\alpha) = 0.
\end{equation}
For any assigned value of $\alpha$ and $m$, equation~\eqref{eq:7-95} determines $\alpha/a$; equation~\eqref{eq:7-94} then determines $p$; thus
\begin{equation}
\label{eq:7-96}
\frac{p}{\Omega} = \frac{\sigma}{\Omega} - m = \pm\frac{2}{\sqrt{1+\alpha^2/a^2}} - m.
\end{equation}
In this way the dispersion relation between $\alpha$ and $p$ can be established.

The case $m = 0$ is particularly simple.  In this case $\alpha$ must be a zero
of $J_0'(x)$, i.e.\ of $J_1(x)$.  If $\alpha_{0,j}$ denotes the $j$th zero of $J_1(x)$, the explicit
form of the dispersion relation is

\figurePlaceholder{63}{The characteristic frequencies of oscillation of a rotating column of
fluid of radius $R_2$ for various values of $m$ and $j$; $a$ measures the wave number
in the unit $1/R_2$.}

\begin{equation}
\label{eq:7-97}
p = \pm\frac{2\Omega}{\sqrt{1+\alpha_{0,j}^2/a^2}}.
\end{equation}
The dispersion relations for $m = 0$, 1, 2, and 3 are illustrated in Fig.~63.
If waves of the type we have been considering should occur as standing
waves in a rotating column of liquid, then their half-wavelengths should
divide the height $H$ of the column an integral number of times; and we
must have
\begin{equation}
\label{eq:7-98}
a = n\pi R_2/H,
\end{equation}
where $n$ is an integer.  By inserting this value of $a$ in equations~\eqref{eq:7-96} and~\eqref{eq:7-97}, we shall obtain the frequencies of the various natural modes of
oscillation of the column.

In some recent experiments, Fultz has shown how these natural modes
of oscillation of a rotating column of liquid can be excited and maintained.  In these experiments the modes $m = 0$ in a rotating cylinder of
water (of diameter 4.2~cm and height 15.3~cm) were excited by means of
a small disk on the axis of the cylinder.  This disk could be moved up
and down the axis with a frequency which could be controlled.  To excite
a particular mode, the disk was so placed that its mean height was at the
position of the topmost antinode of the mode (e.g.\ at the middle of the
column to excite the fundamental mode).  Fultz was able to detect with
great sensitivity whether the disk was set at the resonant frequency by
following the motions of the fluid by means of dye suitably introduced.
In this way Fultz determined the frequencies of some of the lower modes
of oscillation belonging to $m = 0$; and he found that the observed
resonant frequencies agreed remarkably well with those given by
equations~\eqref{eq:7-97} and~\eqref{eq:7-98}.  In Fig.~64 we reproduce some of Fultz's photographs which exhibit in a very striking manner these natural modes of
oscillation of a rotating fluid.

\figurePlaceholder{64}{Photographs illustrating the oscillations of a rotating column of fluid in the modes $m=0$, $j=1$ (photographs 1 and 2), $m=0$, $j=2$ (photographs 3, 4, 5, and 6), and $m=0$, $j=3$ (photographs 7 and 8).  The data for the different photographs are as follows: Photographs 1,2: Rotation period 2.295~sec, observed $1/p$ = 1.316, theoretical 1.318, interval between photographs 1.25~sec. Photographs 3,4: 3.947, 0.788, 0.7886, 1.00. Photographs 5,6: 4.040, 0.789, 0.7886, 2.00. Photographs 7,8: 4.058, 0.646, 0.6444, 1.25. (The height of the column is 8.25~cm in all cases.)}

\subsubsection*{(ii) \emph{The case $m = 0$}}
When $m = 0$, the equation governing $\xi_r$ is (cf.\ equation~\eqref{eq:7-58})
\begin{equation}
\label{eq:7-99}
\frac{d^2\xi_r}{dr^2}+\frac{1}{r}\frac{d\xi_r}{dr}
+ \left(a^2\!\left(\frac{4\Omega^2}{p^2}-1\right)
- \frac{1}{r^2}\right)\xi_r = 0.
\end{equation}
The solution of this equation which vanishes at $r = 1$ and $\eta$ can be
written in the form
\begin{equation}
\label{eq:7-100}
\xi_r = \text{constant}\{Y_1(\alpha\eta)J_1(\alpha r) - J_1(\alpha\eta)Y_1(\alpha r)\},
\end{equation}
where $\alpha$ is a root of the equation
\begin{equation}
\label{eq:7-101}
Y_1(\alpha\eta)J_1(\alpha) - J_1(\alpha\eta)Y_1(\alpha) = 0.
\end{equation}
If $\alpha_j$ ($j = 1, 2, \ldots$) denote the roots of this equation, then
\begin{equation}
\label{eq:7-102}
a^2(4\Omega^2/p^2-1) = \alpha_j^2
\end{equation}
or
\begin{equation}
\label{eq:7-103}
p = \pm\frac{2\Omega}{\sqrt{1+\alpha_j^2/a^2}}.
\end{equation}

The values of $\alpha_j$ for a few values of $\eta$ are given in Table~XXXI.  For
$\eta = 0$ and $0.5$, the values of $\alpha_2$ and $\alpha_3$ are also given.

\subsection*{(b) \emph{The case $\Omega = A+B/r^4$ and $m=0$}}
A further case of some interest is when the distribution of the angular
velocity is that which obtains in a viscous liquid.  In that case (cf.\
equation~\eqref{eq:7-20}),
\begin{equation}
\label{eq:7-104}
\Phi = 4A(A+B/r^4);
\end{equation}

\begin{table}[htbp]
\centering
\caption{The roots $\alpha_j$ of equation~\eqref{eq:7-101} for some values of $\eta$}
\label{tab:XXXI}
\begin{tabular}{cccc}
\hline
$\eta$ & $\alpha_1$ & $\alpha_2$ & $\alpha_3$ \\
\hline
$0{\cdot}0$ & $3{\cdot}8317$ & $7{\cdot}01559$ & $10{\cdot}17347$ \\
$0{\cdot}2$ & $4{\cdot}2357$ & & \\
$0{\cdot}3$ & $4{\cdot}79578$ & & \\
$0{\cdot}4$ & $5{\cdot}30116$ & & \\
$0{\cdot}5$ & $6{\cdot}49316$ & $12{\cdot}64470$ & $18{\cdot}88893$ \\
$0{\cdot}8$ & $7{\cdot}93999$ & & \\
$0{\cdot}95$ & $15{\cdot}73587$ & & \\
\hline
\end{tabular}
\end{table}

and equation~\eqref{eq:7-58} governing $\xi_r$ (in case $m = 0$) can be written in the form
\begin{equation}
\label{eq:7-105}
(DD_* - k^2)\xi_r = -4\frac{k^2}{p^2}\,A\!\left(A+\frac{B}{r^4}\right)\xi_r,
\end{equation}
where
\begin{equation}
\label{eq:7-106}
D = \frac{d}{dr} \quad \text{and} \quad
D_* = \frac{d}{dr}+\frac{1}{r},
\end{equation}
and the boundary conditions are
\begin{equation}
\label{eq:7-107}
\xi_r = 0 \quad \text{for } r = R_1 \text{ and } R_2.
\end{equation}

\subsubsection*{(i) \emph{The solution for a narrow gap}}
Reid has recently shown how equation~\eqref{eq:7-105} can be solved explicitly
in terms of known, tabulated functions in the case when
\begin{equation}
\label{eq:7-108}
d = (R_2-R_1) \ll \tfrac{1}{2}(R_2+R_1).
\end{equation}
When~\eqref{eq:7-108} obtains, we need not distinguish between $D$ and $D_*$ and we
can further replace $A+B/r^4$ which occurs on the right-hand side of
equation~\eqref{eq:7-105} by
\begin{equation}
\label{eq:7-109}
\Omega_1\!\left[1-(1-\mu)\frac{r-R_1}{R_2-R_1}\right].
\end{equation}

In rewriting equation~\eqref{eq:7-105} in the framework of the foregoing approximations, we shall find it convenient to measure radial distances from the
surface of the inner cylinder in the unit $d = R_2-R_1$.  Thus, letting
\begin{equation}
\label{eq:7-110}
\zeta = \frac{r-R_1}{d} \quad \text{and} \quad a = k(R_2-R_1),
\end{equation}
we have to solve
\begin{equation}
\label{eq:7-111}
(D^2-a^2)\xi_r = -a^2\lambda[1-(1-\mu)\zeta]\xi_r,
\end{equation}
where $D$ now stands for $d/d\zeta$ and (cf.\ equations~\eqref{eq:7-3})
\begin{equation}
\label{eq:7-112}
\lambda = \frac{4A\Omega_1}{p^2} = \frac{4\Omega_1^2}{p^2}\,\eta^4\,
\frac{1-\mu/\eta^2}{1-\eta^2},
\end{equation}
and the boundary conditions are
\begin{equation}
\label{eq:7-113}
\xi_r = 0 \quad \text{for } \zeta = 0 \text{ and } 1.
\end{equation}

In the case $\mu = 1$, equation~\eqref{eq:7-111} admits of elementary integration
and the proper solutions which satisfy the boundary conditions~\eqref{eq:7-113} are
given by
\begin{equation}
\label{eq:7-114}
\xi_r = \text{constant}\;\sin n\pi\zeta,
\end{equation}
where $n$ is an integer.  The corresponding characteristic equation is
\begin{equation}
\label{eq:7-115}
\frac{1}{\lambda} = \frac{a^2}{a^2+n^2\pi^2},
\end{equation}
or, since $\mu = 1$,
\begin{equation}
\label{eq:7-116}
p = \frac{2\Omega_1}{\sqrt{1+n^2\pi^2/a^2}};
\end{equation}
this is clearly the asymptotic form of~\eqref{eq:7-103} for $\eta \to 1$.

For $\mu < 1$, we can reduce equation~\eqref{eq:7-111} to the standard form
\begin{equation}
\label{eq:7-117}
\frac{d^2\xi_r}{dx^2} = x\xi_r,
\end{equation}
by the substitution
\begin{equation}
\label{eq:7-118}
x = \left[\frac{a}{\lambda(1-\mu)}\right]^{2/3}\!
\{\lambda - \lambda[1-(1-\mu)\zeta]\}.
\end{equation}
The boundary conditions, then, are
\begin{equation}
\label{eq:7-119}
\xi_r = 0 \quad \text{for } x = x_1 \text{ and } x = x_2,
\end{equation}
where
\begin{equation}
\label{eq:7-120}
x_1 = \left[\frac{a}{\lambda(1-\mu)}\right]^{2/3}\!(1-\lambda)
\quad \text{and} \quad
x_2 = \left[\frac{a}{\lambda(1-\mu)}\right]^{2/3}\!(1-\lambda\mu).
\end{equation}
The general solution of equation~\eqref{eq:7-117} can be expressed in terms of
Bessel functions of order $\frac{1}{3}$, or more conveniently in terms of Airy's
functions\footnote{The definitions of these functions are
\[
\mathrm{Ai}(x) = \tfrac{1}{3}x^{1/2}\{I_{-1/3}(y) - I_{1/3}(y)\}
\]
and
\[
\mathrm{Bi}(x) = (x/3)^{1/2}\{I_{-1/3}(y) + I_{1/3}(y)\},
\]
where $y = \frac{2}{3}x^{3/2}$.} $\mathrm{Ai}(x)$ and $\mathrm{Bi}(x)$; thus
\begin{equation}
\label{eq:7-121}
\xi_r = A\,\mathrm{Ai}(x) + B\,\mathrm{Bi}(x),
\end{equation}
where $A$ and $B$ are constants.  The boundary conditions~\eqref{eq:7-119} lead to
the characteristic equation
\begin{equation}
\label{eq:7-122}
\frac{\mathrm{Ai}(x_1)}{\mathrm{Bi}(x_1)} = \frac{\mathrm{Ai}(x_2)}{\mathrm{Bi}(x_2)}.
\end{equation}
Equation~\eqref{eq:7-122} determines $x_2$ as a function of $x_1$ for the different modes;
the relationship for the lowest mode is illustrated in Fig.~65.

\figurePlaceholder{65}{The $(x_2,x_1)$-relation for the lowest mode
as determined by equation~\eqref{eq:7-122}.}

$(x_2,x_1)$-relationship determined in this fashion, the required relation
between $a$ and $\sqrt{\lambda}$ is given, parametrically, by
\begin{equation}
\label{eq:7-123}
a = (x_2-x_1)\sqrt{\frac{x_2-\mu x_1}{1-\mu}}
\quad \text{and} \quad
\frac{1}{\sqrt{\lambda}} = \sqrt{\frac{x_2-\mu x_1}{x_2-x_1}}.
\end{equation}
The dispersion relations for $\mu = 1$, 0, and $-1$ are illustrated in Fig.~66.
In Fig.~67\,$a$ the limiting form of velocity profile for the lowest mode of
instability when $\mu \to -\infty$ is illustrated; and the corresponding cell
pattern is shown in Fig.~67\,$b$.

\figurePlaceholder{66}{The dispersion relations in the limit of zero viscosity
for $\mu=1$, 0, and $-1$ for the lowest modes of the flow between two cylinders with a narrow gap $d$.  The abscissa gives
the wave number in the unit $1/d$ and the ordinate gives the
growth rate (or the circular frequency) in the unit $\sqrt{4A\Omega_1}$.}

\figurePlaceholder{67}{The limiting form of the velocity profile (a) and the cell pattern (b) for
lowest inviscid mode of instability as $\mu\to-\infty$ and $a/(1-\mu) = 2{\cdot}03$.}

\subsubsection*{(ii) \emph{The formal solution for a wide gap}}
When we wish to allow for the finite difference between $R_1$ and $R_2$,
we must consider equation~\eqref{eq:7-105} without making any approximations.

Now measuring $r$ in units of $R_2$ and letting $a = kR_2$, we have the
equation (cf.\ equations~\eqref{eq:7-22} and~\eqref{eq:7-23})
\begin{equation}
\label{eq:7-124}
(DD_*-a^2)\xi_r = a^2\frac{4\Omega_1^2}{p^2}\,\eta^4\,
\frac{(1-\mu)(1-\mu/\eta^2)}{(1-\eta^2)^2}
\!\left(\frac{1}{r^4}-\kappa\right)\xi_r.
\end{equation}

Letting
\begin{equation}
\label{eq:7-125}
Q_1^2 = -\frac{4\Omega_1^2}{p^2}\,\eta^4\,
\frac{(1-\mu)(1-\mu/\eta^2)}{(1-\eta^2)^2}
\end{equation}
and
\begin{equation}
\label{eq:7-126}
Q_1^2\kappa = -Q_1^2\,\kappa,
\end{equation}
we can rewrite equation~\eqref{eq:7-124} in the form
\begin{equation}
\label{eq:7-127}
\frac{d^2\xi_r}{dr^2}+\frac{1}{r}\frac{d\xi_r}{dr}
+ \left(a^2 Q_1^2 - 1\right)\frac{1-a^2 Q_1^2}{r^4}\,\xi_r = 0.
\end{equation}
The general solution of this equation is
\begin{equation}
\label{eq:7-128}
\xi_r = \mathscr{C}_\nu(ar),
\end{equation}
where
\begin{equation}
\label{eq:7-129}
\nu = \sqrt{1-a^2 Q_1^2}
\quad \text{and} \quad
\alpha = a\sqrt{Q_1^2-1},
\end{equation}
and $\mathscr{C}_\nu$ represents a general cylinder function of order $\nu$.

The boundary conditions require $\mathscr{C}_\nu(\alpha r)$ to vanish at $r = 1$ and $\eta$;
and these conditions will determine $\alpha$.  The required solution can, in fact,
be expressed in the form
\begin{equation}
\label{eq:7-130}
\xi_r = \text{constant}\{J_{-\nu}(\alpha\eta)J_\nu(\alpha r)
- J_\nu(\alpha\eta)J_{-\nu}(\alpha r)\}.
\end{equation}
Defined in this manner, $\xi_r$ clearly vanishes at $r = \eta$; and the condition
that it also vanishes at $r = 1$ leads to the equation
\begin{equation}
\label{eq:7-131}
J_{-\nu}(\alpha\eta)J_\nu(\alpha)
- J_\nu(\alpha\eta)J_{-\nu}(\alpha) = 0.
\end{equation}

With $\alpha$ determined as a root of this equation for some assigned $\nu$, the
required relation between $a^2$ and $p^2$ follows from equations~\eqref{eq:7-126} and~\eqref{eq:7-129}.

From the general theory of \S\S~66 and 67, we know that
\begin{equation}
\label{eq:7-132}
p^2 > 0 \text{ for } \mu > \eta^2 \quad \text{and} \quad
p^2 < 0 \text{ for } \mu < \eta^2.
\end{equation}
Therefore,
\begin{equation}
\label{eq:7-133}
Q_1^2 > 0 \quad \text{or} \quad < 0 \quad \text{according as} \quad
\mu < 1 \quad \text{or} \quad > 1.
\end{equation}
But $Q_1^2$ is positive or negative according as $p^2$ is positive or negative,
i.e.\ according as $\mu > \eta^2$ or $\mu < \eta^2$.  Thus, for the unstable modes, $Q_1^2 < 0$
and $\alpha$ becomes imaginary.  If $\nu$ were real, $\mathscr{C}_\nu(\alpha r)$ would become a linear
combination of $I_\nu(|\alpha|r)$ and $K_\nu(|\alpha|r)$.  But no linear combination of
$I_\nu(x)$ and $K_\nu(x)$ can vanish at two specified points.  Therefore, for the
unstable modes, $\nu$ must be imaginary and the cylinder functions, in terms
of which the solutions are expressed, are of imaginary order and argument.

\section{On viscous Couette flow}\label{sec:7-69}
We now turn our attention to stationary viscous flow between two
rotating coaxial cylinders.

In cylindrical polar coordinates, the Navier--Stokes equations for
viscous incompressible fluids take the forms:
\begin{equation}
\label{eq:7-134}
\frac{\partial u_r}{\partial t}
+ (\mathbf{u}.\mathrm{grad})\,u_r - \frac{u_\theta^2}{r}
= -\frac{\partial}{\partial r}\!\left(\frac{p}{\rho}\right)
+ \nu\!\left(\nabla^2 u_r - \frac{2}{r^2}\frac{\partial u_\theta}{\partial \theta}
- \frac{u_r}{r^2}\right),
\end{equation}
\begin{equation}
\label{eq:7-135}
\frac{\partial u_\theta}{\partial t}
+ (\mathbf{u}.\mathrm{grad})\,u_\theta + \frac{u_r u_\theta}{r}
= -\frac{1}{r}\frac{\partial}{\partial r}\!\left(\frac{p}{\rho}\right)
+ \nu\!\left(\nabla^2 u_\theta + \frac{2}{r^2}\frac{\partial u_r}{\partial \theta}
- \frac{u_\theta}{r^2}\right),
\end{equation}
\begin{equation}
\label{eq:7-136}
\frac{\partial u_z}{\partial t}
+ (\mathbf{u}.\mathrm{grad})\,u_z
= -\frac{\partial}{\partial z}\!\left(\frac{p}{\rho}\right)
+ \nu \nabla^2 u_z,
\end{equation}
where
\begin{equation}
\label{eq:7-137}
\mathbf{u}.\mathrm{grad} = u_r \frac{\partial}{\partial r}
+ \frac{u_\theta}{r}\frac{\partial}{\partial \theta}
+ u_z \frac{\partial}{\partial z}
\end{equation}
and
\begin{equation}
\label{eq:7-138}
\nabla^2 = \frac{\partial^2}{\partial r^2}
+ \frac{1}{r}\frac{\partial}{\partial r}
+ \frac{1}{r^2}\frac{\partial^2}{\partial \theta^2}
+ \frac{\partial^2}{\partial z^2}.
\end{equation}
We have also the equation of continuity,
\begin{equation}
\label{eq:7-139}
\frac{\partial u_r}{\partial r}
+ \frac{u_r}{r}
+ \frac{1}{r}\frac{\partial u_\theta}{\partial \theta}
+ \frac{\partial u_z}{\partial z} = 0.
\end{equation}
These equations allow a stationary solution of the form
\begin{equation}
\label{eq:7-140}
u_r = u_z = 0 \quad \text{and} \quad u_\theta = V(r),
\end{equation}
provided
\begin{equation}
\label{eq:7-141}
\frac{d}{dr}\!\left(\frac{p}{\rho}\right) = \frac{V^2}{r}
\end{equation}
and
\begin{equation}
\label{eq:7-142}
\nu\!\left(\nabla^2 V - \frac{V}{r^2}\right)
= \nu\,\frac{d}{dr}\!\left(\frac{d}{dr}+\frac{1}{r}\right)V = 0.
\end{equation}
From equation~\eqref{eq:7-142}, it is apparent that the most general form of $V(r)$
which is compatible with $\nu \neq 0$ is
\begin{equation}
\label{eq:7-143}
V = Ar + \frac{B}{r},
\end{equation}
where $A$ and $B$ are two arbitrary constants.  The corresponding expression for the angular velocity is
\begin{equation}
\label{eq:7-144}
\Omega = A + \frac{B}{r^2}.
\end{equation}
This is the solution which was quoted in \S~65.  And as we remarked at
that time, the appearance of two arbitrary constants in the solution~\eqref{eq:7-144}
corresponds to the possibility of assigning to the two cylinders, confining
the fluid, angular velocities at will.  The constants $A$ and $B$ can, therefore, be related to the angular velocities $\Omega_1$ and $\Omega_2$ with which the two
cylinders are rotated.  We have (cf.\ equations~\eqref{eq:7-3} and~\eqref{eq:7-4})
\begin{equation}
\label{eq:7-145}
A = -\Omega_1\,\eta^2\,\frac{1-\mu/\eta^2}{1-\eta^2}
\quad \text{and} \quad
B = \Omega_1\,\frac{R_1^2(1-\mu)}{1-\eta^2},
\end{equation}
where
\begin{equation}
\label{eq:7-146}
\mu = \Omega_2/\Omega_1 \quad \text{and} \quad \eta = R_1/R_2.
\end{equation}
As we have seen in \S~66, Rayleigh's criterion for the stability of inviscid
Couette flow applied to the distribution~\eqref{eq:7-144} yields
\begin{equation}
\label{eq:7-147}
\mu > \eta^2 \quad \text{for stability.}
\end{equation}
On general grounds, we should expect the effect of viscosity is to
postpone the onset of instability; and that the new criterion for stability
must essentially take the form of specifying the extent to which the
condition~\eqref{eq:7-147} can be violated before instability will manifest itself.

The precise form which the criterion takes was first investigated, both
theoretically and experimentally, by G.~I.~Taylor.  In the case when the
gap $R_2-R_1$ between the two cylinders can be considered as small
compared to the mean radius $\frac{1}{2}(R_1+R_2)$, Taylor found an explicit
analytical expression for the criterion; and he was able to confirm by
experiments that the marginal state is stationary and exhibits a break-up
of the basic flow into a cellular pattern.  Fig.~68, which is a reproduction of one of Taylor's photographs, clearly shows the manner of the
onset of instability.  In \S~74 we return to a more detailed account of the
experimental work which has been carried out since Taylor's early
pioneering work.

\figurePlaceholder{68}{One of G.~I.~Taylor's photographs illustrating the onset of instability
in the flow between rotating cylinders.  The radii of the two cylinders in this
case are 4.035~cm and 3.25~cm respectively; the cylinders are rotating in the
same direction.}

\section{The perturbation equations}\label{sec:7-70}
We shall now investigate the stability of the flow described by equations~\eqref{eq:7-141} and~\eqref{eq:7-143}.  Let the perturbed state be characterized by
\begin{equation}
\label{eq:7-148}
u_r, \quad V\!+\!u_\theta, \quad u_z, \quad \text{and} \quad
\delta p/\rho = \varpi.
\end{equation}
Assuming that the various perturbations are axisymmetric\footnote{The general non-axisymmetric case has not been investigated for this problem.} and independent of $\theta$, we obtain from~\eqref{eq:7-134}--\eqref{eq:7-136} the linearized equations
\begin{equation}
\label{eq:7-149}
\frac{\partial u_r}{\partial t}
- \frac{2V}{r}\, u_\theta
= -\frac{\partial \varpi}{\partial r}
+ \nu\!\left(\nabla^2 u_r - \frac{u_r}{r^2}\right),
\end{equation}
\begin{equation}
\label{eq:7-150}
\frac{\partial u_\theta}{\partial t}
+ \left(\frac{dV}{dr}+\frac{V}{r}\right) u_r
= \nu\!\left(\nabla^2 u_\theta - \frac{u_\theta}{r^2}\right),
\end{equation}
and
\begin{equation}
\label{eq:7-151}
\frac{\partial u_z}{\partial t}
= -\frac{\partial \varpi}{\partial z}
+ \nu \nabla^2 u_z,
\end{equation}
where $\nabla^2$ has now the meaning
\begin{equation}
\label{eq:7-152}
\nabla^2 = \frac{\partial^2}{\partial r^2}
+ \frac{1}{r}\frac{\partial}{\partial r}
+ \frac{\partial^2}{\partial z^2}.
\end{equation}
Also, for axisymmetric motions, the equation of continuity reduces to
\begin{equation}
\label{eq:7-153}
\frac{\partial u_r}{\partial r}
+ \frac{u_r}{r}
+ \frac{\partial u_z}{\partial z} = 0.
\end{equation}

By analysing the disturbance into normal modes, we seek solutions
of the foregoing equations which are of the forms
\begin{equation}
\label{eq:7-154}
\left.\begin{aligned}
u_r &= e^{pt}u(r)\cos kz;\quad &
u_\theta &= e^{pt}v(r)\sin kz\\
u_z &= e^{pt}w(r)\cos kz;\quad &
\varpi &= e^{pt}\varpi(r)\cos kz
\end{aligned}\right\},
\end{equation}
where $k$ is the wave number of the disturbance in the axial direction
and $p$ is a constant which can be complex.  For solutions of the form~\eqref{eq:7-154}, equations~\eqref{eq:7-149}--\eqref{eq:7-153} become

\begin{equation}
\label{eq:7-155}
\nu\!\left(DD_*-k^2-\frac{p}{\nu}\right)u
+ 2\frac{V}{r}\,v = \frac{d\varpi}{dr},
\end{equation}
\begin{equation}
\label{eq:7-156}
\nu\!\left(D_* D - k^2 - \frac{p}{\nu}\right)v - (D_*V)\,u = 0,
\end{equation}
\begin{equation}
\label{eq:7-157}
\nu\!\left(D_* D-k^2-\frac{p}{\nu}\right)w = -k\varpi,
\end{equation}
\begin{equation}
\label{eq:7-158}
\nabla^2 = \left(\frac{d}{dr}+\frac{1}{r}\right)\frac{1}{r}\frac{d}{dr}\, r - k^2
= D_* D - k^2,
\end{equation}
and
\begin{equation}
\label{eq:7-159}
D_* u = -kw.
\end{equation}
[In the foregoing equations, $D$ and $D_*$ have the same meanings as in
\S~68\,($b$), equation~\eqref{eq:7-106}.]

Eliminating $w$ between equations~\eqref{eq:7-157} and~\eqref{eq:7-159}, we have
\begin{equation}
\label{eq:7-160}
\frac{\nu}{k^2}\!\left(D_* D-k^2-\frac{p}{\nu}\right)D_*\,u = \varpi.
\end{equation}
Inserting this expression for $\varpi$ in equation~\eqref{eq:7-155}, we find after some
rearranging,
\begin{equation}
\label{eq:7-161}
\frac{\nu}{k^2}\!\left(DD_*-k^2-\frac{p}{\nu}\right)(DD_*-k^2)u
= 2\frac{V}{r}\,v.
\end{equation}
This equation must be considered together with
\begin{equation}
\label{eq:7-162}
\nu\!\left(DD_*-k^2-\frac{p}{\nu}\right)v = (D_*V)u.
\end{equation}
These equations are general and do not depend on any particular form
of $V(r)$.

Measuring $r$ in units of the radius $R_2$ of the outer cylinder and writing
\begin{equation}
\label{eq:7-163}
k^2 = a^2/R_2^2 \quad \text{and} \quad
\sigma = p R_2^2/\nu,
\end{equation}
equations~\eqref{eq:7-161} and~\eqref{eq:7-162} become (when $V(r)$ has the particular form~\eqref{eq:7-143}):
\begin{equation}
\label{eq:7-164}
(DD_*-a^2)(DD_*-a^2)u = a^2\frac{2B}{\nu}\!\left(\frac{1}{r^2}+\frac{A\,R_2^4}{B}\right)v,
\end{equation}
and
\begin{equation}
\label{eq:7-165}
(DD_*-a^2-\sigma)v = \frac{2A}{\nu}\,R_2^2\, u.
\end{equation}
It is convenient to make the transformation
\begin{equation}
\label{eq:7-166}
\frac{2A\,R_2^2}{\nu}\,u \to u;
\end{equation}
the equations then take the more convenient forms
\begin{equation}
\label{eq:7-167}
(DD_*-a^2-\sigma)(DD_*-a^2)u = -Ta^2\!\left(\frac{1}{r^2}-\kappa\right)v,
\end{equation}
and
\begin{equation}
\label{eq:7-168}
(DD_*-a^2-\sigma)v = u,
\end{equation}
where (cf.\ equations~\eqref{eq:7-21}, \eqref{eq:7-23}, and~\eqref{eq:7-32})
\begin{equation}
\label{eq:7-169}
T = -\frac{4AB}{R_2^4} = \frac{4\Omega_1^2\,R_2^4}{\nu^2}\,
\frac{(1-\mu)(1-\mu/\eta^2)}{(1-\eta^2)^2}
\end{equation}
and
\begin{equation}
\label{eq:7-170}
\kappa = -\frac{A\,R_2^4}{B} = \frac{1-\mu/\eta^2}{1-\mu}.
\end{equation}

Solutions of equations~\eqref{eq:7-167} and~\eqref{eq:7-168} must be sought which satisfy
the boundary conditions appropriate for no slip on the cylindrical walls
at $r = 1$ and $\eta$.  These conditions are that all three components of the
velocity vanish on the walls; thus,
\begin{equation}
\label{eq:7-171}
u = v = 0 \quad \text{and} \quad Du = 0
\quad \text{for } r = 1 \text{ and } \eta,
\end{equation}
where the last of the three boundary conditions is equivalent to $w = 0$
(cf.\ equation~\eqref{eq:7-159}).

\subsection*{(a) \emph{The stability of the flow for $\mu > \eta^2$}}
We shall now show that when Rayleigh's criterion $\mu > \eta^2$ is satisfied,
the flow is indeed stable.

First, we may notice certain elementary integral properties of the
operator $DD_*$.  If $f(r)$ and $g(r)$ are any two functions and if one of them,
say $f(r)$, vanishes at the limits of integration,
\begin{equation}
\label{eq:7-172}
\int f\, DD_*\, g\, dr
= -\int\!\left(r\frac{df}{dr}\frac{dg}{dr}+\frac{fg}{r}\right)dr,
\end{equation}
and if the derivative of $f$ also vanishes at the limits,
\begin{equation}
\label{eq:7-173}
\int f\, DD_*\, g\, dr = \int \eta\, DD_*\, f\, dr.
\end{equation}
These relations follow by successive integrations by parts.  Thus, by
writing
\begin{equation}
\label{eq:7-174}
\int f\, DD_*\, g\, dr
= \int\!\left\{
\frac{1}{r}\frac{d}{dr}\!\left(r\frac{dg}{dr}\right)
- \frac{g}{r}\right\} f\, dr,
\end{equation}
the truth of~\eqref{eq:7-172} becomes self-evident; and a further integration by parts
leads to~\eqref{eq:7-173}.

Now returning to equations~\eqref{eq:7-167} and~\eqref{eq:7-168}, multiply~\eqref{eq:7-167} by $ru^*$
and integrate over the range of $r$.  We have
\begin{equation}
\label{eq:7-175}
\int ru^*\{(DD_*-a^2)^2u - \sigma(DD_*-a^2)u\}\,dr
= -\mathscr{T}a^2 \int r\phi(r)vu^*\,dr,
\end{equation}
where, for brevity, we have written
\begin{equation}
\label{eq:7-176}
\mathscr{T} = \frac{T}{a^2} = \frac{4\Omega_1^2\,R_2^4\,(1-\mu/\eta^2)}{\nu^2\,(1-\eta^2)^2}
\quad \text{and} \quad
\phi = (1-\mu)\!\left(\frac{1}{r^2}-\kappa\right).
\end{equation}
Since $u$ and its derivative vanish at $r = 1$ and $\eta$, the integrals on the
left-hand side of~\eqref{eq:7-175} can be transformed to positive definite forms by
making use of~\eqref{eq:7-172} and~\eqref{eq:7-173}.  Thus,
\begin{equation}
\label{eq:7-177}
\int ru^*\{(DD_*-a^2)^2u - \sigma(DD_*-a^2)u\}\,dr
= \int r\{(DD_*-a^2)u\}^2\,dr
+ \left(\frac{1}{r}+a^2 r\right)|u|^2\,dr.
\end{equation}
(This is not quite right as written. The full expansion gives separate integral terms.)

Next, substituting for $u^*$ (from equation~\eqref{eq:7-168}) in the integrand on
the right-hand side of equation~\eqref{eq:7-175}, we obtain
\begin{equation}
\label{eq:7-178}
\int r\phi(r)vu^*\,dr
= \int r\phi(r) v(DD_*-a^2-\sigma^*)v^*\,dr
+ \int r\phi(r)vDD_*v^*\,dr.
\end{equation}
Again, by making use of~\eqref{eq:7-172}, we have
\begin{equation}
\label{eq:7-179}
\int r\phi(r)v DD_*v^*\,dr
= -\int \phi(r)\!\left(r\left|\frac{dv}{dr}\right|^2
+ \frac{|v|^2}{r}\right)dr
+ 2(1-\mu)\int \frac{v}{r}\frac{dv^*}{dr}\,dr.
\end{equation}
Now combining equations~\eqref{eq:7-175}, \eqref{eq:7-177}, \eqref{eq:7-178}, and~\eqref{eq:7-179}, we obtain
\begin{equation}
\label{eq:7-180}
\sigma I_1 + I_2 = \mathscr{T}a^2(a^2+\sigma^*)I_3 + I_4\},
\end{equation}
where
\begin{equation}
\label{eq:7-181}
I_1 = \int \phi(r)\!\left(r\left|\frac{du}{dr}\right|^2
+ \left(\frac{1}{r}+a^2 r\right)|u|^2\right)dr,
\end{equation}
\begin{equation}
\label{eq:7-182}
I_2 = \int |(DD_*-a^2)u|^2\,r\,dr,
\end{equation}
\begin{equation}
\label{eq:7-183}
I_3 = \int \phi(r)|v|^2\,dr,
\end{equation}
and
\begin{equation}
\label{eq:7-184}
I_4 = \int \phi(r)\!\left(r\left|\frac{dv}{dr}\right|^2
+ \frac{|v|^2}{r}\right)dr
- 2(1-\mu)\int \frac{v}{r}\frac{dv^*}{dr}\,dr.
\end{equation}
The integrals $I_2$ and $I_3$ are clearly positive definite.  For $\mu > 0$, $\phi(r) > 0$
(see equation~\eqref{eq:7-26}); so that in this case, $I_1$ is also positive definite.  The
first of the two integrals included in $I_4$ is positive definite for $\mu > 0$; but
the second is complex.  However, the real part of $I_4$ is positive definite
for $\mu > 0$; in fact,
\begin{equation}
\label{eq:7-185}
\mathrm{re}(I_4) = \int r\phi(r)\left|\frac{dv}{dr}
- \frac{v}{r}\right|^2 dr.
\end{equation}
For, expanding the integrand in~\eqref{eq:7-185}, we have
\begin{equation}
\label{eq:7-186}
\int r\phi(r)\left|\frac{dv}{dr}
- \frac{v}{r}\right|^2 dr
= \int \phi(r)\!\left(r\left|\frac{dv}{dr}\right|^2
+ \frac{|v|^2}{r}\right)dr
- \int \phi(r)\frac{d}{dr}|v|^2\,dr;
\end{equation}
but
\begin{equation}
\label{eq:7-187}
\int \phi(r)\frac{d}{dr}|v|^2\,dr
= (1-\mu)\int \frac{1}{r}\!\left(\frac{1}{r^2}-\kappa\right)|v|^2\,dr
= (1-\mu)\int \frac{1}{r^2}\frac{d|v|^2}{dr}\,dr.
\end{equation}
Therefore, the right-hand side of~\eqref{eq:7-186} is, indeed, the real part of $I_4$.
Returning to equation~\eqref{eq:7-180} and equating the real parts of this
equation, we obtain
\begin{equation}
\label{eq:7-188}
\mathrm{re}(\sigma)\{I_1-\mathscr{T}a^2 I_3+I_4-\mathscr{T}a^2 \sigma^* I_3 + \mathrm{re}(I_4)\} = 0.
\end{equation}
When $\mu > \eta^2$, $\mathscr{T} < 0$ and the coefficient of $\mathrm{re}(\sigma)$ in equation~\eqref{eq:7-188} is
positive definite; and so also are the remaining terms in the equation.
Therefore,
\begin{equation}
\label{eq:7-189}
\mathrm{re}(\sigma) < 0 \quad \text{for } \mu > \eta^2.
\end{equation}
and the flow is stable; this result is entirely to be expected on physical
grounds.  Nevertheless, it appears to be the only one which can be
established by general analytical arguments.  In particular, it does not
seem that one can deduce the general validity of the principle of the
exchange of stabilities for this problem.  For example, by equating the
imaginary parts of equation~\eqref{eq:7-180}, we obtain (cf.\ equation~\eqref{eq:7-176})
\begin{equation}
\label{eq:7-190}
\mathrm{im}(\sigma)(I_1+\mathscr{T}a^2 I_3)
= -2\mathscr{T}a^2\,\mathrm{im}\int\frac{v}{r}\frac{dv^*}{dr}\,dr,
\end{equation}
and no general conclusions can be drawn from this equation; when $\mu < 0$,
even $I_4$ is not positive definite!

\section{The solution for the case of a narrow gap when the marginal
state is stationary}\label{sec:7-71}
If the gap $R_2-R_1$ between the two cylinders is small compared to
their mean radius $\frac{1}{2}(R_1+R_2)$, we need not (as in \S~68\,($b$)) distinguish
between $D$ and $D_*$ in equations~\eqref{eq:7-161} and~\eqref{eq:7-162}; and we can also replace
$(A+B/r^4)$ which occurs on the right-hand side of equation~\eqref{eq:7-161} by
\begin{equation}
\label{eq:7-191}
\Omega_1\!\left[1-(1-\mu)\frac{r-R_1}{R_2-R_1}\right].
\end{equation}

In rewriting equations~\eqref{eq:7-161} and~\eqref{eq:7-162} in the framework of these
approximations, it will be convenient to measure radial distances from the
surface of the inner cylinder in the unit $d = R_2-R_1$.  Thus, letting
\begin{equation}
\label{eq:7-192}
\zeta = (r-R_1)/d, \quad k = a/d, \quad \text{and} \quad
\sigma = pd^2/\nu,
\end{equation}
we have to consider the equations
\begin{equation}
\label{eq:7-193}
(D^2-a^2-\sigma)(D^2-a^2)u = -Ta^2[1-(1-\mu)\zeta]v
\end{equation}
and
\begin{equation}
\label{eq:7-194}
(D^2-a^2-\sigma)v = \frac{2Ad^2}{\nu}\,u.
\end{equation}
By the further transformation
\begin{equation}
\label{eq:7-195}
u \to \frac{2\Omega_1\,d^2 a^2}{\nu}\,u,
\end{equation}
the equations become
\begin{equation}
\label{eq:7-196}
(D^2-a^2-\sigma)(D^2-a^2)u = (1+\alpha\zeta)v
\end{equation}
and
\begin{equation}
\label{eq:7-197}
(D^2-a^2-\sigma)v = -Ta^2 u,
\end{equation}
where, now,
\begin{equation}
\label{eq:7-198}
T = \frac{4A\Omega_1}{\nu^2}\,d^4
\end{equation}
and
\begin{equation}
\label{eq:7-199}
\alpha = -(1-\mu).
\end{equation}
Equations~\eqref{eq:7-196} and~\eqref{eq:7-197} must be considered together with the boundary
conditions
\begin{equation}
\label{eq:7-200}
u = Du = v = 0 \quad \text{for } \zeta = 0 \text{ and } 1.
\end{equation}
We are primarily interested in the solutions of equations~\eqref{eq:7-196} and~\eqref{eq:7-197} (subject to the boundary conditions~\eqref{eq:7-200}) for various values of $\alpha$
for which the real part of $\sigma$ is zero.  The method described in Chapter~III,
\S~31 in a different connexion is applicable to this problem.  Thus, we
must obtain solutions for two cases: when $\sigma$ is zero and the marginal
state is stationary and when $\sigma$ is imaginary and the marginal state is
oscillatory.  In the latter case for each value of $\alpha$, $\sigma$ must be determined
by the condition that $T$ is real.\footnote{It is, of course, possible that under certain circumstances solutions with this property
do not exist.}  In either case, we must find the minimum of $T$ as a function of $a$; and depending on which of the two minima
is lower, we shall have the onset of instability as a stationary secondary
flow or as overstability.  Careful experiments on the onset of instability
by Taylor and others have failed to reveal any suggestions of overstability.  For this reason, the case $\sigma = 0$ is the only one which has been
considered in the literature.  However, as no general arguments for the
validity of the principle of the exchange of stabilities have been found
for this problem, the case of overstability requires investigation.  We
return to this question in \S~72.

When the marginal state is stationary, the equations to be solved are
\begin{equation}
\label{eq:7-201}
(D^2-a^2)^2 u = (1+\alpha\zeta)v
\end{equation}
and
\begin{equation}
\label{eq:7-202}
(D^2-a^2)v = -Ta^2 u,
\end{equation}
together with the boundary conditions~\eqref{eq:7-200}.

\subsection*{(a) \emph{The solution of the characteristic value problem for the case
$\alpha = 0$}}
It can be readily verified that the characteristic value problem presented by equations~\eqref{eq:7-201}--\eqref{eq:7-202} is not self-adjoint in the usual sense.
For this reason, the method to be described below was patterned after the
ones which have been found successful in cases where the problems are
self-adjoint.  However, Roberts has recently found a variational basis
for the method; this is considered in Appendix~IV.

The method of solution we shall adopt is the following.

Since $v$ is required to vanish at $\zeta = 0$ and $1$, we expand it in a sine
series of the form
\begin{equation}
\label{eq:7-203}
v = \sum_{m=1}^{\infty} C_m \sin m\pi\zeta.
\end{equation}
Having chosen $v$ in this manner, we next \emph{solve} the equation,
\begin{equation}
\label{eq:7-204}
(D^2-a^2)^2 u = (1+\alpha\zeta)\sum_{m=1}^{\infty} C_m \sin m\pi\zeta,
\end{equation}
obtained by inserting~\eqref{eq:7-203} in~\eqref{eq:7-201}, and arrange that the solution
satisfies the four remaining boundary conditions on $u$.  With $u$ determined
in this fashion and $v$ given by~\eqref{eq:7-203}, equation~\eqref{eq:7-202} will lead, as we shall
presently see, to a secular equation for $T$.

The solution of equation~\eqref{eq:7-204} is straightforward.  The general
solution can be written in the form
\begin{multline}
\label{eq:7-205}
u = \sum_{m=1}^{\infty} \frac{C_m}{m^2\pi^2+a^2}
\biggl\{
A_1^{(m)}\cosh a\zeta + A_2^{(m)}\,\zeta\cosh a\zeta
+ A_1^{(m)}\,\zeta^2\cosh a\zeta + \\
+ B_1^{(m)}\,\zeta\sinh a\zeta
+ A_2^{(m)}\,\zeta\sinh a\zeta
+ B_2^{(m)}\,\zeta^2\sinh a\zeta + \\
+ (1+\alpha\zeta)\sin m\pi\zeta
+ \frac{4\alpha m\pi}{m^2\pi^2+a^2}\cos m\pi\zeta
\biggr\},
\end{multline}
where $A_1^{(m)}, \ldots, B_2^{(m)}$ are constants of integration to be
determined by the boundary conditions~\eqref{eq:7-218}.  These latter conditions lead to the equations:
\begin{align}
\label{eq:7-206}
A_1^{(m)} &= -\frac{4m\pi\alpha}{m^2\pi^2+a^2}, \nonumber\\[4pt]
B_1^{(m)} &= \frac{m\pi}{m^2\pi^2+a^2}
\{\alpha + \beta_m(\sinh a + a\cosh a) - \gamma_m\sinh a\}, \nonumber\\[4pt]
A_2^{(m)} &= -\frac{m\pi}{m^2\pi^2+a^2}
\{(\sinh^2 a + a)(\sinh a + a\cosh a) - \gamma_m \sinh a\}, \nonumber\\[4pt]
B_2^{(m)} &= \frac{m\pi}{m^2\pi^2+a^2}
\{(\sinh a\cosh a - a)(1 + \alpha)(-1)^{m+1} - \\
&\qquad -\gamma_m(\alpha\cosh a - \alpha\cosh a)\},
\end{align}
where
\begin{equation}
\label{eq:7-207}
\Delta = \sinh^2 a - a^2,
\end{equation}
\begin{equation}
\beta_m = \frac{4\alpha}{m^2\pi^2+a^2}\{(-1)^{m+1}\cosh a\},
\end{equation}
and
\begin{equation}
\gamma_m = (-1)^{m+1}(1+\alpha) + \frac{4\alpha}{m^2\pi^2+a^2}\, a\sinh a.
\end{equation}

Now substituting for $v$ and $u$ from equations~\eqref{eq:7-203} and~\eqref{eq:7-205} in equation~\eqref{eq:7-202}, we obtain
\begin{multline}
\label{eq:7-209}
\sum_{n=1}^{\infty} C_n(n^2\pi^2+a^2)\sin n\pi\zeta \\
= Ta^2 \sum_{m=1}^{\infty}\frac{C_m}{m^2\pi^2+a^2}
\biggl\{
A_1^{(m)}\cosh a\zeta + B_1^{(m)}\sinh a\zeta
+ A_2^{(m)}\,\zeta\cosh a\zeta + \\
+ B_2^{(m)}\,\zeta\sinh a\zeta
+ (1+\alpha\zeta)\sin m\pi\zeta
+ \frac{4\alpha m\pi}{m^2\pi^2+a^2}\cos m\pi\zeta
\biggr\}.
\end{multline}

Multiplying equation~\eqref{eq:7-209} by $\sin n\pi\zeta$ and integrating over the range of
$\zeta$, we obtain a system of linear homogeneous equations for the constants
$\mathscr{C}_m = C_m(m^2\pi^2+a^2)^{1/2}$; and the requirement that these constants are
not all zero leads to the secular equation
\begin{multline}
\label{eq:7-210}
\left|\frac{n\pi}{(n^2\pi^2+a^2)^{1/2}}
\{[1+(-1)^{n+1}\cosh a]\,A_1^{(m)}+\right. \\
+(-1)^{n+1}\left[\cosh a
- \frac{2a}{n^2\pi^2+a^2}\,\sinh a\right]\,A_2^{(m)} + \\
+\left[(-1)^{n+1}\sinh a
- \frac{2a}{n^2\pi^2+a^2}
\{1+(-1)^{n+1}\cosh a\}\right] B_2^{(m)} + \\
\left.+\alpha X_{mn}+\tfrac{1}{2}\delta_{mn}
- \tfrac{1}{2}(n^2\pi^2+a^2)\frac{\delta_{mn}}{a^2}\right| = 0,
\end{multline}
where
\begin{equation}
\label{eq:7-211}
X_{mn} = \begin{cases}
0 & \text{if } m+n \text{ is even and } m \neq n,\\
\frac{1}{2} & \text{if } m = n,\\
\dfrac{4mn}{n^2-m^2}\left(\dfrac{3}{m^2\pi^2+a^2}
- \dfrac{2}{n^2-m^2}\right.\cdot
\left.\dfrac{1}{\pi(n^2-m^2)}\right)
& \text{if } m+n \text{ is odd.}
\end{cases}
\end{equation}

On using the first two equations of~\eqref{eq:7-206}, equation~\eqref{eq:7-210} simplifies to
the form
\begin{multline}
\label{eq:7-212}
\left|\frac{n\pi}{n^2\pi^2+a^2}\frac{4mm\pi x}{m^2\pi^2+a^2}
\{(-1)^{m+n}-1\} - \right.\\
-\frac{2a}{n^2\pi^2+a^2}
\{(-1)^{n+1}\{A_2^{(m)}\sinh a + B_1^{(m)}\cosh a\} +
B_2^{(m)}\} + \\
\left.+\alpha X_{mn}+\tfrac{1}{2}\delta_{mn}
-\tfrac{1}{2}(n^2\pi^2+a^2)\frac{\delta_{mn}}{a^2}\right| = 0;
\end{multline}
and on substituting for the constants $A_2^{(m)}$ and $B_2^{(m)}$ their explicit solutions
given in~\eqref{eq:7-207}, we find that equation~\eqref{eq:7-212} simplifies greatly and we are
left with
\begin{multline}
\label{eq:7-213}
\left|\frac{4mm'n^2\alpha}{(n^2\pi^2+a^2)(m^2\pi^2+a^2)}
\{(-1)^{m+n}-1\} -\right.\\
- \frac{2amm'\pi^2\alpha}{(m^2\pi^2+a^2)^2(\sinh a)^2}
\{(\sinh a\cosh a - a)[1+(1+\alpha)(-1)^{m+n}] + \\
+(\sinh a - a\cosh a)[(-1)^{m+1} + (1+\alpha)(-1)^{n+1}]
\} - \\
- \frac{4a\alpha\sinh a}{(m^2\pi^2+a^2)^2}
\{[\sinh a + a(-1)^{n+1}][(-1)^{m+n}-1]\} + \\
\left.+\tfrac{1}{2}\delta_{mn}+\alpha X_{mn}
-\tfrac{1}{2}(n^2\pi^2+a^2)\frac{\delta_{mn}}{a^2T}\right| = 0.
\end{multline}

A first approximation to the solution of equation~\eqref{eq:7-213} is obtained by
setting the $(1,1)$-element of the matrix equal to zero.  We find
\begin{equation}
\label{eq:7-214}
\tfrac{1}{2}(\pi^2+a^2)\frac{1}{Ta^2} = \tfrac{1}{2}a+\tfrac{1}{2} -
\frac{2a\pi^2(2+\alpha)}{(\pi^2+a^2)^2}
[(\sinh a\cosh a - a) + (\sinh a - a\cosh a)].
\end{equation}
On further simplification, this gives
\begin{equation}
\label{eq:7-215}
T = \frac{2}{2+\alpha}\,\frac{(\pi^2+a^2)^3}{a^2[1 -
16a^2\pi^2\cosh^2\!\frac{1}{2}a/(\pi^2+a^2)^2(\sinh a + a)]}.
\end{equation}
We observe that, apart from the factor $2/(2+\alpha)$, this expression for $T$
is identical with what was found in Chapter~II (\S~17, equation~(311)) for
the Rayleigh number for the simple B\'enard problem in the first
approximation for the case of two rigid boundaries.

Consequently, in this approximation (cf.\ equation~\eqref{eq:7-199}),
\begin{equation}
\label{eq:7-216}
T_c = \frac{2}{2+\alpha}\times 1715 = \frac{3430}{1+\mu}
\quad \text{and} \quad
a_{\min} = 3{\cdot}12.
\end{equation}

We shall see below (\S\,(b)) that for $0 < \mu < 1$ equation~\eqref{eq:7-216} gives
values for the critical Taylor number which do not differ from those
obtained in the higher approximations by more than one per cent.  The
reason for this relatively high accuracy of the solution in the first
approximation will be made apparent in \S\,(d).

\subsection*{(b) \emph{Numerical results}}
A method of solving the infinite order characteristic equation which~\eqref{eq:7-213} provides for $T$ would be to set the determinant formed by the first
$n$ rows and columns of the secular matrix equal to zero and let $n$ take
increasingly larger values.  In practice, the usefulness of this method will
depend largely on how rapidly the lowest positive root of the resulting
equation of order $n$ tends to its limit as $n \to \infty$.  It appears that for the
problem on hand, the process converges quite rapidly.

In Table~XXXII the values of $T$ obtained with the aid of equation~\eqref{eq:7-213} in the different approximations are listed for those values of $a$
(for the assigned $\mu$'s) at which it was found (by trial and error) that $T$
attained its minimum value.  From an examination of this table, it
would appear that for $\mu \geqslant -1{\cdot}0$, the third approximation provides $T$
to well within one per cent of the true value.  For $-3{\cdot}0 \leqslant \mu \leqslant -1{\cdot}0$,
the calculations were carried out to as high approximations as seemed

\begin{table}[htbp]
\centering
\caption{The critical Taylor numbers and the associated wave numbers for different
values of $\mu$}
\label{tab:XXXII}
\begin{tabular}{ccccccccc}
\hline
& & First & Second & Third & Fourth & Fifth & Sixth \\
$\mu$ & $a$ & approxi- & approxi- & approxi- & approxi- & approxi- & approxi- \\
& & mation & mation & mation & mation & mation & mation \\
\hline
$1{\cdot}00$ & $3{\cdot}12$ & $1715{\cdot}1$ & $1725{\cdot}2$ & $1708$ & & & \\
$0{\cdot}50$ & $3{\cdot}12$ & $1286{\cdot}7$ & & & & & \\
$0{\cdot}25$ & $3{\cdot}12$ & $2744{\cdot}1$ & $2736{\cdot}4$ & & & & \\
$-0{\cdot}25$ & $3{\cdot}12$ & $3430{\cdot}9$ & $3430{\cdot}6$ & $3399{\cdot}3$ & & & \\
$-0{\cdot}50$ & $3{\cdot}25$ & $4573{\cdot}4$ & $4472{\cdot}5$ & & & & \\
$-0{\cdot}50$ & $3{\cdot}25$ & $6800{\cdot}2$ & $6430{\cdot}6$ & $6417{\cdot}6$ & & $6148$ \\
$-0{\cdot}75$ & $3{\cdot}34$ & & $9422{\cdot}7$ & $9434{\cdot}6$ & & & \\
$-1{\cdot}00$ & $3{\cdot}50$ & $11514$ & $11840$ & $11500$ & & & \\
$-1{\cdot}25$ & $3{\cdot}75$ & & & $16847$ & & & \\
$-1{\cdot}50$ & $3{\cdot}75$ & & $18015$ & $18735$ & & $18077$ & \\
$-1{\cdot}75$ & $4{\cdot}00$ & & & $30497$ & $30458$ & & \\
$-2{\cdot}00$ & $4{\cdot}00$ & & & $45000$ & $40992$ & & \\
$-2{\cdot}50$ & $4{\cdot}00$ & & & $67189$ & & & \\
$-3{\cdot}00$ & & & & & $195810$ & $177107$ & $177108$ \\
$-3{\cdot}00$ & & & & & $330748$ & $301116$ & $302496$ \\
$-3{\cdot}09$ & & & & & & & $304452$ \\
\hline
\end{tabular}
\end{table}

necessary.  It would appear that for $\mu < -3{\cdot}0$, one should go to approximations higher than the sixth to achieve comparable accuracy.  Fortunately, the considerations to be set out in \S\,(c) below make it unnecessary for us to obtain the solutions for $\mu < -3{\cdot}0$ by the present method.

\figurePlaceholder{69}{The critical Taylor number $T_c$ for the onset of instability as a function
of $\mu = \Omega_2/\Omega_1$.}

\figurePlaceholder{70}{The wavelength of the disturbance (in units of the gap width, $d$)
manifested at onset of instability as a function of $\mu = \Omega_2/\Omega_1$.}

In Table~XXXIII the coefficients $\mathscr{C}_m$ ($= C_m/(m^2\pi^2+a^2)^{1/2}$) for the marginal state in the expansion~\eqref{eq:7-205} for $u$ are given.  This table also includes
the critical Taylor numbers and the associated wave numbers.  The
$(T_c, \mu)$ and the $[a(T_c), \mu]$-relationships are further illustrated in Figs.~69
and 70.

\begin{table}[htbp]
\centering
\caption{The critical Taylor numbers and related constants for various values of $\mu$}
\label{tab:XXXIII}
\small
\begin{tabular}{cccccccc}
\hline
$\mu$ & $a$ & $T_c$ & $\mathscr{C}_1/\mathscr{C}_1$ & $\mathscr{C}_2/\mathscr{C}_1$ & $\mathscr{C}_3/\mathscr{C}_1$ & $\mathscr{C}_4/\mathscr{C}_1$ & $\mathscr{C}_5/\mathscr{C}_1$ \\
\hline
$1{\cdot}00$ & $3{\cdot}12$ & $1{\cdot}208\times 10^3$ & $1$ & $0{\cdot}001324$ & & $0{\cdot}001147$ & \\
$0{\cdot}50$ & $3{\cdot}12$ & $2{\cdot}273\times 10^3$ & & $0{\cdot}001146$ & & & \\
$0{\cdot}25$ & $3{\cdot}12$ & $2{\cdot}900\times 10^3$ & & $0{\cdot}001304$ & & & \\
$0$ & $3{\cdot}14$ & $3{\cdot}4 \times 10^3$ & & $0{\cdot}000537$ & & $0{\cdot}000144$ & \\
$-0{\cdot}25$ & $3{\cdot}14$ & $4{\cdot}617\times 10^3$ & & $0{\cdot}001156$ & & $0{\cdot}000177$ & \\
$-0{\cdot}50$ & $3{\cdot}20$ & $6{\cdot}417\times 10^3$ & & $0{\cdot}001856$ & & & \\
$-0{\cdot}60$ & $3{\cdot}30$ & $7{\cdot}688\times 10^3$ & & $0{\cdot}003159$ & & $0{\cdot}001194$ & \\
$-0{\cdot}75$ & $3{\cdot}30$ & $1{\cdot}082\times 10^4$ & & $0{\cdot}003195$ & & & \\
$-0{\cdot}90$ & $3{\cdot}40$ & $1{\cdot}494\times 10^4$ & & $0{\cdot}04742$ & & & \\
$-1{\cdot}00$ & $3{\cdot}50$ & $1{\cdot}196\times 10^4$ & & $0{\cdot}00463$ & & & \\
$-1{\cdot}00$ & $4{\cdot}00$ & $1{\cdot}868\times 10^4$ & & $0{\cdot}02139$ & & $0{\cdot}000929$ & $-0{\cdot}000439$ \\
$-1{\cdot}25$ & $4{\cdot}00$ & & & $0{\cdot}04916$ & & $+0{\cdot}000421$ & \\
$-1{\cdot}25$ & $4{\cdot}06$ & $4{\cdot}019\times 10^4$ & & $+0{\cdot}02381$ & & $+0{\cdot}000116$ & $-0{\cdot}00012$ \\
$-1{\cdot}50$ & $4{\cdot}06$ & & & & & & \\
$-1{\cdot}75$ & $4{\cdot}06$ & $6{\cdot}528\times 10^4$ & & $+0{\cdot}04009$ & & & \\
$-2{\cdot}00$ & $4{\cdot}50$ & $9{\cdot}558\times 10^4$ & & $+0{\cdot}04396$ & & $-0{\cdot}001808$ & \\
$-2{\cdot}50$ & $5{\cdot}00$ & $1{\cdot}721\times 10^5$ & & $+0{\cdot}05804$ & & $+0{\cdot}001777$ & $-0{\cdot}000927$ \\
$-3{\cdot}00$ & & $3{\cdot}015\times 10^5$ & & $+0{\cdot}07499$ & & $+0{\cdot}001045$ & \\
$-3{\cdot}09$ & $5{\cdot}14$ & & & & & & $+0{\cdot}001045$ \\
\hline
\end{tabular}
\end{table}

Fig.~71\,$a$ shows the profile of the velocity $u$ for $\mu = -3{\cdot}0$; and in Fig.~71\,$b$ the corresponding cell pattern is exhibited.  For $\mu = -3{\cdot}0$, the

\figurePlaceholder{71}{The velocity profile (a) and the cell pattern (b) for $\mu=-3$ and
$a=8{\cdot}14$.  In (a) the velocity is normalized to unit maximum amplitude; in
(b) the line separating the positive and the negative values of the stream
function $\psi$ is indicated by the dashed line (note that this does not occur at
the nodal surface which is at $\zeta=0{\cdot}25$); (c) is the corresponding cell pattern
for the lowest inviscid mode for the same value of $a$ and $\mu$.}

nodal surface occurs at $\zeta = 0{\cdot}25$; and it is of interest to observe how
rapidly the amplitude of $u$ decreases as we go beyond this point.  The
velocity profile illustrated in Fig.~71 should be contrasted with that
in Fig.~67 which corresponds to the lowest inviscid mode for the same
value of $a$ and $\mu$.

\subsection*{(c) \emph{An alternative method of solution}}
Equations~\eqref{eq:7-201}--\eqref{eq:7-202} present the first example of a characteristic
value problem that we have encountered which is not self-adjoint.  For
this reason, it may be of interest to describe an alternative method of
solution to determine the principles which might guide one in devising
methods of solution which may be as equally convergent as the one we
have followed in \S\,(b).

Eliminating $u$ between equations~\eqref{eq:7-201} and~\eqref{eq:7-202}, we obtain
together with the boundary conditions
\begin{equation}
\label{eq:7-217}
(D^2-a^2)^3 v = -Ta^2(1+\alpha\zeta)v
\end{equation}
\begin{equation}
v = (D^2-a^2)v = D(D^2-a^2)v = 0
\quad \text{for } \zeta = 0 \text{ and } 1.
\label{eq:7-218}
\end{equation}
Rewriting equation~\eqref{eq:7-217} in the manner
\begin{equation}
\label{eq:7-219}
(D^2-a^2)^3 v = (1+\alpha\zeta)\psi
\end{equation}
and
\begin{equation}
\label{eq:7-220}
\psi = -Ta^2 v,
\end{equation}
we expand $\psi$ and $v$ in the forms
\begin{equation}
\label{eq:7-221}
\psi = \sum_{m=1}^{\infty} C_m \sin m\pi\zeta
\quad \text{and} \quad
v = \sum_{m=1}^{\infty} C_m\, v_m,
\end{equation}
where $v_m$ is the solution of the equation
\begin{equation}
\label{eq:7-222}
(D^2-a^2)^3 v_m = (1+\alpha\zeta)\sin m\pi\zeta
\end{equation}
which satisfies the boundary conditions~\eqref{eq:7-218}.  We then insert the
expansions~\eqref{eq:7-221} in~\eqref{eq:7-220} to obtain the secular equation.  The general solution of equation~\eqref{eq:7-222} is readily seen to be
\begin{multline}
\label{eq:7-223}
v_m = -\frac{1}{(m^2\pi^2+a^2)^2}
\biggl\{
A_1^{(m)}\cosh a\zeta + A_2^{(m)}\,\zeta\cosh a\zeta
+ A_1^{(m)}\,\zeta^2\cosh a\zeta + \\
+ B_1^{(m)}\sinh a\zeta
+ B_2^{(m)}\,\zeta\sinh a\zeta
+ B_2^{(m)}\,\zeta^2\sinh a\zeta + \\
+(1+\alpha\zeta)\sin m\pi\zeta
+ \frac{6\alpha m\pi}{m^2\pi^2+a^2}\cos m\pi\zeta
\biggr\},
\end{multline}
where $A_1^{(m)},\ldots, B_2^{(m)}$ are constants of integration to be determined by the
boundary conditions~\eqref{eq:7-218}.  These latter conditions lead to the equations:
\begin{align}
\label{eq:7-224}
A_1^{(m)} &= -\frac{6m\pi\alpha}{m^2\pi^2+a^2}, &
aA_1^{(m)}+3\,B_2^{(m)} &= \frac{m\pi}{2a}(m^2\pi^2+a^2),\nonumber\\
A_2^{(m)} + a\,B_2^{(m)} &= 2m\pi\alpha, & & \nonumber\\[3pt]
A_1^{(m)}\sinh a + A_2^{(m)}(&2a\sinh a + \cosh a) + \nonumber\\
+ B_1^{(m)}&a\cosh a + B_2^{(m)}(2a\cosh a + 3\sinh a)\,a\nonumber\\
&= (-1)^{m+1}\frac{6m\pi\alpha}{m^2\pi^2+a^2},\\[3pt]
A_1^{(m)}a\cosh a + A_2^{(m)}&(2a\cosh a + 3\sinh a)\,a + \nonumber\\
+ B_1^{(m)}&a\sinh a + B_2^{(m)}(2a\sinh a + 3\cosh a)\,a \nonumber\\
&= \frac{m\pi}{2a}(m^2\pi^2+a^2)(1+\alpha)(-1)^m.
\end{align}
On solving these equations, we find that
\begin{align}
\label{eq:7-207a}
A_1^{(m)} &= -\frac{4m\pi\alpha}{m^2\pi^2+a^2}, \nonumber\\
B_1^{(m)} &= \frac{m\pi}{m^2\pi^2+a^2}\,
\{\alpha+\beta_m(\sinh a + a\cosh a)-\gamma_m\sinh a\}, \nonumber\\
A_2^{(m)} &= -\frac{m\pi}{m^2\pi^2+a^2}
\{(\sinh^2 a + a)(\sinh a + a\cosh a)-\gamma_m \sinh a\}, \nonumber\\
B_2^{(m)} &= \frac{m\pi}{m^2\pi^2+a^2}
\{(\sinh a\cosh a - a)(1+\alpha)(-1)^{m+1} - \gamma_m(\alpha\cosh a - \alpha\cosh a)\},
\end{align}
where
$\Delta = \sinh^2 a - a^2$,
$\beta_m = \frac{4\alpha}{m^2\pi^2+a^2}\{(-1)^{m+1}\cosh a\}$,
and
$\gamma_m = (-1)^{m+1}(1+\alpha) + \frac{4\alpha}{m^2\pi^2+a^2}\, a\sinh a$.

Now substituting for $\psi$ and $v$ their expansions in equation~\eqref{eq:7-220}, we
obtain
\begin{equation}
\label{eq:7-225}
\frac{1}{Ta^2}\sum_{m=1}^{\infty} C_m \sin m\pi\zeta
= -\sum_{m=1}^{\infty} C_m\, v_m;
\end{equation}
and this clearly leads to the secular equation
\begin{equation}
\label{eq:7-226}
\left|\frac{1}{Ta^2}\,\delta_{mn}
+2\langle n|m\rangle\right| = 0,
\end{equation}
where $\langle n|m\rangle$ stands for the matrix element
\begin{equation}
\label{eq:7-227}
\langle n|m\rangle = \int_0^1 v_m \sin n\pi\zeta\,d\zeta.
\end{equation}
(Observe that the matrix $\langle k|x|j\rangle$ is not Hermitian.)

On evaluating the matrix element $\langle n|m\rangle$, we find after some lengthy but
elementary reductions that the secular equation takes the explicit form
\begin{multline}
\label{eq:7-228}
\left|
\frac{nmm'\pi^2}{n^2\pi^2+a^2}\left(\frac{6\alpha}{m^2\pi^2+a^2}
+ \frac{4\alpha}{n^2\pi^2+a^2}\right)
\{(-1)^{m+n}-1\} + \right. \\
+ \frac{8a^2}{(m^2\pi^2+a^2)^2}\left\{
\frac{4\alpha^2}{4a^2}\right\}
\{1+(-1)^{m+n}+(-1)^m\}
\cosh a\{(-1)^{m+n}+(-1)^n\} + \\
+\alpha(-1)^m\,
\frac{1+(-1)^{m+n}\cosh a}{\sinh a + (-1)^{n+1}a}
+\alpha\left\{(-1)^{m+n+1}\frac{\sinh a}{\sinh a + (-1)^{n+1}a}\right\} + \\
\left.+\tfrac{1}{2}\delta_{mn}+\alpha X_{mn}
-\tfrac{1}{2}(n^2\pi^2+a^2)^2\frac{\delta_{mn}}{a^2 T}\right| = 0,
\end{multline}
where
\begin{equation}
\label{eq:7-229}
X_{mn} = \begin{cases}
0 & \text{if } m+n \text{ is even and } m \neq n,\\
\frac{1}{2} & \text{if } m = n,\\
\dfrac{4mn}{n^2-m^2}\left(\dfrac{3}{m^2\pi^2+a^2}
- \dfrac{1}{\pi(n^2-m^2)}\right) & \text{if } m+n \text{ is odd.}
\end{cases}
\end{equation}

One might have thought that, since equation~\eqref{eq:7-228} is based on the
solution of an equation of order six, the characteristic values derived
from the present equation in the different approximations would show
an even more rapid convergence than has been found by the method of
\S\,(a).  However, this is not the case: the results given by equations~\eqref{eq:7-213}
and~\eqref{eq:7-228} are practically indistinguishable.  Thus:
\begin{equation}
\label{eq:7-230}
\left.
\begin{aligned}
T &= 6417{\cdot}8 &&\text{(from equation~\eqref{eq:7-228} in 3rd approximation)}\\
  &= 6417{\cdot}1 &&\text{(from equation~\eqref{eq:7-213} in 3rd approximation)}
\end{aligned}
\right\}\\
\text{for } \mu = -0{\cdot}5 \text{ and } a = 3{\cdot}20
\end{equation}
and
\begin{equation*}
T\left\{
\begin{aligned}
&= 95624 &&\text{(from equation~\eqref{eq:7-228} in 4th approximation)}\\
&= 95625 &&\text{(from equation~\eqref{eq:7-213} in 4th approximation)}
\end{aligned}
\right\}\\
\end{equation*}
for $\mu = -2{\cdot}0$ and $a = 6{\cdot}05$.

The explanation for this near efficacy of the two methods probably
derives from the fact that the advantage gained in the present method
by solving an equation of order six has been lost in the artificial treatment
of the equation and the boundary conditions.  In the method described
in \S\,(a), the variables used had direct physical meanings, and the
boundary conditions relevant to each and the relations between them
were satisfied identically; and this is perhaps the real key to the problem.

\subsection*{(d) \emph{An approximate solution for $\mu \to 1$}}
From the results given in Table~XXXIII, it appears that the formula
\begin{equation}
\label{eq:7-231}
T_c = \frac{3416}{1+\mu}
\end{equation}
gives a very good approximation to the true values for $0 \leqslant \mu \leqslant 1$;
and, moreover, the wave number ($\sim 3{\cdot}12$) at which instability sets in,
hardly seems to depend on $\mu$ in the same range.  Formula~\eqref{eq:7-231} was
derived by Taylor by a method of solution very different from the ones
we have described; and we have also seen how effectively the same
formula is given by the \emph{first} approximation to the solution of the secular
equation~\eqref{eq:7-213}.

We shall now show how one can derive~\eqref{eq:7-231} by a simple perturbation
procedure applied to the solution~\eqref{eq:7-217}; at the same time, we
shall obtain a correction term to~\eqref{eq:7-231} which extends its range of validity.

It is convenient for our present purposes to translate the origin of the
coordinate system to be midway between the two cylinders.  Thus, let
\begin{equation}
\label{eq:7-232}
\zeta' = x + \tfrac{1}{2}.
\end{equation}
so that the limits of $x$ are $\pm\frac{1}{2}$.  In terms of $x$, equation~\eqref{eq:7-217} becomes
\begin{equation}
\label{eq:7-233}
(D^2-a^2)^3 v = -\lambda(1+\epsilon x)v,
\end{equation}
where
\begin{equation}
\label{eq:7-234}
\lambda = \tfrac{1}{2}Ta^2(1+\mu) \quad \text{and} \quad
\epsilon = -2\,\frac{1-\mu}{1+\mu},
\end{equation}
and the boundary conditions are
\begin{equation}
\label{eq:7-235}
u = (D^2-a^2)v = D(D^2-a^2)v = 0
\quad \text{for } x = \pm\tfrac{1}{2}.
\end{equation}
When $\epsilon = 0$, the characteristic value problem presented by equations~\eqref{eq:7-233} and~\eqref{eq:7-235} reduces to the one for the simple B\'enard problem (for
the case of two rigid boundaries) when formulated in terms of the
amplitude $\Theta$ of the fluctuations in the temperature (see Chapter~II,
\S\S~12 and 13\,(b)).  Let $\lambda_j = R_j a^2$ in the notation of Chapter~II for
$j = 0, 1, \ldots$, denote the different characteristic values of the equation
\begin{equation}
\label{eq:7-236}
(D^2-a^2)^3\Theta = -\lambda\Theta,
\end{equation}
for the boundary conditions
\begin{equation}
\label{eq:7-237}
\Theta = (D^2-a^2)\Theta = D(D^2-a^2)\Theta = 0
\quad \text{for } x = \pm\tfrac{1}{2}.
\end{equation}
and let the functions belonging to $\lambda_j$ be distinguished by a subscript $j$.
We know from Chapter~II, \S~13\,(b) that the functions\footnote{As defined here $\Theta$ and $W$ are the same functions which describe the perturbations
in the temperature and in the normal component of the velocity for the B\'enard problem
(see equations II~(90)).}
\begin{equation}
\label{eq:7-238}
\Theta_j \quad \text{and} \quad W_j = (D^2-a^2)\Theta_j
\end{equation}
are mutually orthogonal\footnote{Equation~II~(161) directly implies that
\[
\int \{(D\Theta_j)(D\Theta_k)+a^2\Theta_j\Theta_k\}\,dx = 0 \quad \text{if } j \neq k.
\]
Integrating by parts the first term in the integrand, we obtain
\[
\int \Theta_j(D^2-a^2)\Theta_k\,dx = \int \Theta_j\,W_k\,dx = 0 \quad (j \neq k).
\]}
when $j \neq k$; without loss of generality, we
may suppose that they are so normalized that
\begin{equation}
\label{eq:7-239}
\int_{-1/2}^{+1/2}\Theta_j\, W_k\, dx = \delta_{jk}.
\end{equation}

Turning to the solution of~\eqref{eq:7-233}, we first observe that according to
the results given in Table~III (Chapter~II) the lowest value of $\lambda/a^2$,
when $\epsilon = 0$, is 1708 and occurs for $a = 3{\cdot}117$.  By~\eqref{eq:7-234} the corresponding value of $T$ is $3416/(1+\mu)$; and this is exactly the same as~\eqref{eq:7-231}.
Thus, formula~\eqref{eq:7-231} appears as the zero-order term in a perturbation
series.  To obtain the higher-order terms, we expand the various quantities in powers of $\epsilon$.  Thus, we write
\begin{equation}
\label{eq:7-240}
\left.\begin{aligned}
v &= \Theta_0+\epsilon v^{(1)}+\epsilon^2 v^{(2)}+\cdots\\
\lambda &= \lambda_0+\epsilon\lambda^{(1)}+\epsilon^2\lambda^{(2)}+\cdots
\end{aligned}\right\}.
\end{equation}
Substituting these expansions in equation~\eqref{eq:7-233} and equating the
different powers of $\epsilon$, we obtain:
\begin{align}
\label{eq:7-241}
(D^2-a^2)^3 \Theta_0 &= -\lambda_0\Theta_0, \\
\label{eq:7-242}
(D^2-a^2)^3 v^{(1)} &= -(\lambda_0 v^{(1)}+\lambda_0 x\Theta_0+\lambda^{(1)}\Theta_0), \\
\label{eq:7-243}
(D^2-a^2)^3 v^{(2)} &= -(\lambda_0 v^{(2)}+\lambda_0 xv^{(1)}+\lambda^{(1)} v^{(1)}
+\lambda^{(2)}\Theta_0+\lambda^{(1)}x\Theta_0),
\end{align}
etc.

Equation~\eqref{eq:7-241} is clearly satisfied.  To solve equation~\eqref{eq:7-242}, we assume
for $v^{(1)}$ a solution of the form
\begin{equation}
\label{eq:7-244}
v^{(1)} = \sum_{j=0}^{\infty} A_j^{(1)}\Theta_j.
\end{equation}
Inserting this solution in~\eqref{eq:7-242}, we obtain
\begin{equation}
\label{eq:7-245}
\sum_{j=0}^{\infty} A_j^{(1)}\lambda_j\,\Theta_j
= \lambda_0 \sum_{j=0}^{\infty} A_j^{(1)}\Theta_j
+ \lambda_0 x\Theta_0 + \lambda^{(1)}\Theta_0.
\end{equation}
Multiply this equation by $W_k$ and integrate over the range of $x$.  Making
use of the orthogonality relation~\eqref{eq:7-239}, we obtain
\begin{equation}
\label{eq:7-246}
A_k^{(1)}\lambda_k = \lambda_0\, A_k^{(1)}
+ \lambda_0\langle k|x|0\rangle + \lambda^{(1)}\delta_{k0},
\end{equation}
where the notation
\begin{equation}
\label{eq:7-247}
\langle k|x|j\rangle = \int_{-1/2}^{+1/2} W_k\, x\,\Theta_j\, dx
\end{equation}
has been adopted.  (Observe that the matrix $\langle k|x|j\rangle$ is not Hermitian.)
For $k = 0$, equation~\eqref{eq:7-246} gives
\begin{equation}
\label{eq:7-248}
\lambda^{(1)} = -\lambda_0\langle 0|x|0\rangle = -\lambda_0
\int_{-1/2}^{+1/2} W_0\, x\, \Theta_0\, dx.
\end{equation}
Since the proper functions $W_0$ and $\Theta_0$ belonging to $\lambda_0$ are even, the matrix
element $\langle 0|x|0\rangle$ is zero; and we conclude
\begin{equation}
\label{eq:7-249}
\lambda^{(1)} = 0.
\end{equation}
Therefore, formula~\eqref{eq:7-231} is in error by a term which is of the \emph{second
order} in $\epsilon = -2(1-\mu)/(1+\mu)$; and this explains its success in representing the results of more exact calculations as well as it does.

For $k \neq 0$, equation~\eqref{eq:7-246} gives
\begin{equation}
\label{eq:7-250}
A_k^{(1)} = \frac{\lambda_0}{\lambda_k-\lambda_0}\,\langle k|x|0\rangle.
\end{equation}
The coefficient $A_0^{(1)}$ is left unspecified; but as we shall presently see, this
does not lead to any ambiguities in the solution for the physical problem.

Consider next equation~\eqref{eq:7-243}.  Since $\lambda^{(1)} = 0$, this equation reduces to
\begin{equation}
\label{eq:7-251}
(D^2-a^2)^3 v^{(2)} = -(\lambda_0 v^{(2)}
+ \lambda_0 xv^{(1)} + \lambda^{(1)}xv^{(1)} + \lambda^{(2)}\Theta_0).
\end{equation}
In solving this equation, we again expand $v^{(2)}$ in terms of $\Theta_j$.  Thus, with
the assumption
\begin{equation}
\label{eq:7-252}
v^{(2)} = \sum_{j=0}^{\infty} A_j^{(2)}\Theta_j,
\end{equation}
equation~\eqref{eq:7-251} gives
\begin{equation}
\label{eq:7-253}
\sum_{j=0}^{\infty} A_j^{(2)}\lambda_j\,\Theta_j
= \lambda_0\sum_{j=0}^{\infty} A_j^{(2)}\Theta_j
+ \lambda_0 xv^{(1)} + \lambda^{(2)}\Theta_0.
\end{equation}
Multiplying this equation by $W_k$ and integrating over the range of $x$, we
obtain
\begin{equation}
\label{eq:7-254}
A_k^{(2)}\lambda_k = \lambda_0\, A_k^{(2)}
+ \lambda_0 \int W_k\, xv^{(1)}\, dx + \lambda^{(2)}\delta_{k0}.
\end{equation}
Hence
\begin{equation}
\label{eq:7-255}
\lambda^{(2)} = -\lambda_0 \int_{-1/2}^{+1/2} W_0\, xv^{(1)}\, dx.
\end{equation}
Substituting for $v^{(1)}$ from equation~\eqref{eq:7-244}, we obtain
\begin{equation}
\label{eq:7-256}
\lambda^{(2)} = -\lambda_0 \sum_{j=0}^{\infty} A_j^{(1)}
\int_{-1/2}^{+1/2} W_0\, x\,\Theta_j\, dx.
\end{equation}
With the notation~\eqref{eq:7-247} for the matrix elements of $x$, we can write
\begin{equation}
\label{eq:7-257}
\lambda^{(2)} = -\lambda_0 \sum_{j=0}^{\infty} A_j^{(1)}\langle 0|x|j\rangle.
\end{equation}
Since $\langle 0|x|0\rangle = 0$, the term $j = 0$ in~\eqref{eq:7-257} makes no contribution; and
substituting for $A_j^{(1)}$ ($j \neq 0$) from equation~\eqref{eq:7-250}, we have
\begin{equation}
\label{eq:7-258}
\lambda^{(2)} = -\lambda_0^2\sum_{j \neq 0}
\frac{\langle 0|x|j\rangle\langle j|x|0\rangle}{\lambda_j-\lambda_0}.
\end{equation}
This expression for the second-order change in the characteristic
value is very similar to the expression one has in the quantum theory for the
second-order change in the energy levels of an atomic system caused by
a perturbation.  An important difference, however, is the non-Hermitian
character of the matrix representing the perturbation.

In the particular problem to which equation~\eqref{eq:7-258} refers, the `energy
levels' are so widely spaced (e.g.\ $\lambda_1 \sim 15\lambda_0$) that it will be sufficient to
retain only the first term in the infinite sum representing $\lambda^{(2)}$.  Thus, we
may write
\begin{equation}
\label{eq:7-259}
\lambda^{(2)} \simeq -\frac{\lambda_0^2}{\lambda_1-\lambda_0}
\langle 0|x|1\rangle\langle 1|x|0\rangle,
\end{equation}
To the second order in $\epsilon$, the desired expression for the characteristic
value is
\begin{equation}
\label{eq:7-260}
\lambda = \lambda_0\!\left\{1
- \epsilon^2\frac{\lambda_0}{\lambda_1-\lambda_0}
\langle 0|x|1\rangle\langle 1|x|0\rangle\right\}.
\end{equation}
Substituting for $\lambda$ from equation~\eqref{eq:7-234} and remembering that $\lambda_j = R_j\, a^2$,
we can rewrite equation~\eqref{eq:7-260} in the form
\begin{equation}
\label{eq:7-261}
T = \frac{2R_0}{1+\mu}\!\left\{1
- \epsilon^2\frac{R_0}{R_1-R_0}\langle 0|x|1\rangle\langle 1|x|0\rangle\right\},
\end{equation}
where $R_0$ and $R_1$ are the Rayleigh numbers for the lowest even and odd
modes for the assigned $a^2$.

To obtain the critical Taylor number for the onset of instability, we
must minimize the quantity on the right-hand side of equation~\eqref{eq:7-261}
as a function of $a$.  The value of $a$ at which $R_0$ ($= 1708$, by $a_{\min}^{(0)} = 3{\cdot}117$) at which $R_0$ attains its
minimum value, $R_c$ ($= 1708$), by an amount of order $\epsilon^2$.  At the displaced position of the new minimum, $R_0$ will differ from $R_c$ only by
an amount of order $\epsilon^4$.  Therefore, to order $\epsilon^2$, we may write
\begin{equation}
\label{eq:7-262}
T_c = \frac{2R_c}{1+\mu}\!\left\{1
- 4\!\left(\frac{1-\mu}{1+\mu}\right)^{\!2}
\frac{R_0}{R_1^*-R_c}\,
\langle 0|x|1\rangle\langle 1|x|0\rangle\right\},
\end{equation}
where $R_1^*$ ($= 2{\cdot}4982\times 10^4$) is the Rayleigh number for the first odd
mode for $a_{\min}^{(0)}$ (see Table~I, p.~38).  [In~\eqref{eq:7-262} we have substituted for $\epsilon$ its value
$-2(1-\mu)/(1+\mu)$.]

The solution for the first odd mode which is needed for the evaluation
of the matrix elements in~\eqref{eq:7-262} can be found by the method described in
Chapter~II, \S~15\,(b).  And it is found that the numerical form of equation~\eqref{eq:7-262} is
\begin{equation}
\label{eq:7-263}
T_c = \frac{3416}{1+\mu}
\!\left\{1-7{\cdot}61\times 10^{-2}
\!\left(\frac{1-\mu}{1+\mu}\right)^{\!2}\right\}.
\end{equation}
This formula gives $T_c = 3{\cdot}390\times 10^3$ and $6{\cdot}36\times 10^3$ for $\mu = 0$ and $-0{\cdot}5$,
respectively; these values should be compared with $3{\cdot}390\times 10^3$ and
$6{\cdot}42\times 10^3$ given by the `exact' calculations.

\subsection*{(e) \emph{The asymptotic behaviour for $(1-\mu) \to \infty$}}
From an examination of the results given in Table~XXXIII, it appears
that for $(1-\mu) \to \infty$ the following asymptotic relations obtain:
\begin{equation}
\label{eq:7-264}
T_c \to \tau(1-\mu)^4 \quad \text{and} \quad
a(T_c) \to q(1-\mu) \quad \text{as } (1-\mu) \to \infty,
\end{equation}
where $\tau$ and $q$ are certain constants.  Thus, from the tabulated values
we find:
\begin{equation}
\label{eq:7-265}
\frac{T_c}{(1-\mu)^4}\left\{
\begin{aligned}
&= 1180{\cdot}5\\
&= 1180{\cdot}2\\
&= 1181{\cdot}6
\end{aligned}
\right.
\quad \text{and} \quad
\frac{a(T_c)}{1-\mu}\left\{
\begin{aligned}
&= 2{\cdot}03\\
&= 2{\cdot}03\\
&= 2{\cdot}035
\end{aligned}
\right.
\quad \text{for } 1-\mu\left\{
\begin{aligned}
&= 3{\cdot}0\\
&= 3{\cdot}5\\
&= 4{\cdot}0.
\end{aligned}
\right.
\end{equation}
The fact that asymptotic relations such as~\eqref{eq:7-264} should obtain can be
readily understood.  For, it will be recalled that the origin of the instability, in the case when the two cylinders are in counter rotation, lies
in the violation of the Rayleigh criterion in the layers adjoining the inner
cylinder and extending not much beyond the nodal surface.  On these
grounds, one may expect that as $(1-\mu) \to \infty$ the critical Taylor
number for the onset of instability must be comparable to the critical Taylor
number for the case $\mu = 0$, and for a gap width equal to the distance
$d_0$ of the nodal surface to the inner cylinder.  For the distribution of the
angular velocity given by~\eqref{eq:7-191},
\begin{equation}
\label{eq:7-266}
d_0 = d/(1-\mu).
\end{equation}
Since the gap width occurs with the fourth power in the definition of $T$,
the foregoing arguments will lead one to expect an asymptotic behaviour
of the type
\begin{equation}
\label{eq:7-267}
T_c \to \tau(1-\mu)^4,
\end{equation}
where $\tau$ is comparable to the Taylor number for $\mu = 0$.  For the same
reasons, one may also expect that the wave number of the associated
disturbance will show the behaviour
\begin{equation}
\label{eq:7-268}
a(T_c) \to q(1-\mu),
\end{equation}
where $q$ is comparable to the wave number appropriate for $\mu = 0$.

The fact that the constants $\tau$ ($\sim 1182$) and $q$ ($\sim 2{\cdot}035$) indicated by
the calculations are substantially different from the values 3416 and
$3{\cdot}12$ appropriate for $\mu = 0$ should not cause much surprise; for, at
$\zeta = d_0/d$, the conditions which prevail are far from those at $\zeta = 1$ for
$\mu = 0$; and it is also not true, even asymptotically, that the disturbance
does not extend beyond the nodal surface.

While the foregoing arguments give physical grounds for expecting
asymptotic relations of the forms~\eqref{eq:7-267} and~\eqref{eq:7-268}, we can equally trace
their origin to the analytical structure of the underlying characteristic
value problem.  Thus, with the substitution
\begin{equation}
\label{eq:7-269}
x = a\zeta,
\end{equation}
equation~\eqref{eq:7-217} becomes
\begin{equation}
\label{eq:7-270}
(D^2-1)^3 v = -\frac{T}{a^4}\!\left(1-\frac{1-\mu}{a}\,x\right)v;
\end{equation}
and the corresponding boundary conditions are
\begin{equation}
\label{eq:7-271}
v = (D^2-1)v = D(D^2-1)v = 0
\quad \text{for } x = 0 \text{ and } a.
\end{equation}
Suppose now that
\begin{equation}
\label{eq:7-272}
\frac{T}{a^4} \to \lambda \quad \text{and} \quad
\frac{1-\mu}{a} \to \gamma \quad
\text{as } (1-\mu) \to \infty \text{ and } a \to \infty.
\end{equation}
In this limit equations~\eqref{eq:7-270} and~\eqref{eq:7-271} become
\begin{equation}
\label{eq:7-273}
(D^2-1)^3 v = -\lambda(1-\gamma x)v
\end{equation}
and
\begin{equation}
\label{eq:7-274}
v = (D^2-1)v = D(D^2-1)v = 0
\quad \text{for } x = 0 \text{ and } \infty.
\end{equation}
For an assigned value of $\gamma$, equations~\eqref{eq:7-273} and~\eqref{eq:7-274} will determine
a sequence of possible values of $\lambda$; and to determine the constants $\tau$
and $q$ in the asymptotic relations~\eqref{eq:7-267} and~\eqref{eq:7-268}, we must find the
minimum of $\lambda/\gamma^4$ as a function of $\gamma$.  Thus,
\begin{equation}
\label{eq:7-275}
\tau = \min_\gamma\!\left\{\frac{\lambda}{\gamma^4}\right\},
\end{equation}
and $q$ is the reciprocal of the value of $\gamma$ at which $\lambda/\gamma^4$ attains its minimum
value.

The characteristic value problem presented by equations~\eqref{eq:7-273} and~\eqref{eq:7-274} has not been solved satisfactorily at the present time: for various
reasons, the methods which have proved successful in other cases seem
to fail when applied to the present problem.

% ========================================================================
% [MERGED] Content below was originally in chapter_7_part2.tex
% ========================================================================

%% §72 %%
\section{On the principle of the exchange of stabilities}\label{sec:7-72}

As we have remarked in \S\,71 and as we shall see in detail in \S\,74,
experiments confirm that the instability of the Couette flow sets in as a
secondary stationary flow.  Without attempting a full theoretical
justification of the principle of the exchange of stabilities, we shall give
some reasons why we may expect its validity in this instance.  For this
purpose, we return to equations~\eqref{eq:7-196} and~\eqref{eq:7-197} in which the time
constant~$\sigma$ is still retained.  Translating the origin of the coordinate
system to be midway between the two cylinders and replacing~$\mathfrak{u}$ by
$\mathfrak{u}(1+\mu)/2$, we obtain (cf.\ equations~\eqref{eq:7-201}--\eqref{eq:7-204})
\begin{equation}
\label{eq:7-276}
(D^2-a^2)(D^2_*-a^2-\sigma)\mathfrak{u} = (1+\varepsilon x)v
\end{equation}
and
\begin{equation}
\label{eq:7-277}
(D^2-a^2-\sigma)v = -\mathscr{T}a^2\mathfrak{u},
\end{equation}
where
\begin{equation}
\label{eq:7-278}
\mathscr{T} = \tfrac{1}{4}(1+\mu)T \quad \text{and} \quad
\varepsilon = -a\,\frac{1-\mu}{1+\mu}.
\end{equation}
And the boundary conditions are the same as hitherto.

Consider first the case $\mu > 0$.  As we have seen in \S\,71\,(d), in this case,
when $\sigma = 0$, the lowest characteristic values of the problem hardly
depend on the term linear in~$x$ on the right-hand side of equation~\eqref{eq:7-276}.
In fact, to order~$\varepsilon$, the characteristic values are the same as when $\varepsilon = 0$;
and, moreover, on account of the wide spacing of the characteristic
values of the `unperturbed problem', the coefficient of~$\varepsilon^2$ in a perturbation
series for~$\mathscr{T}$ is very small (see equation~\eqref{eq:7-263}).  It would, therefore,
appear that even when $\sigma \ne 0$, we can to the \emph{first} order in~$\varepsilon$ ignore the term
in~$\varepsilon x$ on the right-hand side of equation~\eqref{eq:7-276}.  We shall, then, be left with
\begin{equation}
\label{eq:7-279}
(D^2-a^2)(D^2-a^2-\sigma)\mathfrak{u} = v,
\end{equation}
and
\begin{equation}
\label{eq:7-280}
(D^2-a^2-\sigma)v = -\mathscr{T}a^2\mathfrak{u},
\end{equation}
together with the boundary conditions
\begin{equation}
\label{eq:7-281}
v = \mathfrak{u} = D\mathfrak{u} = 0 \quad \text{for } x = \pm\tfrac{1}{2}.
\end{equation}

Now multiply equation~\eqref{eq:7-279} by~$\mathfrak{u}^*$ and integrate over the range of~$x$.
We obtain by an integration by parts
\begin{equation}
\label{eq:7-282}
\int_{-\frac{1}{2}}^{+\frac{1}{2}} v\mathfrak{u}^*\,dx
= \int_{-\frac{1}{2}}^{+\frac{1}{2}}
  \bigl\{\mathfrak{u}^*\!\bigl[(D^2-a^2)^2 - \sigma(D^2-a^2)\bigr]\mathfrak{u}\bigr\}\,dx
= \int_{-\frac{1}{2}}^{+\frac{1}{2}}
  \bigl\{|(D\mathfrak{u}|^2+a^2|\mathfrak{u}|^2\bigr)^2\bigr\}\,dx
+ \sigma\!\int_{-\frac{1}{2}}^{+\frac{1}{2}}
  \bigl(|D\mathfrak{u}|^2+a^2|\mathfrak{u}|^2\bigr)\,dx.
\end{equation}
On the other hand, from equation~\eqref{eq:7-280} we have
\begin{equation}
\label{eq:7-283}
-\mathscr{T}a^2\!\int_{-\frac{1}{2}}^{+\frac{1}{2}} v\mathfrak{u}^*\,dx
= \int_{-\frac{1}{2}}^{+\frac{1}{2}}
  v(D^2-a^2-\sigma)v^*\,v^*\,dx
= -\int_{-\frac{1}{2}}^{+\frac{1}{2}}
  \bigl\{|Dv|^2+(a^2+\sigma^*)|\,v|^2\bigr\}\,dx.
\end{equation}
Combining equations~\eqref{eq:7-282} and~\eqref{eq:7-283}, we obtain
\begin{multline}
\label{eq:7-284}
\mathscr{T}a^2\!\int_{-\frac{1}{2}}^{+\frac{1}{2}}
  \bigl\{|(D^2-a^2)\mathfrak{u}|^2\,dx
  + \sigma\!\int_{-\frac{1}{2}}^{+\frac{1}{2}}
    (|D\mathfrak{u}|^2+a^2|\mathfrak{u}|^2)\,dx\bigr\} \\
= \int_{-\frac{1}{2}}^{+\frac{1}{2}}
  \bigl(|Dv|^2+a^2|v|^2\bigr)\,dx
  + \sigma^*\!\int_{-\frac{1}{2}}^{+\frac{1}{2}} |v|^2\,dx.
\end{multline}
On equating the imaginary parts of equation~\eqref{eq:7-284}, we obtain
\begin{equation}
\label{eq:7-285}
\operatorname{im}(\sigma)\!\Biggl\{
\mathscr{T}a^2\!\int_{-\frac{1}{2}}^{+\frac{1}{2}}
  \bigl(|D\mathfrak{u}|^2+a^2|\mathfrak{u}|^2\bigr)\,dx
  + \int_{-\frac{1}{2}}^{+\frac{1}{2}} |v|^2\,dx\Biggr\}
= 0.
\end{equation}
From this equation it follows that
\begin{equation}
\label{eq:7-286}
\operatorname{im}(\sigma) = 0 \quad \text{when } \mathscr{T} > 0;
\end{equation}
and we conclude that the principle of the exchange of stabilities is valid.
In view of the fact that for $\mu > 0$ the first-order formula~\eqref{eq:7-231} gives
values which differ from the results of more exact calculations by less
than one per cent, it would appear that the foregoing arguments effectively exclude the occurrence of overstability for $\mu > 0$.

As~$\mu$ turns negative, the replacement of equations~\eqref{eq:7-276} and~\eqref{eq:7-277} by
\eqref{eq:7-279} and~\eqref{eq:7-280} will lead to rapidly increasing errors, and the possibility
of overstability occurring cannot be excluded in the same way.  On the
other hand, from the discussion in \S\,71\,(e), it would appear that for
$(1-\mu) \to \infty$ the basic physical phenomenon is not very different from
what happens when $\mu = 0$; the only difference is that it all happens in
the layers immediately adjoining the inner cylinder; so that in this case
also, overstability would seem unlikely.

While the foregoing arguments are not conclusive, it is not, in principle,
difficult to settle the question by actual computation.  Thus, the characteristic value problem presented by equations~\eqref{eq:7-196},~\eqref{eq:7-197}, and~\eqref{eq:7-200}\footnote{We have replaced~$\sigma$ in equations~\eqref{eq:7-196} and~\eqref{eq:7-197} by~$i\sigma$ to emphasize that we are now primarily interested in marginal states which are strictly oscillatory.}
can be solved by the method described in \S\,71\,(a).  We find that equation
\eqref{eq:7-213} is replaced by
\begin{multline}
\label{eq:7-287}
\Biggl[\frac{a(m^2\pi^2+a^2+b^2)(2m^2\pi^2+a^2+b^2)}{(n^2\pi^2+a^2)(n^2\pi^2+a^2+b^2)(m^2\pi^2+a^2+b^2)}(-1)^{m+n}-1\Biggr] - {} \\
\frac{mn\pi^2(a^2-b^2)}{(n^2+a^2)(m^2\pi^2+a^2+b^2)}\bigl\{(\sigma\sinh a - a\sinh b)[(-1)^{m+1}+(1+\alpha)(-1)^{m+1}] + {} \\
{}+(\alpha\cosh b\sinh a - a\sinh b\cosh a)[1+(1+\alpha)(-1)^{m+n}] + {} \\
\frac{2\alpha(2m^2\pi^2+a^2+b^2)}{(m^2\pi^2+a^2)(m^2\pi^2+b^2)}
  \bigl[ab(\cosh a - \cosh b)[(-1)^{m+1}-(-1)^{m+1}] + {} \\
{}+\tfrac{1}{2}(a^2-b^2)\sinh a\sinh b[(-1)^{m+n}-1]\bigr]\Biggr\} + {} \\
{}+\tfrac{1}{2}\delta_{m,n} + \alpha X_{m,n}
  - \tfrac{1}{4}(n^2\pi^2+a^2)(n^2\pi^2+a^2+b^2)\frac{2m\pi}{|_{v_n}|}
= 0,
\end{multline}
where
\begin{equation}
\label{eq:7-288}
b^2 = a^2 + i\sigma,
\end{equation}
\begin{equation}
\label{eq:7-289}
\Delta = 2ab(1-\cosh a\cosh b) + (a^2+b^2)\sinh a\sinh b,
\end{equation}
and
\begin{equation}
\label{eq:7-290}
X_{m,n} = \begin{cases}
0 & \text{if } m+n \text{ is even and } m \ne n, \\
\frac{1}{2} & \text{if } m = n, \\
\dfrac{4mn}{n^2-m^2}\!\left(
  \dfrac{2m^2n^2+a^2+b^2}{(m^2\pi^2+a^2)(m^2\pi^2+b^2)}
  - \dfrac{1}{n^2(n^2-m^2)}\right)
& \text{if } m+n \text{ is odd.}
\end{cases}
\end{equation}
And the question is: Is there a real~$\sigma$ for which (for some given~$a$) the
characteristic root of the secular equation~\eqref{eq:7-287} is real?  If the answer is
in the affirmative, then overstability is possible.

A detailed examination of the roots of equation~\eqref{eq:7-287} has not been
undertaken.  It would be particularly worthwhile to explore the case
$\mu = -1$.

%% §73 %%
\section{The solution for a wide gap when the marginal state is stationary}\label{sec:7-73}

We now return to a consideration of the equations~\eqref{eq:7-167}--\eqref{eq:7-170} without
making any assumptions about the relative width of the gap $R_2-R_1$
between the cylinders.  But we shall continue to restrict ourselves to the
case when the onset of instability is as a stationary pattern of secondary
flow.  The relevant equations then are
\begin{equation}
\label{eq:7-291}
(DD_*-a^2)^2\mathfrak{u} = -Ta^2\!\left(\tfrac{1}{2}-\kappa\right)v
\end{equation}
and
\begin{equation}
\label{eq:7-292}
(DD_*-a^2)v = \mathfrak{u}.
\end{equation}
Eliminating~$\mathfrak{u}$ between equations~\eqref{eq:7-291} and~\eqref{eq:7-292}, we obtain
\begin{equation}
\label{eq:7-293}
(DD_*-a^2)^3 v = -Ta^2\!\left(\tfrac{1}{2}-\kappa\right)v;
\end{equation}
and the appropriate boundary conditions are
\begin{equation}
\label{eq:7-294}
v = (DD_*-a^2)v = D(DD_*-a^2)v = 0
\quad\text{for } r = 1 \text{ and } \eta.
\end{equation}

A solution of equation~\eqref{eq:7-293} patterned after the method described in
\S\,71\,(a) proves impracticable.  An alternative method based on expansions in terms of orthogonal functions satisfying four boundary conditions
appears well suited to the problem at hand.  The method is the following.

Letting
\begin{equation}
\label{eq:7-295}
G = (DD_*-a^2)v,
\end{equation}
we observe that the boundary conditions~\eqref{eq:7-294} require that both~$G$
and its derivative vanish at $r = 1$ and~$\eta$.  Accordingly, we may expand
$G$ in terms of the set of orthogonal functions, $\mathscr{G}_i(\alpha_i r)$, determined by the
characteristic value problem specified by the equation
\begin{equation}
\label{eq:7-296}
(DD_*)^2 y = \left(\frac{d^2}{dr^2}+\frac{1}{r}\frac{d}{dr}-\frac{1}{r^2}\right)^2 y_i = \alpha_i^4 y,
\end{equation}
and the boundary conditions
\begin{equation}
\label{eq:7-297}
y = 0 \quad\text{and}\quad dy/dr = 0 \quad\text{for } r = 1 \text{ and } \eta.
\end{equation}
The required characteristic functions are expressible as linear combinations of the Bessel functions $J_1$, $Y_1$, $I_1$, and~$K_1$ in the form
\begin{equation}
\label{eq:7-298}
\mathscr{G}_i(\alpha_i r) = A_i J_1(\alpha_i r) + B_i Y_1(\alpha_i r) + C_i I_1(\alpha_i r)
  + D_i K_1(\alpha_i r),
\end{equation}
where~$\alpha_i$ is a root of a certain transcendental equation (given in Appendix
V, equation~(31)) and $A_i$, $B_i$, $C_i$, and~$D_i$ are constants determined apart
from an arbitrary constant of proportionality.

The basic idea, then, is to expand~$G$ in terms of the functions
$\mathscr{G}_i(\alpha_i r)$.  Thus, we assume
\begin{equation}
\label{eq:7-299}
G = (DD_*-a^2)v = \sum_i P_i\,\mathscr{G}_i(\alpha_i r),
\end{equation}
where the coefficients~$P_i$ in the expansion are left unspecified at this stage.
Having expressed~$G$ in this form, we next solve equation~\eqref{eq:7-299} as a
differential equation for~$v$ and arrange that it satisfies the remaining
boundary conditions, namely, that $v = 0$ for $r = 1$ and~$\eta$.  With~$v$
determined in this fashion, equation~\eqref{eq:7-291} will lead, as we shall presently
see, to a secular equation for~$T$.

\subsection*{(a) \textit{The characteristic equation}}

We shall now obtain the explicit form of the characteristic equation
for~$T$.  For this purpose it is convenient to write $\mathscr{G}_i(\alpha_i r)$ in the manner
\begin{equation}
\label{eq:7-300}
\mathscr{G}_i(\alpha_i r) = \mathfrak{u}_i(r) + v_i(r),
\end{equation}
where
\begin{align}
\label{eq:7-301}
\mathfrak{u}_i(r) &= A_i J_1(\alpha_i r) + B_i Y_1(\alpha_i r) \notag\\
\intertext{and}
v_i(r) &= C_i I_1(\alpha_i r) + D_i K_1(\alpha_i r).
\end{align}
Clearly,
\begin{equation}
\label{eq:7-302}
DD_*\,\mathfrak{u}_i = -\alpha_i^2\,\mathfrak{u}_i
\quad\text{and}\quad
DD_*\,v_i = +\alpha_i^2\,v_i.
\end{equation}
Making use of the relations~\eqref{eq:7-302}, we can readily verify that the general
solution of equation~\eqref{eq:7-299} is
\begin{equation}
\label{eq:7-303}
v = \sum_i P_i\!\left[
  p_i\,J_1(\alpha_i r) + q_i\,K_1(\alpha_i r)
  - \frac{\mathfrak{u}_i(r)}{\alpha_i^2+a^2}
  + \frac{v_i(r)}{\alpha_i^2-a^2}
\right],
\end{equation}
where the constants of integration $p_i$ and~$q_i$ are to be determined by
the boundary conditions on~$v$, namely, that~$v$ vanishes at $r = 1$ and~$\eta$.  These
latter conditions lead to the equations
\begin{equation}
\label{eq:7-304}
p_i\,J_1(a) + q_i\,K_1(a)
  = \frac{\mathfrak{u}_i(1)}{\alpha_i^2+a^2}
  - \frac{v_i(1)}{\alpha_i^2-a^2}
\end{equation}
and
\begin{equation}
p_i\,J_1(a\eta) + q_i\,K_1(a\eta)
  = \frac{\mathfrak{u}_i(\eta)}{\alpha_i^2+a^2}
  - \frac{v_i(\eta)}{\alpha_i^2-a^2}.
\end{equation}
In obtaining these equations from~\eqref{eq:7-303}, we have made use of the fact
that
\begin{equation}
\label{eq:7-305}
\mathfrak{u}_i(1) = -v_i(1) \quad\text{and}\quad
\mathfrak{u}_i(\eta) = -v_i(\eta).
\end{equation}
On solving equations~\eqref{eq:7-304}, we find:
\begin{align}
\label{eq:7-306}
p_i &= \frac{2\alpha_i^2}{\Delta(\alpha_i^2-a^2)}\bigl[
  +\mathfrak{u}_i(1)K_1(a\eta) - \mathfrak{u}_i(\eta)K_1(a)\bigr], \notag\\
q_i &= \frac{2\alpha_i^2}{\Delta(\alpha_i^2-a^2)}\bigl[
  -\mathfrak{u}_i(1)J_1(a\eta) + \mathfrak{u}_i(\eta)J_1(a)\bigr],
\end{align}
where
\begin{equation}
\label{eq:7-307}
\Delta = J_1(a\eta)K_1(a) - J_1(a)K_1(a\eta).
\end{equation}

Now substituting for~$v$ from equation~\eqref{eq:7-303} in equation~\eqref{eq:7-293}, we obtain
\begin{multline}
\label{eq:7-308}
\sum_i P_i\bigl[(\alpha_i^2+a^2)^2\mathfrak{u}_{ij} + (\alpha_i^2-a^2)^2v_{ij}\bigr] \\
= Ta^2\!\left(\tfrac{1}{2}-\kappa\right)\sum_j \sum_i P_i\!\left[
  p_i\,J_1(\alpha_i r) + q_i\,K_1(\alpha_i r)
  + \frac{\alpha_i^2(\mathfrak{u}_i + v_i)}{(\alpha_i^2+a^2)(\alpha_i^2-a^2)}
\right].
\end{multline}
Finally, multiplying this equation by $r(\mathfrak{u}_j + v_j)$ and integrating over the
range of~$r$, we obtain
\begin{multline}
\label{eq:7-309}
P_i(\alpha_i^2+a^2)N_i^2 + 2a^2\sum_j P_j \alpha_j^2 \Delta_i^{(j)} \\
= Ta^2\sum_j P_j\bigl[p_j \alpha \varepsilon I_j^{(+1)}(a)
  - I_j^{(-1)}(a)\bigr]
  + q_j\bigl[\alpha \varepsilon K_j^{(+1)}(a) - K_j^{(-1)}(a)\bigr] \\
  + \frac{\alpha_j^2}{(\alpha_j^2+a^2)}
    \bigl[\alpha N_{ij}^2 - \Delta_{ij}^{(-)}\bigr]
  - \frac{\alpha_j^2}{(\alpha_j^2-a^2)} M_{ij}\Biggr],
\end{multline}
where we have made use of the orthogonality relation
\begin{equation}
\label{eq:7-310}
\int (\mathfrak{u}_i+v_i)(\mathfrak{u}_j+v_j)\,dr = N_i^2\,\delta_{ij}.
\end{equation}
and introduced the abbreviations
\begin{align}
\label{eq:7-311}
\Delta_i^{(j)} &= \int (\mathfrak{u}_i-v_j)(\mathfrak{u}_i+v_j)a^{2+1}\,dr, \notag\\
I_j^{(\pm 1)}(a) &= \int J_1(ar)(\mathfrak{u}_j+v_j)a^{a\pm 1}\,dr, \notag\\
K_j^{(\pm 1)}(a) &= \int K_1(ar)(\mathfrak{u}_j+v_j)a^{a\pm 1}\,dr,
\end{align}
and
\begin{equation}
\label{eq:7-312}
M_{ij} = \int (\mathfrak{u}_i+v_i)(\mathfrak{u}_j+v_j)\frac{dr}{r}.
\end{equation}
Equation~\eqref{eq:7-309} leads to the secular equation
\begin{multline}
\label{eq:7-313}
\Biggl|N_i^2(\alpha_i^2+a^2)\,\delta_{ij}
  + 2a^2\alpha_j^2\Delta_i^{(j)} \\
-T\Bigl[p_j\,\alpha\,\varepsilon\,I_j^{(+1)}(a)
  - I_j^{(-1)}(a)\bigr]
  + q_j\bigl[\alpha\,\varepsilon\,K_j^{(+1)}(a) - K_j^{(-1)}(a)\bigr] \\
  - \frac{\alpha_j^2}{(\alpha_j^2+a^2)}
    \bigl[\alpha N_{ij}^2 - \Delta_{ij}^{(-)}\bigr]
  - \frac{\alpha_j^2}{(\alpha_j^2-a^2)}M_{ij}\Biggr|
= 0.
\end{multline}

Of the various matrix elements defined in equations~\eqref{eq:7-311} and~\eqref{eq:7-312},
those with the superscript~$(+1)$ in~\eqref{eq:7-311} can be evaluated explicitly.
The remaining matrix elements must be evaluated numerically.

\subsection*{(b) \textit{Numerical results for the case} $\eta = \frac{1}{2}$}

First, we may note that when $\eta = \frac{1}{2}$ the definitions of~$T$ and~$\kappa$ are
\begin{equation}
\label{eq:7-314}
\mathscr{T} = \frac{64}{9}\!\left(\frac{\Omega_1 R_1^2}{\nu}\right)^{\!2}
  (1-\mu)(1-4\mu)
\end{equation}
and
\begin{equation}
\label{eq:7-315}
\kappa = (1-4\mu)/(1-\mu).
\end{equation}
And Rayleigh's criterion for stability in this case is
\begin{equation}
\label{eq:7-316}
\mu > \tfrac{1}{4} \quad \text{(for stability).}
\end{equation}

\begin{table}[ht]
\centering
\caption{Taylor numbers for various assigned values of $\kappa$ and~$a$ ($\eta = \frac{1}{2}$)}
\label{tab:XXXIV}
\medskip
\small
\begin{tabular}{cccccc}
\hline
& & \textit{First} & \textit{Second} & \textit{Third} \\
$\kappa$ & $a$ & \textit{approxi-} & \textit{approxi-} & \textit{approxi-} \\
& & \textit{mation} & \textit{mation} & \textit{mation} \\
\hline
  & $6{\cdot}0$ & $1{\cdot}5530\times10^4$ & $1{\cdot}5470\times10^4$ & $1{\cdot}5370\times10^4$ \\
0 & $6{\cdot}2$ & $1{\cdot}5456\times10^4$ & $1{\cdot}5434\times10^4$ & $1{\cdot}5338\times10^4$ \\
  & $6{\cdot}4$ & $1{\cdot}5468\times10^4$ & $1{\cdot}5414\times10^4$ & $1{\cdot}5315\times10^4$ \\
\hline
    & $6{\cdot}0$ & $1{\cdot}5833\times10^4$ & $1{\cdot}5728\times10^4$ & \\
0.4 & $6{\cdot}2$ & $1{\cdot}5738\times10^4$ & $1{\cdot}5686\times10^4$ & $1{\cdot}5593\times10^4$ \\
    & $6{\cdot}4$ & $1{\cdot}5806\times10^4$ & $1{\cdot}5642\times10^4$ & $1{\cdot}5542\times10^4$ \\
\hline
    & $6{\cdot}0$ & $1{\cdot}6800\times10^4$ & & \\
0.6 & $6{\cdot}2$ & $1{\cdot}6604\times10^4$ & $1{\cdot}5642\times10^4$ & \\
    & $6{\cdot}4$ & $1{\cdot}6411\times10^4$ & $1{\cdot}6119\times10^4$ & \\
\hline
    & $5{\cdot}6$ & $1{\cdot}3890\times10^4$ & $1{\cdot}3386\times10^4$ & $1{\cdot}3110\times10^4$ \\
1.0 & $5{\cdot}8$ & $1{\cdot}3598\times10^4$ & $1{\cdot}3390\times10^4$ & $1{\cdot}3110\times10^4$ \\
    & $6{\cdot}0$ & $1{\cdot}3455\times10^4$ & $1{\cdot}3392\times10^4$ & \\
\hline
     & $5{\cdot}6$ & $1{\cdot}3468\times10^4$ & $1{\cdot}3492\times10^4$ & $1{\cdot}3352\times10^4$ \\
1.333 & $5{\cdot}8$ & $1{\cdot}5068\times10^4$ & $1{\cdot}3380\times10^4$ & \\
      & $6{\cdot}0$ & $1{\cdot}5316\times10^4$ & $1{\cdot}3307\times10^4$ & $1{\cdot}3045\times10^4$ \\
\hline
    & $5{\cdot}0$ & $1{\cdot}5035\times10^4$ & $1{\cdot}3807\times10^4$ & \\
1.6 & $5{\cdot}5$ & $1{\cdot}4961\times10^4$ & $1{\cdot}3007\times10^4$ & \\
    & $6{\cdot}0$ & $1{\cdot}1935\times10^4$ & $1{\cdot}1005\times10^4$ & $9{\cdot}8831\times10^3$ \\
    & $6{\cdot}5$ & $2{\cdot}0024\times10^4$ & $1{\cdot}4006\times10^4$ & \\
\hline
    & $7{\cdot}0$ & $7{\cdot}507\times10^4$ & & \\
1.8 & $7{\cdot}5$ & $7{\cdot}843\times10^4$ & $3{\cdot}0862\times10^4$ & \\
    & $8{\cdot}0$ & $8{\cdot}061\times10^4$ & & \\
\hline
    & $8{\cdot}4$ & & $2{\cdot}024$ & \\
1.9 & $8{\cdot}8$ & & $2{\cdot}065\times10^4$ & \\
    & $9{\cdot}2$ & & & \\
\hline
$\frac{9}{4}$ & $9{\cdot}6$ & & $5{\cdot}093\times10^4$ & $4{\cdot}305\times10^4$ \\
\hline
\end{tabular}
\end{table}

In Table~XXXIV the values of~$T$ obtained with the aid of the characteristic equation~\eqref{eq:7-313} in the different approximations (the `order' of the
approximation being the order of the determinant which is set equal
to zero in the determination of~$T$) are listed for various assigned values
of~$a$ and~$\kappa$.  For each value of~$\kappa$, values of~$a$ were chosen (by trial and
error) in the range in which~$T$ (as a function of~$a$) attains its minimum.
From an examination of this table, it appears that for $\kappa < 1{\cdot}6$ the
required values of~$T$ have been determined to within one per cent of the
true values.  For $\kappa = 1{\cdot}8$, $1{\cdot}9$, and~$2{\cdot}0$, the results obtained in the second
and the third approximations differ appreciably.  Nevertheless, it is
unlikely that even in these cases the results of the third approximation
are in error by more than 3~per cent: this estimate of the error does not
appear unreasonable when we observe that for $\kappa = 1{\cdot}4$, while the results
of the first and the second approximations differ by 20~per cent, the

\figurePlaceholder{72}{The critical Taylor number $T_c$ for the onset of instability as a function of $\kappa = (1-4\mu)/(1-\mu)$.  A scale of $\mu$ ($= \Omega_2/\Omega_1$) is also shown.}

results of the second and the third approximations do not differ by more
than 2~per cent.

In Table~XXXV the critical Taylor numbers (derived from the data
of Table~XXXIV) appropriate for the different values of~$\kappa$ are given;
the values of~$a$ at which these minimum Taylor numbers are attained are
also given.  The derived $(T_c,\,\kappa)$ and~$[a(T_c),\,\kappa]$-relationships are further
illustrated in Figs.~72 and~73.  From Fig.~72 it appears that in a $(\log T_c,\,\kappa)$-plot, the relationship is very nearly a linear one in the neighbourhood
of $\kappa = 2$.

Table~XXXV also includes the values
\begin{equation}
\label{eq:7-317}
\omega_R = \frac{\Omega_1 R_1^2}{\nu}
  = 0{\cdot}375\sqrt{\frac{T}{(1-\mu)(1-4\mu)}}
  = \frac{0{\cdot}375}{1-\mu}\sqrt{\frac{T}{\kappa}};
\end{equation}
and $\omega_R$ and $\omega_R$, are, therefore, the angular velocities $\Omega_1$ and~$\Omega_2$ measured in
the unit $\nu/R_1^2$.  In the $(\omega_R,\,\omega_R)$-plane, the locus determined by equations
\eqref{eq:7-317} separates the regions of stability from the regions of instability
(see Fig.~74).  In this plane, Rayleigh's criterion is represented by the
straight line
\begin{equation}
\label{eq:7-318}
\omega_R = \tfrac{1}{4}\omega_1 \quad \text{(Rayleigh's criterion).}
\end{equation}

\figurePlaceholder{73}{The variation of the wave number~$a$ (in the unit $1/R_2$) of the disturbance at which instability first sets in as a function of~$\kappa$ and~$\mu$.  The uncertainty in the determination of~$a$ is indicated by the height of the lines.}

\begin{table}[ht]
\centering
\caption{Critical Taylor numbers and related constants for various values of $\kappa$
  ($\eta = \frac{1}{2}$)}
\label{tab:XXXV}
\medskip
\small
\begin{tabular}{ccccccc}
\hline
$\mu$ & $\kappa$ & $a$ & $T_c$ & $\omega_R$ & $\omega_R$ \\
\hline
0       & $+0{\cdot}25$ & $6{\cdot}1$ & $1{\cdot}333\times10^4$ &         & \\
$10{\cdot}1$   & $+0{\cdot}12076$ & $6{\cdot}1$ & $1{\cdot}631\times10^4$ & $190{\cdot}3$ & $+42{\cdot}30$ \\
$10{\cdot}3$   & $-0{\cdot}19036$ & $6{\cdot}1$ & $1{\cdot}719\times10^4$ & $139{\cdot}3$ & $+49{\cdot}21$ \\
$10{\cdot}3$   & $+0{\cdot}31676$ & $6{\cdot}1$ & $1{\cdot}810\times10^4$ & $114{\cdot}0$ & \\
0.4     & $-0{\cdot}16667$ & $6{\cdot}2$ & $1{\cdot}954\times10^4$ & $69{\cdot}5$ & $+16{\cdot}38$ \\
$0{\cdot}6$ & $-0{\cdot}17647$ & $6{\cdot}2$ & $2{\cdot}864\times10^4$ & $81{\cdot}2$ & $-9{\cdot}73$ \\
\hline
$1{\cdot}0$ & $-0{\cdot}6$ & $5{\cdot}8$ & $1{\cdot}310\times10^4$ & $80{\cdot}4$ & \\
\hline
$1{\cdot}333$ & $-0{\cdot}425$ & $6{\cdot}0$ & $3{\cdot}335\times10^4$ & $66{\cdot}5$ & $-8{\cdot}310$ \\
$1{\cdot}5$ & & $6{\cdot}5$ & $6{\cdot}813\times10^4$ & $74{\cdot}15$ & \\
\hline
$1{\cdot}7$ & $-0{\cdot}30438$ & & $1{\cdot}377\times10^5$ & $81{\cdot}95$ & $+44{\cdot}60$ \\
$1{\cdot}8$ & & $7{\cdot}5$ & $1{\cdot}993\times10^5$ & $91{\cdot}10$ & $-23{\cdot}30$ \\
\hline
$1{\cdot}95$ & $-0{\cdot}30530$ & & $2{\cdot}435\times10^5$ & $97{\cdot}10$ & $-39{\cdot}42$ \\
$1{\cdot}9$ & $-0{\cdot}45573$ & & $2{\cdot}6$ & $1{\cdot}424\times10^5$ & $-44{\cdot}72$ \\
\hline
$1{\cdot}95$ & $-0{\cdot}49415$ & & $2{\cdot}550\times10^5$ & $109{\cdot}4$ & $-27{\cdot}71$ \\
$1{\cdot}99$ & $-0{\cdot}090$ & $9{\cdot}6$ & $4{\cdot}386\times10^5$ & & $-37{\cdot}79$ \\
\hline
\end{tabular}
\end{table}

Since (see the first entry in Table~XXXV)
\begin{equation}
\label{eq:7-319}
T_c \to 1{\cdot}553\times10^4 \quad \text{as } \mu \to \tfrac{1}{4}
  \text{ and } \kappa \to 0,
\end{equation}
it is clear that in the neighbourhood of the line~\eqref{eq:7-318}, as $\omega_R \to \infty$, the
asymptotic behaviour of the true $(\omega_R,\,\omega_R)$-locus is given by (cf.\ equation
\eqref{eq:7-317})
\begin{equation}
\label{eq:7-320}
\omega_2 = 0{\cdot}375\sqrt{\frac{1{\cdot}553\times10^4}{\kappa(1-4\mu)}}
  = \frac{53{\cdot}81}{\sqrt{\kappa}} \quad (\mu \to \tfrac{1}{4}).
\end{equation}

\figurePlaceholder{74}{The regions of stability and instability in the $(\omega_R,\,\omega_R)$-plane; $\omega_1$ and~$\omega_2$ are the angular velocities of the inner and the outer cylinder in the units $\nu/R_1^2$ and~$\nu/R_2^2$, respectively.}

Finally, in Figs.~75 and~76, the profiles of the velocity in the radial
direction and the corresponding cell patterns at marginal stability are
shown for two typical cases ($\mu = 0$ and $\mu = -4/11$).

\figurePlaceholder{75}{The velocity profile~(a) and the cell pattern~(b) at the onset of instability for the case $\kappa = 1$, $\mu = 0$, and $a = 6{\cdot}4$.  The stream function $\psi$ (or $v_z \cos ax$) has been normalized to unity and the cell pattern is drawn symmetrically about $x = 0$.  (The unit of length is the radius of the outer cylinder.)}

\figurePlaceholder{76}{The velocity profile~(a) and the cell pattern~(b) at the onset of instability for the case $\kappa = 1{\cdot}8$, $\mu = -4/11$, and $a = 7{\cdot}6$.  The stream function $\psi$ (or $v_z \cos ax$) has been normalized to unity and the cell pattern is drawn symmetrically about $x = 0$.  (The unit of length is the radius of the outer cylinder.)}

%% §74 %%
\section{Experiments on the stability of viscous flow between rotating cylinders}\label{sec:7-74}

The first experiments on the onset of instability in the viscous flow
between rotating coaxial cylinders are those of Taylor.  Experiments
similar to Taylor's, and others related to the same phenomenon, have
since been carried out by Lewis, Terada and Hattori, Wendt, Taylor
himself, Schultz-Grunow and Hein, Donnelly, and Donnelly and Fultz.
We shall restrict our account of the experiments to those of Donnelly,
and Donnelly and Fultz as they are not only typical of the class, but they
are also among the most precise and the most extensive; moreover, they
bear directly on certain aspects of the theory considered in this chapter.

First, we may observe that the experiments one can perform in the
present connexion are of two kinds: those which involve the measure\-ment
of the torque exerted on the outer cylinder (kept at rest) for different
rates of rotation of the inner cylinder, and those in which the motions in
the fluid are observed by suitable tracers.

In the experiments of the first kind, the onset of instability is detected
by a principle not unlike the Schmidt--Milverton principle for the
detection of the onset of thermal instability (cf.\ \S\,18\,($b$)).  In the present
instance, the torque exerted on the outer cylinder (which is suitably
suspended) is measured as a function of the angular velocity $\Omega_1$ of the
inner cylinder.  It is known that so long as the flow between the cylinders
is laminar and conforms to the law (1), the relation is a linear one and is
given by
\begin{equation}\label{eq:7-321}
\text{torque} = 4\pi\mu\,\frac{R_1^2\,R_2^2\,H}{R_2^2 - R_1^2}\,\Omega_1,
\end{equation}
where $H$ is the height of the suspended cylinder.  At the onset of instability,
the effective viscosity will abruptly start increasing, and this
will be reflected in a sudden break in the relation~\eqref{eq:7-321}.

In the experiments of the second kind, the onset of instability is
detected by visual observations of sudden changes in the character of the
fluid motion occurring as the relative angular velocity of the two
cylinders is increased.  While these experiments are not as precise as
those on the torque measurements for determining the critical angular
velocities at which instability arises, they are not restricted to any special
value of $\mu$ ($= \Omega_2/\Omega_1$) such as $\mu = 0$.  Indeed, the important range of
negative $\mu$ is readily accessible to observations only by this second
method.

%% ------------------------------------------------------------
%% (a) Critical Taylor numbers by torque measurements
%% ------------------------------------------------------------
\subsection{The determination of the critical Taylor numbers, for the case $\mu = 0$,
by torque measurements}\label{sec:7-74a}

The most precise experiments on the determination of the critical
Taylor numbers for the onset of instability, for the case $\mu = 0$, are those
of Donnelly.

In his experiments, Donnelly used a rotating cylinder viscometer (see
Fig.~77) in which only a central section of the outer cylinder was sus\-pended;
the fixed end sections served as a guard cylinder for eliminating
end effects.  The viscometer was so designed that if the suspended
cylinder departed by so much as 0$\cdot$005~in., in any direction from exact
concentricity with the guard cylinder, the suspension would not swing
freely.  The system was thus protected from external vibrations and
lateral motions.

\figurePlaceholder{77}{Sketch of the viscometer used in Donnelly's experiments.}

The suspension system for the outer cylinders consisted of the follow\-ing
parts.  The straps by which the outer cylinder was carried were
attached to a long quartz rod (0$\cdot$018~in.\ in diameter and 11~cm long) at
the end of which was a torsion fibre (0$\cdot$005~in.\ in diameter and 10~cm long).
The fibre was in turn attached to a shaft whose angular position could be
read on a dial provided with a vernier.  A small mirror affixed to the
torsion fibre was directly in front of a stationary mirror.  The torque
exerted on the suspended cylinder was measured by a null method by
turning the dial so as to bring into coincidence the image of a cross-hair
(observed through an auto-collimator) reflected in both mirrors.  The
angular deflexion necessary to bring about the coincidence is a measure
of the torque.

The inner cylinder was driven by a shaft coupled to a motor and gear
reducing system.  The period of rotation was determined photo\-electrically
by using suitable scaling circuits.  In the experiments, the
speed of rotation was maintained constant to well within 0$\cdot$1 per cent.

The viscometer provided two interchangeable inner cylinders: one of
them left a gap width ($d/R_2 \sim 0{\cdot}05$) sufficiently small for the theory
of \S\,71 to be applicable; and the other left a gap width equal to the radius
of the inner cylinder so that the results of the experiments could be
compared with the calculations of \S\,73\,($b$).  The relevant dimensions
of the viscometer are given in Table~XXXVI.

To maintain the temperature of the liquid at a constant value during
the experiments, the viscometer was placed in a beaker kept in a constant
temperature bath.

The liquids used in the experiments and their relevant properties are
listed in Table~XXXVII.

\begin{table}[ht]
\centering
\caption{The dimensions and constants of Donnelly's viscometer}
\label{tab:XXXVI}
\medskip
\begin{tabular}{ll}
\hline
\multicolumn{2}{c}{\textit{Narrow gap}} \\
\hline
$R_1$ & $= 1{\cdot}89935 \pm 0{\cdot}0001$~cm \\
$R_2$ & $= 2{\cdot}00023 \pm 0{\cdot}0001$~cm \\
$H$   & $= 4{\cdot}9987 \pm 0{\cdot}0003$~cm \\
$C$   & $= 3{\cdot}93 \times 10^{-4}$ \\
\hline
\multicolumn{2}{c}{\textit{Wide gap}} \\
\hline
$R_1$ & $= 0{\cdot}99963 \pm 0{\cdot}0001$~cm \\
$R_2$ & $= 2{\cdot}00023 \pm 0{\cdot}0001$~cm \\
$H$   & $= 4{\cdot}9987 \pm 0{\cdot}0003$~cm \\
$C$   & $= 0{\cdot}967 \times 10^{-4}$ \\
\hline
\end{tabular}
\end{table}

\begin{table}[ht]
\centering
\caption{List of fluids and their physical properties}
\label{tab:XXXVII}
\medskip
\begin{tabular}{lccc}
\hline
Fluid & Temperature (${}^{\circ}$C) & $\rho$ (g) & $\nu$ (cm$^2$/sec) \\
\hline
CCl$_4$                          & 25$\cdot$00 & 1$\cdot$585 & $(1{\cdot}756 \pm 0{\cdot}03) \times 10^{-2}$ \\
CCl$_4$                          & 20$\cdot$95 & 1$\cdot$591 & $(6{\cdot}009 \pm 0{\cdot}03) \times 10^{-3}$ \\
CS$_2$                           & 25$\cdot$00 & 1$\cdot$255 & $(1{\cdot}935 \pm 0{\cdot}02) \times 10^{-2}$ \\
oil I, lot 111$^{\dag}$          & 25$\cdot$00 & 0$\cdot$8244 & $1{\cdot}330 \times 10^{-1}$ \\
oil D, lot 111$^{\dag}$          & 25$\cdot$00 & 0$\cdot$2813 & $2{\cdot}347 \times 10^{-2}$ \\
\hline
\end{tabular}

\smallskip
{\footnotesize $^{\dag}$ These oils are provided by the National Bureau of Standards (Washington, D.C.)
for calibrating viscometers. The National Bureau of Standards does not state the limits
of error on their measurements.}
\end{table}

If $K$ is the torsion constant of the fibre and $\Phi$ is the deflexion in radians,
then according to equation~\eqref{eq:7-321},
\begin{equation}\label{eq:7-322}
K\Phi = 4\pi\mu\,\frac{R_1^2\,R_2^2\,H}{R_2^2 - R_1^2}\,\Omega_1.
\end{equation}
In Donnelly's apparatus, the dial attached to the torsion fibre was
divided into 1,000 parts so that if $\phi$ is the measured deflexion in dial
divisions,
\begin{equation}\label{eq:7-323}
\nu = \left(\frac{K}{4\pi}\,\frac{R_2^2 - R_1^2}{R_1^2\,R_2^2\,H} \times 10^{-3}\right)\phi P
= C\,\phi\,P \quad (\text{say}),
\end{equation}
where $P$ ($= 2\pi/\Omega_1$) is the period of rotation of the inner cylinder.
Equation~\eqref{eq:7-323} applies when the flow is laminar; when it is not, we may
still regard the quantity on the right-hand side as giving an \textit{effective
viscosity}.

The value of $C$ was determined by calibration with liquids of known
viscosity rather than by an absolute determination of the torsion
constant~$K$.  The values of $C$ given in Table~XXXVI were obtained in
this manner.

In carrying out the experiments, the speed of rotation of the inner
cylinder was slowly increased from the lowest setting, keeping the
outer cylinder near its null position all the time.

\medskip
\noindent (i) \textit{Results of the experiments with the narrow gap}

The definition of the Taylor number appropriate for this case is (see
equation~(198))
\begin{equation}\label{eq:7-324}
T = \frac{4A\Omega_1}{\nu^2}\,d^4,
\end{equation}
where $A$ is given in equation~(146).  For $\mu = 0$, the explicit form of $T$ is
\begin{equation}\label{eq:7-325}
T = \frac{4\,R_1^2\,d^4}{\bigl(R_2^2 - R_1^2\bigr)}\left(\frac{\Omega_1}{\nu}\right)^{\!2}.
\end{equation}
According to the results given in Table~XXXIII, the critical Taylor
number for $\mu = 0$ is $3{\cdot}390 \times 10^3$.  The corresponding critical value of $\Omega_1$
is given by
\begin{equation}\label{eq:7-326}
\Omega_1(\text{critical}) = \left(3{\cdot}390 \times 10^3\,\frac{R_2^2 - R_1^2}{4\,R_1^2\,d^4}\right)^{1/2}\!\nu.
\end{equation}
For the values of $R_1$ and $R_2$ given in Table~XXXVI,
\begin{equation}\label{eq:7-327}
\Omega_1(\text{critical}) = 9{\cdot}448 \times 10^4 \nu.
\end{equation}
For carbon tetrachloride at $25^{\circ}$~C instability is predicted for a period of
rotation
\begin{equation}\label{eq:7-328}
P(\text{critical}) = 1{\cdot}147 \pm 0{\cdot}006~\text{s (calculated)},
\end{equation}
where the limits of error are due to the uncertainty in $\nu$.

The results of Donnelly's experiments are shown in Fig.~78\,$a$, $b$.  In
these figures, $\phi P$ is plotted as a function of $P$.  We observe that as the
period of rotation of the inner cylinder is decreased, the viscosity as
given by equation~\eqref{eq:7-323} begins to increase abruptly at a certain period.
The break occurs with extreme sharpness; how sharply, can be seen from
Fig.~78\,$b$ in which the measurements in the immediate vicinity of the
critical period are exhibited.

The critical period of rotation determined, as intermediate between
the last of the stable and the first of the unstable points of observation, is
\begin{equation}\label{eq:7-329}
P(\text{critical}) = 1{\cdot}1337 \pm 0{\cdot}00095~\text{s (measured)}.
\end{equation}
This agrees very well with the calculated value~\eqref{eq:7-328}.

\figurePlaceholder{78a}{Plot of $\phi P$ (which is proportional to the effective viscosity)
as a function of the period of rotation, $P$, of the inner cylinder for the
narrow-gap case: $R_1 = 1{\cdot}9$~cm, $R_2 = 2{\cdot}0$~cm.  The liquid was carbon
tetrachloride at $25^{\circ}$~C\@.  $C$ is the calculated value for instability.}

\figurePlaceholder{78b}{Detail of measurements, shown in Fig.~78\,$a$,
about the critical speed.  The onset of instability is
characterized by a discontinuity both in the magnitude
of $\phi P$ and in the slope of the effective viscosity curve.}

\medskip
\noindent (ii) \textit{Results of the experiments with the wide gap} ($\eta = \frac{1}{2}$)

The definition of $T$ appropriate for this case is (cf.\ equation~(314))
\begin{equation}\label{eq:7-330}
T = \frac{64}{9}\!\left(\frac{\Omega_1\,R_1^2}{\nu}\right)^{\!2}.
\end{equation}
According to the results given in Table~XXXV, the critical Taylor
number for this case ($\kappa = 1$ and $\mu = 0$) is $3{\cdot}310 \times 10^4$.  The corresponding
critical value of $\Omega_1$ is given by
\begin{equation}\label{eq:7-331}
\Omega_1(\text{critical}) = \left(\frac{9 \times 3{\cdot}310 \times 10^4}{64\,R_1^4}\right)^{1/2}\!\nu.
\end{equation}
For the value of $R_1$ given in Table~XXXVI,
\begin{equation}\label{eq:7-332}
\Omega_1(\text{critical}) = 68{\cdot}28\,\nu.
\end{equation}
The experiments with this wide-gap viscometer were carried out with
an oil (oil~I in Table~XXXVII) supplied by the National Bureau of
Standards for calibrating viscometers.  For this oil, instability of flow
is predicted for a period of rotation
\begin{equation}\label{eq:7-333}
P(\text{critical}) = 0{\cdot}7508 \pm 0{\cdot}0004~\text{s (calculated)}.
\end{equation}
The results of Donnelly's experiments in this case are shown in Fig.~79\,$a$, $b$.  We observe that, again, an increase in the viscosity as given by
equation~\eqref{eq:7-323} begins abruptly at a certain critical period of rotation of

\figurePlaceholder{79a}{Plot of $\phi P$ as a function of $P$ for the wide-gap
case: $R_1 = 1{\cdot}0$~cm, $R_2 = 2{\cdot}0$~cm.  The liquid was oil~I at
$25^{\circ}$~C\@.  $C$ is the calculated value for instability.}

\figurePlaceholder{79b}{Detail of measurements, shown in Fig.~79\,$a$, near
the critical speed.  The onset of instability is characterized
by a discontinuity in the curve, but less pronounced than in the
narrow-gap case (Fig.~78\,$b$).  Curves 1 and 2 were taken
on different days.  $C$ is the calculated value for instability.}

\noindent the inner cylinder.  The measurements in the immediate vicinity of the
critical period exhibited in Fig.~79\,$b$ show how sharply it begins.

The two curves (1) and (2) in Fig.~79\,$b$ show that, while the magnitude
of $\phi P$ at the critical period varies from one sequence of measurements to
another, the critical period itself is reproducible; the value of the critical
period as experimentally determined is
\begin{equation}\label{eq:7-334}
P(\text{critical}) = 0{\cdot}7543 \pm 0{\cdot}0003~\text{s (measured)}.
\end{equation}
Again, this agrees very well with the calculated value~\eqref{eq:7-333}.

%% ------------------------------------------------------------
%% (b) Visual and photographic observations
%% ------------------------------------------------------------
\subsection{The dependence of the critical Taylor number on $\Omega_2/\Omega_1$. The results of
visual and photographic observations}\label{sec:7-74b}

While the experiments described in \S\,(a) above verify the theoretical
predictions regarding the critical Taylor numbers for the onset of
instability for the particular case $\mu = 0$, they do not provide
confirmation for other important aspects of the theory such as the
dependence on $\mu$ ($= \Omega_2/\Omega_1$) of the critical Taylor number and of the
wave number of the associated disturbance.  For these latter purposes,
it is essential that experiments are carried out in an apparatus in which
the two cylinders can be rotated independently of each other.
The first experiments in this category are those of Taylor.  Taylor's
experiments were carried out with an apparatus with a narrow enough
gap for the results of \S\,71\,($b$) to be applicable.  They have been repeated
by Donnelly and Fultz with an apparatus in which the gap was equal to
the radius of the inner cylinder so that the results of \S\,73\,($b$) are
similarly applicable.

Before presenting the results of Taylor's and Donnelly and Fultz's
experiments, we shall give a brief description of the apparatus used by
the latter authors (see Fig.~80).  The outer cylinder was a precision
bore Pyrex tube 94~cm long and inside diameter $R_2 = 6{\cdot}2846 \pm 0{\cdot}0006$~cm.
The inner cylinder was a brass tube turned to a radius
$R_1 = 3{\cdot}1432 \pm 0{\cdot}0001$~cm.  The Pyrex cylinder rested on a groove cut
in a plate which fitted on a brass base mounted on a rotating table.  The
brass cylinder was connected by means of a drive pin and locking collar
to a shaft which could be rotated independently of the brass base.  The
cylinders were rotated by separate pulleys and by two motors with
variable speed transmission.  The two cylinders were aligned to be coaxial
by suitable precision ball bearings.

\figurePlaceholder{80}{Sketch of the apparatus used in the experiments of Donnelly and
Fultz.  The ratio of the radii of the two cylinders is one-half.}

The motion of the fluid was traced by ink emitted from three small
holes at the middle of the inner cylinder.  The ink was a solution of
Nigrosine dye in the liquid used in the experiments and was of such
strength that drops of it neither sank nor rose in the liquid under the
experimental conditions.

The liquids used in the experiments were distilled water and glycerine
water mixtures of different strengths.  Kinematic viscosities in the range
$0{\cdot}01$--$0{\cdot}20$~cm$^2$/sec were thus available for the experiments.

In the experiments, the angular velocity of the outer cylinder was
treated as an independent variable.  The outer cylinder was first brought
to a certain speed and was, thereafter, kept at that speed.  After the
conditions had become stationary, the inner cylinder was set in rotation
and its speed was gradually increased until it was within a few per cent of
the expected critical value (as determined by preliminary observations).
A small amount of ink was then injected and allowed to spread over a
small region of the inner cylinder.  The speed of the inner cylinder was
then increased, slowly and carefully, until a definite sustained three-dimensional
motion was observed which indicated the onset of instability.

%% ------------------------------------------------------------
%% (i) Wave number observations
%% ------------------------------------------------------------
\medskip
\noindent (i) \textit{Observations on the wave numbers of the disturbance manifested at
marginal stability}

In some ways, the most striking phenomenon observed in these
experiments is the sudden appearance of regularly spaced toroidal
vortices of rectangular cross-section.  An example of such manifestation,
as first photographed by Taylor, is shown in Fig.~82.

The height of the vortices, which appear at marginal stability, is
related to the wave number $a$ introduced in the theory by
\begin{equation}\label{eq:7-335}
\text{height of vortex} = R_2\,\pi/a.
\end{equation}
From the measured heights, the wave number $a$ can be deduced.  In
Fig.~81 a comparison is made between the experimentally deduced and
the theoretically predicted values of~$a$.

\figurePlaceholder{81}{Comparison of the calculated (full-line curve and the dots with
vertical lines) and observed (dots) wave numbers at the onset of instability.}

Photographs of the vortices obtained by Donnelly and Fultz under
various conditions are shown in Figs.~83--84.

The onset of instability at $\mu = 0$ is illustrated by the sequence of
photographs $a$, $b$, and $c$ in Fig.~82: ($a$)~shows the liquid in laminar flow
at a speed slightly below critical; ($b$)~shows the beginnings of three-dimensional
motion when the ink gathering together starts to move

\figurePlaceholder{85}{Comparison between the observed and the
predicted $(\Omega_1, \Omega_2)$-relation for the onset of instability
for the case $\eta = \frac{1}{2}$: the dots represent the points at
which instability was observed to occur in the experi\-ments
of Donnelly and Fultz and the full-line curve
represents the theoretical relation.}

\noindent radially outwards to form the first cell; and ($c$)~shows the vortex fully
developed from the initial state shown in ($b$) with no additional change
in the speed.

When both cylinders are rotated in the same direction, the appearance
of the cells is much the same as in the case $\mu = 0$.  An example is shown
in Fig.~82\,$d$.

When the cylinders are rotated in opposite directions, the phenomena
one observes are much more complex.  Theoretically, we expect two cells
in each compartment, with the motion in the outer cell relatively much
weaker (see Fig.~71\,$a$, $b$).  Since the ink is ejected from holes in the inner
cylinder, we should not, under the circumstances, expect to witness
motions in the outer cell.  Also, for reasons which we have explained in
\S\S\,68 and 71\,($b$), we must expect that, as $\mu \to -\infty$, the cells at marginal
stability appear very close to the inner cylinder and become increasingly
tiny.  All these expectations are strikingly confirmed by the sequence of
photographs in Figs.~83 and 84.

\figurePlaceholder{82}{The onset of instability with the outer cylinder at rest, $\mu = 0$. $P_c = 4{\cdot}493$~sec;
(a) laminar flow, $P = 4{\cdot}500$~sec; (b) beginning of radial motion at $P = 4{\cdot}483$~sec;
(c) Appearance of cells with outer cylinder at rest: cells at marginal stability, $P = P_c = 4{\cdot}466$~sec.
(d) Appearance of cells with cylinders rotating in the same direction: $P = 3{\cdot}544$~sec, $\mu = 0{\cdot}1184$.}

\figurePlaceholder{83}{Photographs of cells with cylinders ($\eta = \frac{1}{2}$) rotating in opposite directions:
(a) $\mu = -0{\cdot}4116$; (b) $\mu = -0{\cdot}5$; and (c) $\mu = -1{\cdot}74$.}

\figurePlaceholder{84}{Photographs of cells with cylinders ($\eta = \frac{1}{2}$) rotating in opposite directions:
(a) $\mu = -2{\cdot}44$; (b) $\mu = -3{\cdot}02$; (c) $\mu = -5{\cdot}86$; and (d) $\mu = -6{\cdot}83$.}

%% ------------------------------------------------------------
%% (ii) Comparison with predicted relations
%% ------------------------------------------------------------
\medskip
\noindent (ii) \textit{Comparison between the measured and the predicted $(T_c, \mu)$-relations}

The results of Donnelly and Fultz, obtained with their apparatus on
the critical angular velocities at which instability occurs, are exhibited
in Figs.~85 and 86 as plots analogous to those of Figs.~72 and 74.  The
agreement of the experimental results with the calculations of \S\,73\,($b$)
is very satisfactory.

\figurePlaceholder{86}{The variation of the critical Taylor number $(T_c)$ for the
case $\eta = \frac{1}{2}$ as a function of $\kappa\;(= (1-4\mu)/(1-\mu))$.  The solid line repre\-sents
the theoretically calculated relation and the dashed line is a
linear extrapolation of it.  The dots represent the points at which
instability was observed to occur.  Data to the left of the arrow are
unduly influenced by small errors in $\mu$.}

As we have already remarked, one of Taylor's original sets of measurements
was obtained by an apparatus in which the gap width was a small
fraction ($\sim\!5$ per cent) of the radius of the outer cylinder.  These measurements
can, therefore, be compared with the theoretical results of \S\,71\,($b$).
This comparison is included in Fig.~87 (curve labelled $\eta = 0{\cdot}8418$).
The agreement of the theoretical curve with the results of the measurements
is good over the entire range.

Fig.~87 includes measurements by Taylor for other values of $\eta$; and
also of Donnelly and Fultz for values of $-\mu$ beyond the range of the
calculations of \S\,73\,($b$).  From an examination of the data presented in
Fig.~87, Donnelly and Fultz have inferred that an asymptotic relation
of the form,
\begin{equation}\label{eq:7-336}
(\Omega_1\,R_1^2/\nu)^5 \to \text{constant}\times (|\Omega_2|\,R_2^2/\nu)^3
\quad\text{for}\quad \Omega_2/\Omega_1 \to -\infty,
\end{equation}
exists.

\figurePlaceholder{87}{Comparisons between the observed and the predicted $(\omega_1, \omega_2)$-relations for
the onset of stability: $\omega_1 = \Omega_1\,R_1^2/\nu$ and $\omega_2 = \Omega_2\,R_2^2/\nu$.  The uppermost solid curve
and the one at the extreme left represent the theoretical relations derived in \S\,71\,($b$)
and \S\,73\,($b$), respectively; the dashed-line continuation of the latter represents an
extrapolation of the theoretical relation.  The remaining solid lines are derived from
the asymptotic relation~\eqref{eq:7-342}.  The experimental data for $\eta = 0{\cdot}9418$, $0{\cdot}8798$, and
$0{\cdot}7435$ are those of G.~I.~Taylor; the data for $\eta = 0{\cdot}5$ are those of Donnelly and Fultz.}

Donnelly and Fultz have pointed out that the asymptotic relation
\eqref{eq:7-336} has a simple physical interpretation.  Arguing as in \S\,71\,(e), they
first observe that the onset of instability, in the case when $\mu < 0$, is
primarily to be associated with the violation of Rayleigh's criterion in
the layers immediately adjoining the inner cylinder and extending not
much beyond the nodal surface (at $r = R_n$, say).  Since the nodal
surface approaches the inner cylinder as $\mu \to -\infty$, the fluid which
becomes unstable is essentially confined to a narrow gap between $R_1$
and $R_n$.  Accordingly, we may write (cf.\ equation~(325))
\begin{equation}\label{eq:7-337}
\frac{4\Omega_1^2}{\nu^2}\,\frac{R_1^2\,d_0^4}{R_n^2 - R_1^2} = T_0
\end{equation}
as the condition for the onset of instability, where $d_0 = R_n - R_1$ and
$T_0$ is a number of the order of 1000.  Since the present arguments pre\-suppose
that $d_0 \ll R_1$, an equivalent form of~\eqref{eq:7-337} is
\begin{equation}\label{eq:7-338}
\left(\frac{\Omega_1\,R_1^2}{\nu}\right)^{\!2} = \tfrac{1}{2}\,T_0\!\left(\frac{R_1}{d_0}\right)^{\!3}\!.
\end{equation}
But, according to equation~(29),
\begin{equation}\label{eq:7-339}
d_0 = R_2\,\eta\!\left(\frac{1+|\mu|}{\eta^2+|\mu|}\right)^{\!1/2} - R_1
= R_1\!\left(\sqrt{\frac{1+|\mu|}{\eta^2+|\mu|}} - 1\right)\!.
\end{equation}
For $\mu \to -\infty$,
\begin{equation}\label{eq:7-340}
\frac{d_0}{R_1} \to \frac{1-\eta^2}{2|\mu|}
= \frac{1}{2}\left|\frac{\Omega_1}{\Omega_2}\right|(1-\eta^2).
\end{equation}
Inserting this value of $d_0/R_1$ in equation~\eqref{eq:7-338}, we obtain
\begin{equation}\label{eq:7-341}
\left(\frac{\Omega_1\,R_1^2}{\nu}\right)^{\!2}
= \frac{4\,T_0}{(1-\eta^2)^3}\left|\frac{\Omega_2}{\Omega_1}\right|^{\!3}
\qquad (\mu \to -\infty),
\end{equation}
or, alternatively,
\begin{equation}\label{eq:7-342}
\left(\frac{\Omega_1\,R_1^2}{\nu}\right)^{\!5}
= \frac{4\,\eta^6\,T_0}{(1-\eta^2)^3}\left(\frac{|\Omega_2|\,R_2^2}{\nu}\right)^{\!3}
\qquad (\mu \to -\infty),
\end{equation}
in agreement with~\eqref{eq:7-336}.

The values of $T_0$, empirically determined by Donnelly and Fultz for
various values of $\eta$, are
\begin{equation}\label{eq:7-343}
T_0 = \begin{cases}
450  & \eta = 0{\cdot}9418 \\
530  & \eta = 0{\cdot}8798 \\
670  & \eta = 0{\cdot}7435 \\
1{,}180 & \eta = 0{\cdot}5.
\end{cases}
\end{equation}


\section*{BIBLIOGRAPHICAL NOTES}

The fundamental papers on the stability of Couette flow are those of Lord
Rayleigh and Sir Geoffrey Taylor:

\begin{enumerate}
\item \textsc{Lord Rayleigh}, `On the dynamics of revolving fluids', \textit{Scientific Papers},
\textbf{6}, 447--53, Cambridge, England, 1920.

\item G.~I. \textsc{Taylor}, `Stability of a viscous liquid contained between two rotating
cylinders', \textit{Phil.\ Trans.\ Roy.\ Soc.\ (London)} A, \textbf{223}, 289--343 (1923).

\noindent In reference 1, Rayleigh treated the case of inviscid flow and derived the general
criterion known after him.  In reference~2, Taylor treated for the first time the
corresponding problem in viscid flow, both theoretically and experimentally.

\noindent The following general reference, in which the problem of hydrodynamic stability
is broadly surveyed with special reference to Couette flow, may be noted:

\item J.~L. \textsc{Synge}, `Hydrodynamic stability', 227--69, \textit{Semicentennial Addresses
of the American Mathematical Society}, ii, New York, 1938.

\item C.~C. \textsc{Lin}, \textit{The Theory of Hydrodynamic Stability}, Cambridge, England, 1955.

\noindent We may also refer to the relevant sections in:

\medskip
\noindent\S\,66.  As we have already remarked, Rayleigh developed his criterion for the
stability of inviscid rotational flow in reference~1.  In his discussion, Rayleigh
argued largely in terms of the analogy between the present problem and the
problem of the stability of an incompressible fluid of variable density.  For his
discussion of the latter problem, see:

\item Lord \textsc{Rayleigh}, `Investigation of the character of the equilibrium of an
incompressible heavy fluid of variable density', \textit{Scientific Papers}, \textbf{2},
200--7, Cambridge, England, 1900.

\medskip
\noindent\S\,67.  The analytical derivation of Rayleigh's criterion (for the case $m = 0$)
is due to:

\item J.~L. \textsc{Synge}, `The stability of heterogeneous liquids', \textit{Trans.\ of the Royal
Society of Canada}, \textbf{27}, 1--18 (1933).

\noindent The general case ($m \neq 0$) is considered in:

\item S. \textsc{Chandrasekhar}, `The stability of inviscid flow between rotating
cylinders', \textit{J.\ Indian Math.\ Soc.} (in press).

\noindent The discussion in this section is largely derived from reference~7.

\noindent A useful reference for the Sturmian theory and its later development is:

\item E.~L. \textsc{Ince}, \textit{Ordinary Differential Equations}, chapter~x, Longmans, Green \&
Co.\ Ltd., London, 1927.

\medskip
\noindent\S\,68.  The oscillations of a rotating column of liquid were first considered by
Lord Kelvin:

\item Lord \textsc{Kelvin}, `Vibrations of a columnar vortex', 152--65, \textit{Mathematical
and Physical Papers}, iv, \textit{Hydrodynamics and General Dynamics}, Cambridge,
England, 1910.

\noindent An extensive treatment of this subject with meteorological overtones will be
found in:

\item V. \textsc{Bjerknes}, J. \textsc{Bjerknes}, H. \textsc{Solberg}, and T. \textsc{Bergeron}, \textit{Physikalische
Hydrodynamik}, chapter~11, Springer, Berlin, 1933.

\medskip
\noindent\S\,69\,(a).  Fultz first showed how the natural modes of oscillation of a rotating
cylinder of liquid can be excited and maintained.

\item D. \textsc{Fultz}, `A note on overstability and the elasto-inertia oscillations of
Kelvin, Solberg, and Bjerknes', \textit{J.\ Meteorology}, \textbf{16}, 199--208 (1959).

\noindent The solution of equation~(96) for different values of $m$ was carried out by Miss
Donna Elbert; the results of her calculations are included in Fig.~63.
The roots listed in Table~XXXI (with the exception of those for $\eta = 0$) are taken
from:

\item S. \textsc{Chandrasekhar} and Donna \textsc{Elbert}, `The roots of
\[
Y_n(\lambda\eta)J_n(\lambda) - J_n(\lambda\eta)Y_n(\lambda) = 0\text{'},
\]
\textit{Proc.\ Camb.\ Phil.\ Soc.} \textbf{50}, 266--8 (1954).

\medskip
\noindent\S\,68\,(b).  The solution for the narrow gap described in this section is due to Reid:

\item W.~H. \textsc{Reid}, `Inviscid modes of instability in Couette flow', \textit{J.\ Math.\
Analysis and Applications}, \textbf{1}, 411--22 (1960).

\noindent I am grateful to Dr.~Reid for providing me with copies for the illustrations
(Figs.~65--67) included in this section.

\noindent A convenient reference for Airy's functions, in terms of which the solution for
this case is found, is:

\item J.~C.~P. \textsc{Miller}, `The Airy integral, giving tables of solutions of the
differential equation $y'' = xy$', \textit{British Assoc.\ Math.\ Tables}, Part-volume~B,
Cambridge, England, 1946.

\medskip
\noindent\S\,69.  The illustration (Fig.~68) in this section is taken from Taylor's original
paper (reference~2).

\noindent\S\,70.  The theorem proved in \S\,(a) is due to Synge.

\item J.~L. \textsc{Synge}, `On the stability of a viscous liquid between rotating coaxial
cylinders', \textit{Proc.\ Roy.\ Soc.\ (London)} A, \textbf{167}, 250--6 (1938).

\medskip
\noindent\S\,71\,(a).  The analysis in this section is based on:

\item S. \textsc{Chandrasekhar}, `The stability of viscous flow between rotating cylin\-ders',
\textit{Mathematika}, \textbf{1}, 5--13 (1954).

\noindent Alternative treatments of the characteristic value problem presented by equa\-tions~(200)--(202) have been given by:

\item D. \textsc{Meksyn}, `Stability of viscous flow between rotating cylinders.~I',
\textit{Proc.\ Roy.\ Soc.\ (London)} A, \textbf{187}, 115--28 (1946).

\item \leavevmode\kern-0.3em---\kern0.3em `Stability of viscous flow between rotating cylinders.  II. Cylinders
rotating in opposite directions', ibid.\ 480--91 (1946).

\item \leavevmode\kern-0.3em---\kern0.3em `Stability of viscous flow between rotating cylinders.  III.~Integra\-tion
of a sixth order linear equation', ibid.\ 492--504 (1946).

\item R.~C. \textsc{Di Prima}, `Application of the Galerkin method to problems in
hydrodynamic stability', \textit{Quart.\ Appl.\ Math.} \textbf{13}, 55--62 (1955).

\noindent The methods used by these authors (as well as by Taylor in reference~2) do not
enable as systematic (or as accurate) a solution of the problem as the method
described in the text.

\noindent An extension of the analysis in reference~16 to include a term in $\varepsilon^2$ on the right-hand
side of equation~(201) has been given by:

\item H. \textsc{Steinman}, `The stability of viscous flow between rotating cylinders',
\textit{Quart.\ Appl.\ Math.} \textbf{14}, 27--33 (1956).

\medskip
\noindent\S\,71\,(b).  The numerical results given in reference~16 have been supplemented.
These additional results (even as the original ones in reference~16) are due to
Miss Donna Elbert.

\medskip
\noindent\S\,71\,(c).  The alternative method described in this section has not been published
before.

\medskip
\noindent\S\,71\,(d).  The relation disclosed in this section between the present problem and
the simple B\'enard problem was first noticed by Low:

\item A.~R. \textsc{Low}, `Instability of viscous fluid motion', \textit{Nature}, \textbf{115}, 299--300 (1925).

\noindent The analytical basis for this identity in the characteristic values of the two
problems (in a certain approximation) was provided by:

\item H. \textsc{Jeffreys}, `Some cases of instability in fluid motion', \textit{Proc.\ Roy.\ Soc.\
(London)} A, \textbf{118}, 195--208 (1928).

\noindent The perturbation procedure described in this section is, however, new.  I am
indebted to Dr.~F.~Vandervort for the evaluation of the matrix elements in the
solution~(252).

\medskip
\noindent\S\,71\,(e).  The essential arguments of this section can be found (in varying degrees
of explicitness) in the writings of Meksyn (reference~18), Lin (reference~4), Di
Prima (reference~20), and Donnelly and Fultz (reference~28 below).

\medskip
\noindent\S\,72.  The discussion of overstability in this section has not been published before.

\medskip
\noindent\S\,73.  The analysis in this section is based on:

\item S. \textsc{Chandrasekhar}, `The stability of viscous flow between rotating cylin\-ders',
\textit{Proc.\ Roy.\ Soc.\ (London)} A, \textbf{246}, 301--11 (1958).

\noindent The method of solution derives from:

\item S. \textsc{Chandrasekhar} and W.~H. \textsc{Reid}, `On the expansion of functions which
satisfy four boundary conditions', \textit{Proc.\ Nat.\ Acad.\ Sci.} \textbf{43}, 521--7 (1957).

\noindent The orthogonal functions, $\Psi_i(a, r)$, in terms of which the solution is found, are
tabulated in:

\item S. \textsc{Chandrasekhar} and Donna~D. \textsc{Elbert}, `On orthogonal functions
which satisfy four boundary conditions.  III.~Tables for use in Fourier--Bessel-type
expansions', \textit{Astrophys.\ J.\ Supp.\ Ser.} \textbf{3}, 453--8 (1958).

\medskip
\noindent\S\,74.  The references to the experimental work described in this section are:

\item R.~J. \textsc{Donnelly}, `Experiments on the stability of viscous flow between
rotating cylinders.  I.~Torque measurements', \textit{Proc.\ Roy.\ Soc.\ (London)} A,
\textbf{246}, 312--25 (1958).

\item \leavevmode\kern-0.3em---\kern0.3em and D. \textsc{Fultz}, `Experiments on the stability of viscous flow between
rotating cylinders.  II.~Visual observations', ibid.\ \textbf{258}, 101--23 (1960).

\noindent The results of similar and related experiments will be found in:

\item J.~W. \textsc{Lewis}, `An experimental study of the motion of a viscous liquid
contained between two coaxial cylinders', \textit{Proc.\ Roy.\ Soc.\ (London)} A, \textbf{117},
388--407 (1928).

\item T. \textsc{Terada} and T. \textsc{Hattori}, `Some experiments on motions of fluids, IV',
\textit{Rep.\ Aero.\ Res.\ Inst.\ Tokyo}, \textbf{2}, 287--326 (1926).

\item F. \textsc{Wendt}, `Turbulente Str\"omungen zwischen zwei rotierenden koaxialen
Zylindern', \textit{Ingen.\ Arch.} \textbf{4}, 577--96 (1933).

\item G.~I. \textsc{Taylor}, `Fluid friction between rotating cylinders.  I.~Torque
measurements', \textit{Proc.\ Roy.\ Soc.\ (London)} A, \textbf{157}, 546--64 (1936).

\item \leavevmode\kern-0.3em---\kern0.3em `Fluid friction between rotating cylinders. II. Distribution of
velocity between concentric cylinders when outer one is rotating and inner
one is at rest', ibid.\ 565--78 (1936).

\item F. \textsc{Schultz-Grunow} and H. \textsc{Hein}, `Beitrag zur Couettestr\"omung', \textit{Z.\ f.\
Flugwiss.} \textbf{4}, 28--30 (1956).

\end{enumerate}
